\chapter{实数}
\label{cha:real-numbers}

\index{real numbers|(}%
任何名副其实的数学基础最终都必须处理分析学意义下的实数构造，即作为完备阿基米德有序域的实数。
\index{ordered field}%
完备性有两种概念。Cauchy 意义下的完备要求实数在 Cauchy 序列\index{Cauchy!sequence}的极限下封闭；更强的 Dedekind 意义下的完备则要求在 Dedekind 割\index{cut!Dedekind}下封闭。
这对应于两种实数构造，我们分别在 \cref{sec:dedekind-reals} 与 \cref{sec:cauchy-reals} 中研究。在 \cref{RD-final-field,RC-initial-Cauchy-complete} 中，我们用泛性质刻画这两种构造：Dedekind 实数是终阿基米德有序域，而 Cauchy 实数是初 Cauchy 完备阿基米德有序域。

在传统构造数学中，
\index{mathematics!constructive}%
实数似乎总要做某些折衷。例如，Dedekind 实数更适合与幂集或其他某种非预定性形式配合，而 Cauchy 实数在可数选择存在时表现更好。
\index{axiom!of choice!countable}%
不过我们会给出一种新的 Cauchy 实数构造：把它作为一个高阶归纳-归纳类型；这看起来是第三条道路，既不需要幂集，也不需要可数选择。

在~\cref{sec:comp-cauchy-dedek} 中，我们比较这两种实数构造。Cauchy 实数可嵌入 Dedekind 实数。若排中律或可数选择成立，它们重合；但一般情形下该嵌入可能是真包含。

在~\cref{sec:compactness-interval} 中，我们考察闭区间~$[0,1]$ 的三种紧致性概念。先证明 $[0,1]$ 在度量意义下是紧的\indexdef{metrically compact}\indexdef{compactness!metric}，即完备且全有界，并说明度量紧空间上的一致连续映射具有预期性质。相反，Bolzano--Weierstra\ss{} 性质“每个序列都有收敛子序列”蕴含有限全知原理，它是排中律的一个实例。最后我们讨论 Heine--Borel 紧致性：有限子覆盖性质的朴素表述并不可行，但与证明相关的归纳覆盖概念可行。本节基本上属于标准构造分析。

同伦类型论中的实数与分析发展可以很容易与经典数学兼容。若假设排中律与选择公理，我们就得到标准经典分析：\index{mathematics!classical}\index{classical!analysis}Dedekind 与 Cauchy 实数重合；关于 Dedekind 实数非预定性本质的基础问题消失；区间也达到其可能达到的最强紧致性。

本章最后在 \cref{sec:surreals} 中把 Conway 的超实数构造为一个高阶归纳-归纳类型；
在单价类型论中，这一构造比在经典集合论中更自然。

除 \cref{cha:basics,cha:logic} 的基础理论外，如上所述，我们在 Cauchy 实数与超实数处使用了“高阶归纳-归纳类型”：这是把 \cref{cha:hits} 的思想与 \cref{sec:generalizations} 提到的归纳-归纳类型概念结合起来。
我们还将频繁使用 \cref{subsec:prop-trunc} 所述传统逻辑记号，以及（在 \cref{sec:piw-pretopos} 中证明的）我们的“集合”按预期方式工作的事实。

注意，圆的通用覆盖总空间——它在 \cref{subsec:pi1s1-homotopy-theory} 中扮演了与经典代数拓扑中“实数”相似的角色——\emph{并不是}我们这里要找的实数类型。该类型是可缩的，因此等价于单点类型，故不可能携带非平凡代数结构。



\section{有理数域}
\label{sec:field-rati-numb}

\indexdef{rational numbers}%
\indexsee{number!rational}{rational numbers}%
我们先构造有理数 \Q，因为实数可以看作对~\Q 的完备化。行家会指出，\Q 可由任意近似域替代，
\indexdef{field!approximate}%
即 \Q 的一个子环，其中存在任意精度的近似逆元\index{inverse!approximate}%
。一个例子是二进有理数环，
\index{rational numbers!dyadic}%
即形如 $n/2^k$ 的数。
若我们要在计算机上实现构造数学，
近似域会更合适；但这种精细问题就留给真正关心 $\pi$ 小数位的人。

我们在 \cref{sec:set-quotients} 中把整数 \Z 构造为 $\N\times
\N$ 的一个商，并观察到该商由一个幂等元生成。在
\cref{sec:free-algebras} 中我们看到 \Z 是 \unit 上的自由群；类似地也可证明
它是 \emptyt 上的自由交换环\index{ring}。有理数域 \Q
也沿同样思路构造，即取商
%
\[ \Q \defeq (\Z \times \N)/{\approx} \]
%
其中
\[ (u,a) \approx (v,b) \defeq (u (b + 1) = v (a + 1)). \]
%
换言之，二元组 $(u, a)$ 表示有理数 $u / (1 + a)$。由于我们巧妙地在分母~$a$ 上加了 1，因此不存在除以零。这里同样有代表元的典范选取，即最简分数。因此可应用 \cref{lem:quotient-when-canonical-representatives} 得到集合 \Q，它同样具有可判定相等性。
\index{decidable!equality}%

我们不再写出 \Q 上的算术运算细节，相信读者即便在分母加一的情形下也会做分数运算。
这里只记录结论：有理数有序域 \Q 的构造完全没有问题，并且其相等与次序都可判定。
它还可刻画为初有序域。
\index{initial!ordered field}%

\symlabel{positive-rationals}
\indexdef{rational numbers!positive}%
\indexdef{positive!rational numbers}%
最后，记 $\Qp \defeq \setof{ q : \Q | q > 0 }$ 为正有理数类型。

\section{Dedekind 实数}
\label{sec:dedekind-reals}

\index{real numbers!Dedekind|(}%
先回顾 Dedekind 构造的基本思想。我们使用双边 Dedekind 割，而非常见的单边版本，因为对称性使构造更优雅，并且它在构造性与经典情形下都可行。
\index{mathematics!constructive}%
一个\emph{Dedekind 割}\index{cut!Dedekind}由一对子集 $L, U \subseteq \Q$ 组成，分别称为\emph{下割}与\emph{上割}，并满足：
%
\begin{enumerate}
\item \emph{可居住：}存在 $q \in L$ 与 $r \in U$，
\item \emph{圆整：}$q \in L \Leftrightarrow \exis {r \in \Q} q < r \land r \in L$
  以及 $r \in U \Leftrightarrow \exis {q \in \Q} q \in U \land q < r$，
  \index{rounded!Dedekind cut}
\item \emph{不相交：}$\lnot (q \in L \land q \in U)$，以及
\item \emph{定位：}$q < r \Rightarrow q \in L \lor r \in U$。
  \index{locatedness}%
\end{enumerate}
%
把圆整条件从左到右读，说明割是\emph{开}的，
\index{open!cut}%
从右到左读，说明它们分别是\emph{下}集与\emph{上}集。定位条件说明 $L$ 与 $U$ 之间不存在“大空隙”。由于割总是开集，即使中间点是有理数，也不会包含该“中间点”。典型 Dedekind 割如下：
%
\begin{center}
  \begin{tikzpicture}[x=\textwidth]
    \draw[<-),line width=0.75pt] (0,0) -- (0.297,0) node[anchor=south east]{$L\ $};
    \draw[(->,line width=0.75pt] (0.300, 0) node[anchor=south west]{$\ U$} -- (0.9, 0) ;
  \end{tikzpicture}
\end{center}
%
我们也许会天真地把该非形式定义翻译成类型论：割是映射对 $L, U : \Q \to \prop$。但在 \cref{subsec:prop-subsets} 中我们看到 $\prop$ 是一个含糊\index{typical ambiguity}记号，实际指 $\prop_{\UU_i}$，其中 $\UU_i$ 是某个宇宙。一旦用某个特定 $\UU_i$ 定义割，实数类型就位于下一宇宙 $\UU_{i+1}$，关于实数的性质位于更高一层 $\UU_{i+2}$，关于实数子集的性质位于 $\UU_{i+3}$，依此类推。原则上我们应追踪这些宇宙层级\index{universe level}，特别在证明助手帮助下并非难事，但在这里这么做只会带来我们希望避免的繁琐手续。因此我们作一个简化假设：单一命题类型 $\Omega$ 足以满足本章需求。

事实上，Dedekind 实数构造对逻辑操作相当稳健。把单一类型 $\Omega$ 用好有多种方式：
%
\begin{enumerate}

\item 可把 $\Omega$ 与含糊记号 $\prop$ 视为同一，并追踪定义与构造中出现的全部宇宙。

\item 可假设命题重整公理，
  \index{propositional!resizing}%
  如 \cref{subsec:prop-subsets} 所述，它本质上把各层 $\prop_{\UU_i}$ 压到最低层\index{universe level}，记为 $\Omega$。

\item 若经典数学家不关心类型论宇宙或计算细节，可直接对纯命题假设排中律~\eqref{eq:lem}，令 $\Omega \jdeq \bool$。
  \index{excluded middle}
  这不仅消除了 $\prop$ 层级\index{universe level}问题，也把我们所做的一切变成标准经典\index{mathematics!classical}实数构造。

\item 在另一端，也可追求使构造可行的最小条件。一个纯谓词成为 Dedekind 割这一条件，只用到合取、析取以及对~\Q 的存在量词\index{quantifier!existential}，而 \Q 是可数集合。因此可把 $\Omega$ 取为初\emph{$\sigma$-框架}，
  \index{initial!sigma-frame@$\sigma$-frame}%
  \index{sigma-frame@$\sigma$-frame!initial|defstyle}%
  即带可数并\index{join!in a lattice}且二元交对可数并分配的格\index{lattice}。（初 $\sigma$-框架不可能是两点格 $\bool$，因为除非假设排中律，$\bool$ 对可数并并不封闭。）
  这会把 $\Omega$ 导向一个高阶归纳-归纳类型构造，但在 \cref{sec:cauchy-reals} 做一个这类实验已足够。
\end{enumerate}

在上述所有情形中，$\Omega$ 都是集合。
%
不再赘述，我们把非形式定义翻译到类型论中。
本章将始终使用
\cref{defn:logical-notation} 的逻辑记号。

\begin{defn} \label{defn:dedekind-reals}
  一个\define{Dedekind 割}
  \indexsee{Dedekind!cut}{cut, Dedekind}%
  \indexdef{cut!Dedekind}%
  是一对纯谓词 $L : \Q \to \Omega$ 与 $U
  : \Q \to \Omega$，满足：
  %
  \begin{enumerate}
  \item \label{defn:dedekind-reals-inhabited}
    \emph{可居住（即有界）：}$\exis{q : \Q} L(q)$ 与 $\exis{r : \Q} U(r)$，
  \item \emph{圆整：}对所有 $q, r : \Q$，
    \index{rounded!Dedekind cut}
    %
    \begin{align*}
      L(q) &\Leftrightarrow \exis{r : \Q} (q < r) \land L(r)
      \qquad\text{且}\\
      U(r) &\Leftrightarrow \exis{q : \Q} (q < r) \land U(q),
    \end{align*}
  \item \emph{不相交：}对所有 $q : \Q$，$\lnot (L(q) \land U(q))$，
  \item \emph{定位：}对所有 $q, r : \Q$，$(q < r) \Rightarrow L(q) \lor U(r)$。
  \index{locatedness}%
  \end{enumerate}
  %
  记 $\dcut(L, U)$ 为这些条件的合取。\define{Dedekind 实数}类型是
  \indexsee{Dedekind!real numbers}{real numbers, De\-de\-kind}%
  \indexdef{real numbers!Dedekind}%
  %
  \begin{equation*}
    \RD \defeq \setof{ (L, U) : (\Q \to \Omega) \times (\Q \to \Omega) | \dcut(L,U)}.
  \end{equation*}
\end{defn}

显然 $\dcut(L, U)$ 是纯命题，而 $\Q \to \Omega$ 是集合，所以 Dedekind 实数也构成集合。关于导向扩展实数、上下实数与区间域的 Dedekind 割变体，见
\cref{ex:RD-extended-reals,ex:RD-lower-cuts,ex:RD-interval-arithmetic}。

有一个嵌入 $\Q \to \RD$，把每个有理数 $q : \Q$ 关联到割 $(L_q, U_q)$，其中
%
\begin{equation*}
  L_q(r) \defeq (r < q)
  \qquad\text{且}\qquad
  U_q(r) \defeq (q < r).
\end{equation*}
%
下文将直接把与有理数对应的割 $(L_q, U_q)$ 记作 $q$。

\subsection{Dedekind 实数的代数结构}
\label{sec:algebr-struct-dedek}

Dedekind 实数的代数与序结构按直觉主义逻辑中的常规方式构造即可。我们不展开所有细节，而只指出经典\index{mathematics!classical}与直觉主义语境的差异。记实数 $x : \RD$ 的下割、上割为 $L_x$ 与 $U_x$，则加法定义为%
%
\indexdef{addition!of Dedekind reals}%
\begin{align*}
  L_{x + y}(q) &\defeq \exis{r, s : \Q} L_x(r) \land L_y(s) \land q = r + s, \\
  U_{x + y}(q) &\defeq \exis{r, s : \Q} U_x(r) \land U_y(s) \land q = r + s,
\end{align*}
%
加法逆元定义为
%
\begin{align*}
  L_{-x}(q) &\defeq \exis{r : \Q} U_x(r) \land q = - r, \\
  U_{-x}(q) &\defeq \exis{r : \Q} L_x(r) \land q = - r.
\end{align*}
%
在这些运算下，$(\RD, 0, {+}, {-})$ 是阿贝尔\index{group!abelian}群。乘法稍繁琐：
%
\indexdef{multiplication!of Dedekind reals}%
\begin{align*}
  L_{x \cdot y}(q) &\defeq
  \begin{aligned}[t]
    \exis{a, b, c, d : \Q} & L_x(a) \land U_x(b) \land L_y(c) \land U_y(d) \land {}\\
                           & \qquad q < \min (a \cdot c, a \cdot d, b \cdot c, b \cdot d),
  \end{aligned} \\
  U_{x \cdot y}(q) &\defeq
  \begin{aligned}[t]
    \exis{a, b, c, d : \Q} & L_x(a) \land U_x(b) \land L_y(c) \land U_y(d) \land {}\\
                           & \qquad \max (a \cdot c, a \cdot d, b \cdot c, b \cdot d) < q.
  \end{aligned}
\end{align*}
%
\index{interval!arithmetic}%
这些公式与区间算术中的区间乘法相关：端点有理的区间 $[a,b]$ 与 $[c,d]$ 相乘得
%
\begin{equation*}
  [a,b] \cdot [c,d] =
  [\min(a c, a d, b c, b d), \max(a c, a d, b c, b d)].
\end{equation*}
%
例如，下割公式可读作：当存在包含 $x,y$ 的区间 $[a,b],[c,d]$ 且 $q$ 位于 $[a,b] \cdot [c,d]$ 左侧时，$q < x \cdot y$。
一般而言，把满足 $L_x(a)$ 且 $U_x(b)$ 的区间 $[a,b]$ 看作 $x$ 的近似是有益的，见
\cref{ex:RD-interval-arithmetic}。

现在得到交换幺环\index{ring}
\index{unit!of a ring}%
$(\RD, 0, 1, {+}, {-}, {\cdot})$。讨论乘法逆元前先引入次序。定义 $\leq$ 与 $<$ 为
%
\begin{align*}
  (x \leq y) &\ \defeq \ \fall{q : \Q} L_x(q) \Rightarrow L_y(q), \\
  (x < y)    &\ \defeq \ \exis{q : \Q} U_x(q) \land L_y(q).
\end{align*}

\begin{lem} \label{dedekind-in-cut-as-le}
  对所有 $x : \RD$ 与 $q : \Q$，有 $L_x(q) \Leftrightarrow (q < x)$ 且 $U_x(q)
  \Leftrightarrow (x < q)$。
\end{lem}

\begin{proof}
  若 $L_x(q)$，则由圆整性仅存在 $r > q$ 使 $L_x(r)$，又因 $U_q(r)$，可得 $q < x$。反之若 $q < x$，则存在 $r : \Q$ 使 $U_q(r)$ 且 $L_x(r)$，因此由 $L_x$ 是下集得 $L_x(q)$。另一半证明对称。
\end{proof}

\index{partial order}%
\index{transitivity!of . for reals@of $<$ for reals}
\index{transitivity!of . for reals@of $\leq$ for reals}
\index{relation!irreflexive}
\index{irreflexivity!of . for reals@of $<$ for reals}
关系 $\leq$ 是偏序，$<$ 具传递且反自反性。线性性
\index{order!linear}%
\index{linear order}%
%
\begin{equation*}
  (x < y) \lor (y \leq x)
\end{equation*}
%
在假设排中律时成立；若不假设，则可得弱线性
%
\index{order!weakly linear}
\index{weakly linear order}
\begin{equation} \label{eq:RD-linear-order}
  (x < y) \Rightarrow (x < z) \lor (z < y).
\end{equation}
%
乍看之下~\eqref{eq:RD-linear-order} 与线序关系不明显。但若取 $x \jdeq u - \epsilon$ 与 $y \jdeq u + \epsilon$（$\epsilon > 0$），得到
%
\begin{equation*}
  (u - \epsilon < z) \lor (z < u + \epsilon).
\end{equation*}
%
这就是“允许小数值误差的线性性”：既然期望能以无限精度计算并不合理，那么只能在我们已计算到的有限精度内判定~$<$，也并不奇怪。

要见~\eqref{eq:RD-linear-order} 成立，设 $x < y$。则仅存在 $q : \Q$ 使 $U_x(q)$ 与
$L_y(q)$。由圆整性，仅存在 $r, s : \Q$ 使 $r < q < s$、$U_x(r)$
与 $L_y(s)$。再由定位性，$L_z(r)$ 或 $U_z(s)$。前者给出 $x < z$，后者给出 $z < y$。

经典上，对所有非零数都存在乘法逆元。然而在不假设排中律时，需要更强条件。称 $x, y : \RD$
彼此\define{分离}
\indexdef{apartness}%
（记作 $x \apart y$）当且仅当 $(x < y) \lor (y < x)$：
%
\symlabel{apart}
\begin{equation*}
  (x \apart y) \defeq (x < y) \lor (y < x).
\end{equation*}
%
若 $x \apart y$，则 $\lnot (x = y)$。
反向在假设排中律时成立，但在构造性情形不可证。
\index{mathematics!constructive}%
实际上，若“$\lnot (x = y)$ 蕴含 $x\apart y$”，则可推出一点排中律；见 \cref{ex:reals-apart-neq-MP}。

\begin{thm} \label{RD-inverse-apart-0}
  一个实数可逆，当且仅当它与 $0$ 分离。
\end{thm}

\begin{rmk}
  我们注意到：实数可逆当且仅当它仅可逆。
  事实上在任意环\index{ring}中都如此，因为环是集合，且乘法逆元若存在则唯一。
  见 \cref{cor:UC} 后讨论。
\end{rmk}

\begin{proof}
  设 $x \cdot y = 1$。则仅存在 $a, b, c, d : \Q$ 使
  $a < x < b$、$c < y < d$ 且 $0 < \min (a c, a d, b c, b d)$。由 $0 < a c$ 与 $0 < b c$ 可知
  $a,b,c$ 要么全正要么全负。
  因而要么 $0 < a < x$，要么 $x < b < 0$，故 $x \apart 0$。

  反之，若 $x \apart 0$，则要么 $x > 0$，要么 $x < 0$。
  若 $x > 0$，定义 $x^{-1}$ 为：
  %
  \begin{align*}
    L_{x^{-1}}(q) &\defeq
    (q > 0) \Rightarrow \exis{r : \Q} U_x(r) \land (q r < 1),
    \\
    U_{x^{-1}}(q) &\defeq
    (q > 0) \land \exis{r : \Q} L_x(r) \land (q r > 1).
  \end{align*}
  %
  若 $x < 0$，则定义为
  %
  \begin{align*}
    L_{x^{-1}}(q) &\defeq
    (q < 0) \land \exis{r : \Q} U_x(r) \land (q r > 1),
    \\
    U_{x^{-1}}(q) &\defeq
    (q < 0) \Rightarrow \exis{r : \Q} L_x(r) \land (q r < 1). \qedhere
  \end{align*}
\end{proof}

\index{ordered field!archimedean}%
\index{dense}%
\indexsee{order-dense}{dense}%
阿基米德原理有多种表述。我们认为最有启发的是“\Q 在 \RD 中稠密”的形式。

\begin{thm}[\RD 的阿基米德原理] \label{RD-archimedean}
  %
  对所有 $x, y : \RD$，若 $x < y$，则仅存在 $q : \Q$ 使
  $x < q < y$。
\end{thm}

\begin{proof}
  由 $<$ 的定义。
\end{proof}

在讨论 Dedekind 实数完备性之前，先精确陈述其代数结构。下述定义并非追求最小公理化，而是追求实用且足够的结构与性质。

\begin{defn} \label{ordered-field} 一个\define{有序域}
  \indexdef{ordered field}%
  \indexsee{field!ordered}{ordered field}%
  是一个集合 $F$，并配备下列数据：
  \begin{itemize}
  \item 常元 $0,1$；
  \item 运算 $+, -, \cdot, \min, \max$；
  \item 纯关系 $\leq, <, \apart$。
  \end{itemize}
  这些数据满足：
  %
  \begin{enumerate}
  \item $(F, 0, 1, {+}, {-}, {\cdot})$ 是交换幺环；
    \index{unit!of a ring}%
    \index{ring}%
  \item $x : F$ 可逆当且仅当 $x \apart 0$；
  \item $(F, {\leq}, {\min}, {\max})$ 是格；
  \item 严格序 $<$ 传递、反自反，
    \index{relation!irreflexive}
    \index{irreflexivity!of . in a field@of $<$ in a field}%
    且弱线性（$x < y \Rightarrow x < z \lor z < y$）；\index{transitivity!of . in a field@of $<$ in a field}
    \index{order!weakly linear}
    \index{weakly linear order}
    \index{strict!order}%
    \index{order!strict}%
  \item 分离关系 $\apart$ 反自反、对称、余传递（$x \apart y \Rightarrow x \apart z \lor y \apart z$）；
    \index{relation!irreflexive}
    \index{irreflexivity!of apartness}%
    \indexdef{relation!cotransitive}%
    \index{cotransitivity of apartness}%
  \item 对所有 $x, y, z : F$：
    %
    \begin{align*}
      x \leq y &\Leftrightarrow \lnot (y < x), &
      x < y \leq z &\Rightarrow x < z, \\
      x \apart y &\Leftrightarrow (x < y) \lor (y < x), &
      x \leq y < z &\Rightarrow x < z, \\
      x \leq y &\Leftrightarrow x + z \leq y + z, &
      x \leq y \land 0 \leq z &\Rightarrow x z \leq y z, \\
      x < y &\Leftrightarrow x + z < y + z, &
      0 < z \Rightarrow (x < y &\Leftrightarrow x z < y z), \\
      0 < x + y &\Rightarrow 0 < x \lor 0 < y, &
      0 &< 1.
    \end{align*}
  \end{enumerate}
  %
  每个这样的域都有典范嵌入 $\Q \to F$。若有序域还满足：对所有 $x, y : F$，若 $x < y$ 则仅存在 $q :
  \Q$ 使 $x < q < y$，则称其为\define{阿基米德}有序域。
  \indexdef{ordered field!archimedean}%
  \indexsee{archimedean property}{ordered field, archi\-mede\-an}%
\end{defn}

\begin{thm} \label{RD-archimedean-ordered-field}
  Dedekind 实数构成阿基米德有序域。
\end{thm}

\begin{proof}
  我们略去证明，希望前面的展示已使该定理足够可信。
\end{proof}

\subsection{Dedekind 实数是 Cauchy 完备的}
\label{sec:RD-cauchy-complete}

回忆 $x : \N \to \Q$ 是\emph{Cauchy 序列}\indexdef{Cauchy!sequence}，当它满足
%
\begin{equation} \label{eq:cauchy-sequence}
  \prd{\epsilon : \Qp} \sm{n : \N} \prd{m, k \geq n} |x_m - x_k| < \epsilon.
\end{equation}
%
注意我们\emph{没有}对内层存在量词做截断，因为我们确实希望计算收敛速率——没有误差估计的近似信息价值很小。由 \cref{thm:ttac}，\eqref{eq:cauchy-sequence} 给出函数 $M
: \Qp \to \N$，称为\emph{收敛模}\indexdef{modulus!of convergence}，满足若 $m, k \geq M(\epsilon)$ 则 $|x_m - x_k| < \epsilon$。由此可得对所有 $\delta, \epsilon : \Qp$，有 $|x_{M(\delta/2)} - x_{M(\epsilon/2)}|<
\delta + \epsilon$。事实上，映射 $(\epsilon \mapsto
x_{M(\epsilon/2)}) : \Qp \to \Q$ 与原始 Cauchy 条件~\eqref{eq:cauchy-sequence} 对极限携带同样信息。我们将使用这些近似函数，而不是 Cauchy 序列。

\begin{defn} \label{defn:cauchy-approximation}
  一个\define{Cauchy 近似}
  \indexdef{Cauchy!approximation}%
  是映射 $x : \Qp \to \RD$，满足
  %
  \begin{equation}
    \label{eq:cauchy-approx}
    \fall{\delta, \epsilon :\Qp} |x_\delta - x_\epsilon| < \delta + \epsilon.
  \end{equation}
  %
  Cauchy 近似 $x : \Qp \to \RD$ 的\define{极限}
  \index{limit!of a Cauchy approximation}%
  是某个 $\ell : \RD$，使
  %
  \begin{equation*}
    \fall{\epsilon, \theta : \Qp} |x_\epsilon - \ell| < \epsilon + \theta.
  \end{equation*}
\end{defn}

\begin{thm} \label{RD-cauchy-complete}
  \RD 中每个 Cauchy 近似都有极限。
\end{thm}

\begin{proof}
  注意我们证明的是存在，而非仅存在。
  给定 Cauchy 近似 $x : \Qp \to \RD$，定义
  %
  \begin{align*}
    L_y(q) &\defeq \exis{\epsilon, \theta : \Qp} L_{x_\epsilon}(q + \epsilon + \theta),\\
    U_y(q) &\defeq \exis{\epsilon, \theta : \Qp} U_{x_\epsilon}(q - \epsilon - \theta).
  \end{align*}
  %
  显然 $L_y,U_y$ 可居住且圆整。不相交性由 Cauchy 近似条件推出。为证明定位性，取任意
  $q, r : \Q$ 且 $q < r$。存在 $\epsilon : \Qp$ 使
  $4 \epsilon < r - q$。由于 $q + 2 \epsilon < r - 2 \epsilon$，仅有
  $L_{x_\epsilon}(q + 2 \epsilon)$ 或 $U_{x_\epsilon}(r - 2 \epsilon)$。前者得
  $L_y(q)$，后者得 $U_y(r)$。

  要证明 $y$ 是 $x$ 的极限，取任意 $\epsilon, \theta : \Qp$。由于
  \Q 在 \RD 中稠密，仅存在 $q, r : \Q$ 使
  %
  \begin{narrowmultline*}
    x_\epsilon - \epsilon - \theta < q < x_\epsilon - \epsilon - \theta/2
    < x_\epsilon < \\
    x_\epsilon + \epsilon + \theta/2 < r < x_\epsilon + \epsilon + \theta,
  \end{narrowmultline*}
  %
  因而 $q < y < r$。由此得 $|y - x_\epsilon| < \epsilon + \theta$。
\end{proof}

为完整起见，我们也记录经典表述。

\begin{cor}
  设 $x : \N \to \RD$ 满足 Cauchy 条件~\eqref{eq:cauchy-sequence}。则
  存在 $y : \RD$ 使
  %
  \begin{equation*}
    \prd{\epsilon : \Qp} \sm{n : \N} \prd{m \geq n} |x_m - y| < \epsilon.
  \end{equation*}
\end{cor}

\begin{proof}
  由 \cref{thm:ttac} 存在 $M : \Qp \to \N$，使得 $\bar{x}(\epsilon) \defeq
  x_{M(\epsilon/2)}$ 是 Cauchy 近似。令 $y$ 为其极限（由
  \cref{RD-cauchy-complete} 存在）。给定任意 $\epsilon : \Qp$，令 $n \defeq M(\epsilon/4)$，并注意对任意 $m \geq n$，
  %
  \begin{narrowmultline*}
    |x_m - y| \leq |x_m - x_n| + |x_n - y| =
    |x_m - x_n| + |\bar{x}(\epsilon/2) - y| < \narrowbreak
    \epsilon/4 + \epsilon/2 + \epsilon/4 = \epsilon.\qedhere
  \end{narrowmultline*}
\end{proof}

\subsection{Dedekind 实数是 Dedekind 完备的}
\label{sec:RD-dedekind-complete}

我们把 \RD 构造为 \Q 上的 Dedekind 割类型。但也可以从任意阿基米德有序域 $F$ 出发，在 $F$ 上构造 Dedekind 割\index{cut!Dedekind}。它们仍构成一个阿基米德有序域 $\bar{F}$，称为\define{$F$ 的 Dedekind 完备化}，%
\index{completion!Dedekind}%
\indexsee{Dedekind!completion}{completion, Dedekind}
并且 $F$ 作为子域嵌入其中。若把该构造应用于
\RD，会不会得到更多实数？答案是否定的。事实上我们将证明更强结论：\RD 是终对象。

称有序域~$F$ 对 $\Omega$\define{可容许}
\indexsee{admissible!ordered field}{ordered field, admissible}%
\indexdef{ordered field!admissible}%
，若其严格序
$<$ 是映射 ${<} : F \to F \to \Omega$。

\begin{thm} \label{RD-final-field}
  每个对 $\Omega$ 可容许的阿基米德有序域都可嵌入~\RD 作为子域。
\end{thm}

\begin{proof}
  设 $F$ 是阿基米德有序域。对每个 $x : F$，定义 $L_x, U_x : \Q \to
  \Omega$ 为
  %
  \begin{equation*}
    L_x(q) \defeq (q < x)
    \qquad\text{且}\qquad
    U_x(q) \defeq (x < q).
  \end{equation*}
  %
  （这里用到了 $F$ 对 $\Omega$ 可容许的假设。）
  则 $(L_x, U_x)$ 是 Dedekind 割。\index{cut!Dedekind}确实：割可居住且圆整，因为
  $F$ 是阿基米德的且 $<$ 传递；不相交，因为 $<$ 反自反；定位，因为 $<$ 弱线性。令 $e : F \to \RD$ 为映射 $e(x) \defeq (L_x,
  U_x)$。

  我们断言 $e$ 是保持并反映次序的域嵌入。首先注意，对有理数 $q$ 有 $e(q) = q$。接着对所有 $x, y : F$，有等价
  %
  \begin{narrowmultline*}
    x < y \Leftrightarrow
    (\exis{q : \Q} x < q < y) \Leftrightarrow \narrowbreak
    (\exis{q : \Q} U_x(q) \land L_y(q)) \Leftrightarrow
    e(x) < e(y),
  \end{narrowmultline*}
  %
  因此 $e$ 的确保持并反映次序。$e(x + y) = e(x) + e(y)$ 成立，因为对所有 $q : \Q$，
  %
  \begin{equation*}
    q < x + y \Leftrightarrow
    \exis{r, s : \Q} r < x \land s < y \land q = r + s.
  \end{equation*}
  %
  右推左显然。左推右方面，若 $q < x +
  y$，则仅存在 $r : \Q$ 使 $q - y < r < x$；取 $s \defeq
  q - r$ 即得所需 $r,s$。$e$ 保持乘法留作练习。
\end{proof}

要证明在 \RD 上取 Dedekind 割不会得到新对象，只需再用一个引理。

\begin{lem} \label{lem:cuts-preserve-admissibility}
  若 $F$ 对 $\Omega$ 可容许，则其 Dedekind 完备化也可容许。
  \index{completion!Dedekind}%
\end{lem}

\begin{proof}
  设 $\bar{F}$ 是 $F$ 的 Dedekind 完备化。$\bar{F}$ 上的严格序定义为
  %
  \begin{equation*}
    ((L,U) < (L',U')) \defeq \exis{q : \Q} U(q) \land L'(q).
  \end{equation*}
  %
  由于 $U(q)$ 与 $L'(q)$ 都是 $\Omega$ 的元素，只要 $\Omega$
  对合取与可数存在封闭，该引理即成立；这正是我们一开始所假设的。
\end{proof}


\begin{cor} \label{RD-dedekind-complete}
  %
  \indexdef{complete!ordered field, Dedekind}%
  \indexdef{Dedekind!completeness}%
  Dedekind 实数是 Dedekind 完备的：对每个实值 Dedekind 割 $(L, U)$，
  存在唯一 $x : \RD$ 使得对所有 $y :
  \RD$，有 $L(y) = (y < x)$ 且 $U(y) = (x < y)$。
\end{cor}

\begin{proof}
  由 \cref{lem:cuts-preserve-admissibility}，\RD 的 Dedekind 完备化 $\barRD$
  对 $\Omega$ 可容许，因此由 \cref{RD-final-field} 有嵌入 $\barRD
  \to \RD$，也有嵌入 $\RD \to \barRD$。但这些嵌入必为同构，因为其复合是保持次序的域同态\index{homomorphism!field}，并固定稠密子域~\Q，这意味着它们是恒等。于是该推论立刻由 $\barRD \to \RD$ 是同构得到。
\end{proof}

\index{real numbers!Dedekind|)}%

\section{Cauchy 实数}
\label{sec:cauchy-reals}

\index{real numbers!Cauchy|(}%
\index{completion!Cauchy|(}%
\indexsee{Cauchy!completion}{completion, Cauchy}%
按其本意，Cauchy 实数是 \Q 在 Cauchy 序列极限下的完备化。\index{Cauchy!sequence}
在经典的 Cauchy 实数构造中，我们先考虑 \Q 上所有 Cauchy 序列构成的集合 $\mathcal{C}$，再取一个合适的商 $\mathcal{C}/{\approx}$。
随后，为了说明 $\mathcal{C}/{\approx}$ 是 Cauchy 完备的，我们取一个 Cauchy 序列 $x : \N \to \mathcal{C}/{\approx}$，把它提升为序列的序列 $\bar{x} : \N \to \mathcal{C}$，并利用 $\bar{x}$ 构造 $x$ 的极限。然而，把~$x$ 提升到 $\bar{x}$ 需要
可数选择公理（即~\eqref{eq:ac} 中 $X=\N$ 的情形）或排中律，而我们可能希望避免这些假设。
\indexdef{axiom!of choice!countable}%
凡是以取商作为最后一步的实数构造，都会遇到这一缺陷。
在构造性数学中，常见的应对方式有三种：
\index{mathematics!constructive}%
%
\index{bargaining}%
\begin{enumerate}
\item 假装实数是一个 setoid $(\mathcal{C}, {\approx})$，也就是说，把
  Cauchy 序列类型 $\mathcal{C}$ 连同其重合\index{coincidence, of Cauchy approximations}关系
  以管理性规定附着在一起。这样，实数序列就直接\emph{是}一个代表它们的 Cauchy
  序列序列。
\item 屈从诱惑并接受可数选择公理。毕竟，该公理在多数基于计算观点的构造性数学模型中成立，
  例如可实现性模型。
\item 认为 Cauchy 实数不够理想，转而构造 Dedekind 实数。
  在某些语境下（如层论模型中的构造性数学）这种判断完全正当。
  不过，正如我们在 \cref{sec:dedekind-reals} 所见，构造性的 Dedekind 实数也有它们自己的问题。
\end{enumerate}

不过，借助高阶归纳类型，还有第四种解法。我们认为它优于上述任一方案，而且对经典数学家也同样有趣。
其思想是：Cauchy 实数应当是由~\Q 生成的\emph{自由完备度量空间}。\index{free!complete metric space}
一般地，构造任意种类的自由对象，都需要把该对象的运算反复施加到生成元上许多次。
例如，集合 $X$ 上自由群的元素不只是 $X$ 中元素做一次二元乘积与取逆，而是通过反复迭代乘积和取逆得到的词。
因此，对于 Cauchy 完备化，我们自然也应有同样现象，其中相关“运算”就是“取一个 Cauchy 序列的极限”。
（在这里，迭代必须是超限的，因为即使经过无穷多步，仍会出现新的 Cauchy 序列可取极限。）

上文提到的论证表明：如果排中律或可数选择成立，则 Cauchy 完备化非常特殊；构造一个空间的完备化时，把该运算施加\emph{一步}就够了。
这可类比于自由幺半群和自由群能用（约化）词给出显式描述这一事实。
不过我们在 \cref{sec:free-algebras} 已见，高阶归纳类型允许我们\emph{直接}构造自由对象，而不必依赖是否存在显式描述。
本节将说明，对 Cauchy 实数同样如此（类似技巧可构造任意度量空间的 Cauchy 完备化；见 \cref{ex:metric-completion}）。
具体而言，高阶归纳类型允许我们\emph{同时}添加 Cauchy 序列的极限并对重合关系取商，从而避免“把实数序列提升为代表元序列”的问题。
\index{completion!Cauchy|)}%


\subsection{Cauchy 实数的构造}
\label{sec:constr-cauchy-reals}

把 $\RC$ 作为高阶归纳类型来构造，比 \cref{sec:free-algebras} 中自由代数结构的构造更微妙一些。
我们希望加入一个“取极限”的构造子，其输入是一个实数的 Cauchy 序列；但“实数的 Cauchy 序列”这一概念依赖于某种“距离”测度。
一般而言，两实数之间的距离本身又是一个实数，这会导致潜在的循环定义问题。

然而，对“实数的 Cauchy 序列”我们真正需要的并非一般“距离”概念，而是：对任意 $\epsilon:\Qp$，能够表述“两个实数之间的距离小于 $\epsilon$”。
这可以由一族二元关系表示，记作 $\mathord{\close\epsilon} : \RC\to\RC\to \prop$。
$x \close\epsilon y$ 的意图是 $|x - y| < \epsilon$；但由于我们此时还没有减法、绝对值或不等式的概念（毕竟我们正在定义 $\RC$ 本身），就必须在定义 $\RC$ 的同时定义这些关系 $\close\epsilon$。
并且由于 $\close\epsilon$ 是以两个 $\RC$ 为索引的类型族，普通的互递（高阶）归纳定义不够用；我们需要的是\emph{高阶归纳-归纳定义}。
\index{inductive-inductive type!higher}

回忆 \cref{sec:generalizations}：普通归纳-归纳定义允许我们通过同时归纳，定义一个类型及其以该类型为索引的类型族。
显然，“高阶”版本则允许类型与类型族都含有路径构造子以及点构造子。
我们不打算在此建立高阶归纳-归纳定义的一般理论，但希望接下来对 $\RC$ 与 $\close\epsilon$ 的描述足以使其思想清晰。

\begin{rmk}
  我们也可以考虑\emph{高阶归纳-递归定义}：在对 $\RC$ 做\emph{归纳}定义的同时，用 $\RC$ 的\emph{递归}原理定义 $\close\epsilon$。
  之所以改走归纳-归纳路线，有两个原因。
  第一，从同伦语义角度看，高阶归纳-递归定义似乎更难正当化。
  第二，也是更重要的，归纳-归纳定义给出更强的归纳原理，而我们要发展 Cauchy 实数哪怕最基础的理论都需要它。
\end{rmk}

最后，和 \cref{sec:RD-cauchy-complete} 中讨论 Dedekind 实数的 Cauchy 完备性时一样，我们将使用\emph{Cauchy 近似}（\cref{defn:cauchy-approximation}），而不是 Cauchy 序列。
当然，这里我们的 Cauchy 近似由 Cauchy 实数组成，而不是 Dedekind 实数或有理数。

\begin{defn}\label{defn:cauchy-reals}
  设 $\RC$ 与关系 $\closesym:\Qp \times \RC \times \RC \to \type$ 为如下高阶归纳-归纳类型族。
  \define{Cauchy 实数}类型 $\RC$
  \indexdef{real numbers!Cauchy}%
  \indexsee{Cauchy!real numbers}{real numbers, Cau\-chy}%
  由以下构造子生成：
  \begin{itemize}
  \item \emph{有理点：}
    对任意 $q : \Q$，有一个实数 $\rcrat(q)$。
    \index{rational numbers!as Cauchy real numbers}%
  \item \emph{极限点：}
    对任意 $x : \Qp \to \RC$，若
    %
    \begin{equation}
      \label{eq:RC-cauchy}
      \fall{\delta, \epsilon : \Qp} x_\delta \close{\delta + \epsilon} x_\epsilon
    \end{equation}
    %
    则有一点 $\rclim(x) : \RC$。我们称 $x$ 为一个 \define{Cauchy 近似}。
    \indexdef{Cauchy!approximation}%
    \index{limit!of a Cauchy approximation}%
    %
  \item \emph{路径：}
    对 $u, v : \RC$，若
    %
    \begin{equation}
      \label{eq:RC-path}
      \fall{\epsilon : \Qp} u \close\epsilon v
    \end{equation}
    %
    则有路径 $\rceq(u, v) : \id[\RC]{u}{v}$。
  \end{itemize}
  同时，类型族 $\closesym:\RC\to\RC\to\Qp \to\type$ 由以下构造子生成。
  其中 $q$ 与 $r$ 表示有理数；$\delta$、$\epsilon$ 与 $\eta$ 表示正有理数；$u$ 与 $v$ 表示 Cauchy 实数；$x$ 与 $y$ 表示 Cauchy 近似：
  \begin{itemize}
  \item 对任意 $q,r,\epsilon$，若 $-\epsilon < q - r < \epsilon$，则 $\rcrat(q) \close\epsilon \rcrat(r)$，
  \item 对任意 $q,y,\epsilon,\delta$，若 $\rcrat(q) \close{\epsilon - \delta} y_\delta$，则 $\rcrat(q) \close{\epsilon} \rclim(y)$，
  \item 对任意 $x,r,\epsilon,\delta$，若 $x_\delta \close{\epsilon - \delta} \rcrat(r)$，则 $\rclim(x) \close\epsilon \rcrat(r)$，
  \item 对任意 $x,y,\epsilon,\delta,\eta$，若 $x_\delta \close{\epsilon - \delta - \eta} y_\eta$，则 $\rclim(x) \close\epsilon \rclim(y)$，
  \item 对任意 $u,v,\epsilon$，若 $\xi,\zeta : u \close{\epsilon} v$，则 $\xi=\zeta$（命题截断）。
  \end{itemize}
\end{defn}

\mentalpause

$\RC$ 的第一个构造子说，任何有理数都可以视作实数。
第二个构造子说，从任意一个实数的 Cauchy 近似都能得到一个新的实数，称其“极限”。
第三个构造子表达的是：若两个 Cauchy 近似重合，则它们的极限相等。

$\closesym$ 的前四个构造子分别规定：何时两个有理数接近、何时有理数与极限接近、何时两个极限接近。
在两个有理数的情形下，这就是有理数通常的 $\epsilon$-接近概念；其余情形可由如下直觉导出：每个近似项 $x_\delta$ 预期都在距离极限 $\rclim(x)$ 的 $\delta$ 之内。

我们提醒自己注意证明相关性：由 $\rclim$ 得到的实数表示时不仅包含一个 Cauchy 近似 $x$，还包含~\eqref{eq:RC-cauchy} 的一个证明 $p$，
因此严格写法应是 $\rclim(x,p)$，而不是 $\rclim(x)$。
对 $\rceq$ 与~\eqref{eq:RC-path} 也有类似观察，不过我们仍写
$\rceq : u = v$，而不是 $\rceq(u, v, p) : u = v$。这类记号滥用之所以可接受，
是因为我们省略的是纯命题与易于补全的信息。
同样，$\mathord{\close\epsilon}$ 的最后一个构造子也保证前四个构造子无需命名。

我们立刻就能在 $\RC$ 中填入许多实数。设 $x : \N \to
\Q$ 是一个传统意义下的有理数 Cauchy 序列\index{Cauchy!sequence}，$M : \Qp \to \N$ 是它的
收敛模。则 $\rcrat \circ x \circ M : \Qp \to \RC$ 是一个 Cauchy
近似，所需见证可由 $\closesym$ 的第一个构造子给出。
因此，$\rclim(\rcrat \circ x \circ M)$ 是一个实数。诸如 $\sqrt{2}$、$\pi$、$e$、\dots{} 这类著名
实数都可由这样的有理 Cauchy 序列取极限得到。

\subsection{Cauchy 实数的归纳与递归}
\label{sec:induct-recurs-cauchy}

当然，要对 $\RC$ 做任何有用的事情，我们都需要给出其归纳原理。
与所有“同时归纳定义两个或更多对象”的情形一样，$\RC$ 与 $\closesym$ 的基本归纳原理需要对二者同时归纳。
因此可以预期：给定分别定义在 $\RC$ 与 $\closesym$ 上的两个类型族，以及与各构造子对应的数据，就能得到这两个类型族的截面。
不过由于 $\closesym$ 以两个 $\RC$ 为索引，这两个类型族的精确依赖关系会稍微微妙。
归纳原理将适用于任意如下类型族对：
\begin{align*}
A&:\RC\to\type\\
B&:\prd{x,y:\RC} A(x) \to A(y) \to \prd{\epsilon:\Qp} (x\close\epsilon y) \to \type.
\end{align*}
$A$ 的类型显然；而 $B$ 的类型需要稍加思考。
因为 $B$ 必须依赖于 $\closesym$，而 $\closesym$ 又依赖于两个 $\RC$ 与一个 $\Qp$，所以显然 $B$ 也必须依赖变量 $x,y:\RC$ 与 $\epsilon:\Qp$，以及一个 $(x\close\epsilon y)$ 的元素。
不那么显然的是，$B$ 还必须依赖 $A(x)$ 与 $A(y)$。

若看非依值情形（递归原理）这点会更清楚：此时 $A$ 是一个简单类型（而非类型族）。
在这种情况下，我们会期望 $B$ 不再依赖 $x,y:\RC$ 或 $x\close\epsilon y$。
但递归原理（及其对应的唯一性原理）应当表达：带 $\close\epsilon$ 的 $\RC$ 是某个范畴中的“初对象”，所以此时 $A$ 与 $B$ 的依赖结构应镜像 $\RC$ 与 $\close\epsilon$ 的结构，即应有 $B:A\to A\to \Qp \to \type$。
把这一观察与“在依值情形下 $B$ 还必须依赖 $x,y:\RC$ 及 $x\close\epsilon y$”结合起来，就必然得到上面给出的 $B$ 的类型。

\symlabel{RC-recursion}
把 $B$ 看成在类型 $A(x)$ 与 $A(y)$ 之间按 $\epsilon$ 索引的一族关系会很有帮助。
据此，可把 $B(x,y,a,b,\epsilon,\xi)$ 记作 $(x,a) \bsim_\epsilon^\xi (y,b)$。
由于 $\xi:x\close\epsilon y$ 一旦存在就是唯一的，我们通常省略它，写作 $(x,a) \bsim_\epsilon (y,b)$；只要记住该关系仅在 $x\close\epsilon y$ 时定义，这样写是无害的。
有时还会进一步简写为 $a\bsim_\epsilon b$，其中 $x,y$ 由 $a,b$ 的类型推断；但在某些地方为清晰起见仍需显式写出它们。

\index{induction principle!for Cauchy reals}%
现在，给定类型族 $A:\RC\to\type$ 与如上关系族 $\bsim$，归纳原理所需前提数据如下（对应 $\RC$ 或 $\closesym$ 的每个构造子）：
\begin{itemize}
\item 对任意 $q : \Q$，给定一个元素 $f_q:A(\rcrat(q))$。
\item 对任意 Cauchy 近似 $x$，以及任意 $a:\prd{\epsilon:\Qp} A(x_\epsilon)$，若
  \begin{equation}
    \fall{\delta, \epsilon : \Qp}
    (x_\delta,a_\delta) \bsim_{\delta+\epsilon} (x_\epsilon,a_\epsilon),
    \label{eq:depCauchyappx}
  \end{equation}
  则给定一个元素 $f_{x,a}:A(\rclim(x))$。
  我们称这样的 $a$ 是定义在 $x$ 上的\define{依值 Cauchy 近似}。
  \indexdef{Cauchy!approximation!dependent}%
  \indexsee{approximation, Cauchy}{Cauchy approximation}%
  \indexdef{dependent!Cauchy approximation}%
\item 对 $u, v : \RC$，若 $h:\fall{\epsilon : \Qp} u \close\epsilon v$，并且对任意 $a:A(u)$ 与 $b:A(v)$ 满足
  $\fall{\epsilon:\Qp} (u,a) \bsim_\epsilon (v,b)$，
  则给定依值路径 $\dpath{A}{\rceq(u,v)}{a}{b}$。
\item 对 $q,r:\Q$ 与 $\epsilon:\Qp$，若 $-\epsilon < q - r < \epsilon$，则有
  \narrowequation{(\rcrat(q),f_q) \bsim_\epsilon (\rcrat(r),f_r).}
\item 对 $q:\Q$、$\delta,\epsilon:\Qp$、Cauchy 近似 $y$，以及定义在 $y$ 上的依值 Cauchy 近似 $b$，若 $\rcrat(q) \close{\epsilon - \delta} y_\delta$，则
  \[(\rcrat(q),f_q) \bsim_{\epsilon-\delta} (y_\delta,b_\delta)
  \;\Rightarrow\;
  (\rcrat(q),f_q) \bsim_\epsilon (\rclim(y),f_{y,b}).\]
\item 同样地，对 $r:\Q$、$\delta,\epsilon:\Qp$、Cauchy 近似 $x$，以及定义在 $x$ 上的依值 Cauchy 近似 $a$，若 $x_\delta \close{\epsilon - \delta} \rcrat(r)$，则
  \[(x_\delta,a_\delta) \bsim_{\epsilon-\delta} (\rcrat(r),f_r)
  \;\Rightarrow\;
  (\rclim(x),f_{x,a}) \bsim_\epsilon (\rcrat(q),f_r).
  \]
\item 对 $\epsilon,\delta,\eta:\Qp$ 与 Cauchy 近似 $x,y$，以及分别定义在 $x,y$ 上的依值 Cauchy 近似 $a,b$，若有 $x_\delta \close{\epsilon - \delta - \eta} y_\eta$，则
  \[ (x_\delta,a_\delta) \bsim_{\epsilon - \delta - \eta} (y_\eta,b_\eta)
  \;\Rightarrow\;
  (\rclim(x),f_{x,a}) \bsim_\epsilon (\rclim(y),f_{y,b}).\]
\item 对 $\epsilon:\Qp$、$x,y:\RC$ 与 $\xi,\zeta:x\close{\epsilon} y$，以及 $a:A(x)$ 与 $b:A(y)$，$(x,a) \bsim_\epsilon^\xi (y,b)$ 与 $(x,a) \bsim_\epsilon^\zeta (y,b)$ 的任意两个元素在 $\xi=\zeta$ 之上依值相等。
  如常见情况，这等价于要求 $\bsim$ 取值于纯命题。
\end{itemize}
在这些前提下，可推出函数
\begin{align*}
  f&:\prd{x:\RC} A(x)\\
  g&:\prd{x,y:\RC}{\epsilon:\Qp}{\xi:x\close{\epsilon} y}
  (x,f(x)) \bsim_\epsilon^\xi (y,f(y))
\end{align*}
并按预期计算：
\begin{align}
  f(\rcrat(q)) &\defeq f_q, \label{eq:rcsimind1}\\
  f(\rclim(x)) &\defeq f_{x,(f,g)[x]}. \label{eq:rcsimind2}
\end{align}
这里 $(f,g)[x]$ 表示把 $f$ 与 $g$ 作用于 Cauchy 近似 $x$，从而得到定义在 $x$ 上的依值 Cauchy 近似。
也就是说，定义 $(f,g)[x]_\epsilon \defeq f(x_\epsilon) : A(x_\epsilon)$；然后对任意 $\epsilon,\delta:\Qp$，由 $x$ 是 Cauchy 近似的假设可得 $\xi: x_\epsilon \close{\epsilon+\delta} x_\delta$，并由 $g(x_\epsilon,x_\delta,\epsilon+\delta,\xi)$ 见证 $(f,g)[x]$ 确为依值 Cauchy 近似。

这个记号之后不会再用，不必刻意记忆。
通常我们采用模式匹配约定：像~\eqref{eq:rcsimind1} 与~\eqref{eq:rcsimind2} 这样的等式定义 $f$，其中~\eqref{eq:rcsimind2} 右端可使用符号 $f(x_\epsilon)$，并假设它们组成一个依值 Cauchy 近似。

然而，这个归纳原理确实仍然相当繁复。
为了帮助理解，我们注意到它包含两个特例：分别针对~$\RC$ 与~$\closesym$ 的归纳原理。
第一，若只给定类型族 $A:\RC\to\type$，并令 $\bsim$ 恒为 \unit，则所需数据中大部分自动平凡化，剩下的是：
\begin{itemize}
\item 对任意 $q : \Q$，给定元素 $f_q:A(\rcrat(q))$；
\item 对任意 Cauchy 近似 $x$，以及任意 $a:\prd{\epsilon:\Qp} A(x_\epsilon)$，给定元素 $f_{x,a}:A(\rclim(x))$；
\item 对 $u, v : \RC$ 与 $h:\fall{\epsilon : \Qp} u \close\epsilon v$，以及 $a:A(u)$、$b:A(v)$，给定 $\dpath{A}{\rceq(u,v)}{a}{b}$。
\end{itemize}
在这些数据下，归纳原理给出函数 $f:\prd{x:\RC} A(x)$，满足
\begin{align*}
  f(\rcrat(q)) &\defeq f_q,\\
  f(\rclim(x)) &\defeq f_{x,f(x)}.
\end{align*}
我们称该原理为\define{$\RC$-归纳}；它本质上说明：若把 $\close\epsilon$ 视为既有结构，则 $\RC$ 由其构造子归纳生成。

注意，若 $A$ 是纯性质，则 $\RC$-归纳中的第三个假设自动成立（稍后会看到这两者实际上等价）。
因此，证明实数的纯性质时，只需分别处理有理点与 Cauchy 近似的极限点。
下面是一个例子。

\begin{lem} \label{lem:close-reflexive}
  对任意 $u:\RC$ 与 $\epsilon:\Qp$，有 $u\close\epsilon u$。
\end{lem}
\begin{proof}
  定义 $A(u) \defeq \fall{\epsilon:\Qp} (u\close\epsilon u)$。
  由于这是纯命题（由 $\closesym$ 的最后一个构造子），由 $\RC$-归纳只需在 $u$ 为 $\rcrat(q)$ 与 $\rclim(x)$ 时证明。
  第一种情形中，对任意 $\epsilon$ 显然有 $|q-q|<\epsilon$，故由 $\closesym$ 的第一个构造子得 $\rcrat(q) \close\epsilon \rcrat(q)$。
  %
  第二种情形中，可归纳假设对所有 $\delta,\epsilon:\Qp$ 有 $x_\delta \close\epsilon x_\delta$。
  特别地，$x_{\epsilon/3} \close{\epsilon/3} x_{\epsilon/3}$，于是由 $\closesym$ 的第四个构造子得 $\rclim(x) \close{\epsilon} \rclim(x)$。
\end{proof}

由 \cref{lem:close-reflexive} 可知：直接应用 $\RC$-归纳要成功，$A:\RC\to\type$ 至少应是纯性质。
为此固定 $u:\RC$。
取 $v$ 为 $u$，$\RC$-归纳的第三个假设告诉我们：对任意 $a : A(u)$，有 $\dpath{A}{\rceq(u,u)}{a}{a}$。
若再给定一点 $b : A(u)$，也得到 $\dpath{A}{\rceq(u,u)}{a}{b}$。
由依值路径类型定义可知，这两条路径合起来推出 $a = b$，即 $A(u)$ 中任意两点相等。

\begin{thm}\label{thm:Cauchy-reals-are-a-set}
  $\RC$ 是集合。
\end{thm}
\begin{proof}
  我们刚刚已经证明纯关系
  \narrowequation{P(u,v) \defeq \fall{\epsilon:\Qp} (u\close\epsilon v)}
  具自反性。
  由 $\RC$ 的路径构造子可知它蕴含恒等，因此结论由 \cref{thm:h-set-refrel-in-paths-sets} 得出。
\end{proof}

我们还可证明：虽然 $\RC$ 未必是\emph{有理数} Cauchy 序列集合的商，但它仍然是\emph{实数} Cauchy 序列集合的商。
（当然，这不能作为 $\RC$ 的\emph{定义}，但它是一个有用性质。）
我们定义 Cauchy 近似的类型为
%
\symlabel{cauchy-approximations}%
\index{Cauchy!approximation!type of}%
\begin{equation*}
  \CAP \defeq
  \setof{ x : \Qp \to \RC |
    \fall{\epsilon, \delta : \Qp} x_\delta \close{\delta + \epsilon} x_\epsilon
  }.
\end{equation*}
$\RC$ 的第二个构造子给出函数 $\rclim:\CAP\to\RC$。

\begin{lem} \label{RC-lim-onto}
  每个实数都仅仅是一个极限点：$\fall{u : \RC} \exis{x : \CAP} u = \rclim(x)$。
  换言之，$\rclim:\CAP\to\RC$ 是满射。
\end{lem}
\begin{proof}
  用 $\RC$-归纳对 $u$ 分类。
  若 $u$ 本来就是极限 $\rclim(x)$，命题显然。
  因此只需设 $u$ 是有理点 $\rcrat(q)$；我们主张 $u$ 等于 $\rclim(\lam{\epsilon} \rcrat(q))$。
  由 $\RC$ 的路径构造子，只需证对所有 $\epsilon:\Qp$ 有 $\rcrat(q) \close\epsilon \rclim(\lam{\epsilon} \rcrat(q))$。
  由 $\closesym$ 的第二个构造子，这又只需找到 $\delta:\Qp$ 使 $\rcrat(q)\close{\epsilon-\delta} \rcrat(q)$。
  而由 $\closesym$ 的第一个构造子，只要取任意满足 $\delta<\epsilon$ 的 $\delta:\Qp$ 即可。
\end{proof}

%

\begin{lem} \label{RC-lim-factor}
  若 $A$ 是集合，且 $f : \CAP \to A$ 保持 Cauchy 近似的重合\index{coincidence!of Cauchy approximations}，即
  %
  \begin{equation*}
    \fall{x, y : \CAP} \rclim(x) = \rclim(y) \Rightarrow f(x) = f(y),
  \end{equation*}
  %
  则 $f$ 唯一地经由 $\rclim : \CAP \to \RC$ 因子化。
\end{lem}
\begin{proof}
  由 $\rclim$ 满射及 \cref{lem:images_are_coequalizers}，$\RC$ 是 $\rclim$ 的核对\index{kernel!pair}在 $\CAP$ 上的商。
  这正是引理所述。
\end{proof}

再看归纳原理的第二个特例：改令 $A$ 恒为 $\unit$。
此时 $\bsim$ 只是定义在 $\epsilon$-接近实数对上的一族按 $\epsilon$ 索引的关系，因此可写 $u\bsim_\epsilon v$，而不写 $(u,\ttt)\bsim_\epsilon (v,\ttt)$。
所需数据化简为下列内容，其中 $q, r$ 为有理数，$\epsilon, \delta, \eta$ 为正有理数，$x, y$ 为 Cauchy 近似：
\begin{itemize}
\item 若 $-\epsilon < q - r < \epsilon$，则
  $\rcrat(q) \bsim_\epsilon \rcrat(r)$，
\item 若 $\rcrat(q) \close{\epsilon - \delta} y_\delta$ 且
  $\rcrat(q)\bsim_{\epsilon-\delta} y_\delta$，
  则 $\rcrat(q) \bsim_\epsilon \rclim(y)$，
\item 若 $x_\delta \close{\epsilon - \delta} \rcrat(r)$ 且
  $x_\delta \bsim_{\epsilon-\delta} \rcrat(r)$，
  则 $\rclim(y) \bsim_\epsilon \rcrat(q)$，
\item 若 $x_\delta \close{\epsilon - \delta - \eta} y_\eta$ 且
  $x_\delta\bsim_{\epsilon - \delta - \eta} y_\eta$，
  则 $\rclim(x) \bsim_\epsilon \rclim(y)$。
\end{itemize}
得到的结论是 $\fall{u,v:\RC}{\epsilon:\Qp} (u\close\epsilon v) \to (u \bsim_\epsilon v)$。
我们称该原理为\define{$\closesym$-归纳}；它本质上说明：若把 $\RC$ 视为既有结构，则 $\close\epsilon$（作为类型族）由\emph{它自己的}构造子归纳生成。
例如，我们可用它证明 $\closesym$ 对称。

\begin{lem}\label{thm:RCsim-symmetric}
  对任意 $u,v:\RC$ 与 $\epsilon:\Qp$，有 $(u\close\epsilon v) = (v\close\epsilon u)$。
\end{lem}
\begin{proof}
  两边都是纯命题，因此由对称性只需证明一个方向。
  定义 $(u\bsim_\epsilon v) \defeq (v\close\epsilon u)$。
  用 $\closesym$-归纳，可化归到 $u\close\epsilon v$ 由 $\closesym$ 四个非平凡构造子之一导出的情形。
  第一种（$u,v$ 都是有理点）是平凡的（再用一次第一个构造子即可）。
  其余三种情形中，归纳假设（加上 \Q 加法交换律）恰好给出另一个 $\closesym$ 构造子的输入（第二与第三互换，第四保持不变）。
\end{proof}

因此，一般归纳原理（可称为\define{$(\RC,\closesym)$-归纳}）是一种联合的 $\RC$-归纳与 $\closesym$-归纳。
尤其是它的非依值版本，我们称为\define{$(\RC,\closesym)$-递归}，将是我们最常用的形式。
\index{recursion principle!for Cauchy reals}%
普通的 $\RC$-递归告诉我们：要定义函数 $f : \RC \to A$，只需：
\begin{enumerate}
\item 对每个 $q : \Q$ 构造 $f(\rcrat(q)) : A$，
\item 对每个 Cauchy 近似 $x : \Qp \to \RC$，构造 $f(x) : A$，
  并假设对所有 $\epsilon : \Qp$，$f(x_\epsilon)$ 已定义，
\item 对所有满足 $\fall{\epsilon:\Qp} u\close\epsilon v$ 的 $u, v : \RC$，证明 $f(u) = f(v)$。\label{item:rcrec3}
\end{enumerate}
然而，在不知道 $f$ 如何作用于 $\epsilon$-接近 Cauchy 实数的情况下，一般很难证明~\ref{item:rcrec3}。
增强版的 $(\RC,\closesym)$-递归正是为弥补此缺陷：它允许我们预先指定一个\emph{任意的}“$f$ 在 $\epsilon$-接近 Cauchy 实数上如何作用”的方式，再通过与定义 $f$ 同步进行的归纳来证明它确实成立。
这正是关系族 $\bsim$ 的角色。
由于 $A$ 与 $\RC$ 无关，为简便起见可设 $\bsim$ 只依赖 $A$ 与 $\Qp$，于是写 $a\bsim_\epsilon b$ 而非 $(u,a) \bsim_\epsilon (v,b)$ 不会有歧义。
在这种情形下，用 $(\RC,\closesym)$-递归定义 $f:\RC\to A$ 需要以下分支（以下采用模式匹配记法）。
\begin{itemize}
\item 对每个 $q : \Q$，构造 $f(\rcrat(q)) : A$。
\item 对每个 Cauchy 近似 $x : \Qp \to \RC$，构造 $f(\rclim(x)) : A$；归纳假设对所有 $\epsilon : \Qp$，$f(x_\epsilon)$ 已定义，并且关于 $\bsim$ 形成一个“Cauchy 近似”，即 $\fall{\epsilon,\delta:\Qp} (f(x_\epsilon) \bsim_{\epsilon+\delta} f(x_\delta))$。
\item 证明关系族 $\bsim$ 是\emph{分离的}，即对任意 $a,b:A$，
  \indexdef{relation!separated family of}%
  \indexdef{separated family of relations}%
\narrowequation{(\fall{\epsilon:\Qp} a\bsim_\epsilon b) \Rightarrow (a=b).}
\item 证明：若 $q,r:\Q$ 满足 $-\epsilon< q-r <\epsilon$，则 $f(\rcrat(q))\bsim_\epsilon f(\rcrat(r))$。
\item 对任意 $q:\Q$ 与任意 Cauchy 近似 $y$，证明
\narrowequation{f(\rcrat(q)) \bsim_\epsilon f(\rclim(y)),} 归纳假设是：对某个 $\delta:\Qp$，有 $\rcrat(q)\close{\epsilon-\delta} y_\delta$ 且 $f(\rcrat(q)) \bsim_{\epsilon-\delta} f(y_\delta)$，并且 $\eta \mapsto f(x_\eta)$ 关于 $\bsim$ 为 Cauchy 近似。
\item 对任意 Cauchy 近似 $x$ 与任意 $r:\Q$，证明
\narrowequation{f(\rclim(x)) \bsim_\epsilon f(\rcrat(r)),}
归纳假设是：对某个 $\delta:\Qp$，有 $x_\delta \close{\epsilon-\delta} \rcrat(r)$ 且 $f(x_\delta) \bsim_{\epsilon-\delta} f(\rcrat(r))$，并且 $\eta\mapsto f(x_\eta)$ 关于 $\bsim$ 为 Cauchy 近似。
\item 对任意 Cauchy 近似 $x,y$，证明
\narrowequation{f(\rclim(x)) \bsim_\epsilon f(\rclim(y)),}
归纳假设是：对某些 $\delta,\eta:\Qp$，有 $x_\delta \close{\epsilon-\delta-\eta} y_\eta$ 且 $f(x_\delta) \bsim_{\epsilon-\delta-\eta} f(y_\eta)$，并且 $\theta\mapsto f(x_\theta)$ 与 $\theta\mapsto f(y_\theta)$ 都关于 $\bsim$ 为 Cauchy 近似。
\end{itemize}
注意，最后四项证明中，我们可以自由使用前两项中给出的 $f(\rcrat(q))$ 与 $f(\rclim(x))$ 的具体定义。
但分离性的证明必须适用于 $A$ 的\emph{任意}两个元素，与 $f$ 无关：它是关系族 $\bsim$ 的一种“可接受性”条件。
因此，我们往往在定义好 $\bsim$ 后先验证它，再继续定义 $f(\rcrat(q))$ 与 $f(\rclim(x))$。

在上述前提下，$(\RC,\closesym)$-递归给出函数 $f:\RC\to A$，并且 $f(\rcrat(q))$ 与 $f(\rclim(x))$ 与前两条中给出的定义在判断上相等。
此外，还可推出
\begin{equation}
  \fall{u,v:\RC}{\epsilon:\Qp} (u\close\epsilon v) \to (f(u) \bsim_\epsilon f(v)).\label{eq:RC-sim-recursion-extra}
\end{equation}

作为典型示例，$(\RC,\closesym)$-递归允许我们把定义在 \Q 上的函数延拓到整个 $\RC$，前提是它们有足够连续性。
\index{function!continuous}%

\begin{defn}\label{defn:lipschitz}
  函数 $f:\Q\to\RC$ 称为 \define{Lipschitz}
  \indexdef{function!Lipschitz}%
  \indexdef{Lipschitz!function}%
  \indexdef{Lipschitz!constant}%
  \indexdef{constant!Lipschitz}%
  若存在 $L:\Qp$（称为\define{Lipschitz 常数}）使得
  \[ |q - r|<\epsilon \Rightarrow (f(q) \close{L\epsilon} f(r)) \]
  对所有 $\epsilon:\Qp$ 与 $q,r:\Q$ 成立。
  %
  同样地，$g:\RC\to\RC$ 称为 \define{Lipschitz}，若存在 $L:\Qp$ 使
  \[ (u\close\epsilon v) \Rightarrow (g(u) \close{L\epsilon} g(v)) \]
  对所有 $\epsilon:\Qp$ 与 $u,v:\RC$ 成立。
\end{defn}

特别地，由 $\closesym$ 的第一个构造子可见：若 $f:\Q\to\Q$ 在通常意义下是 Lipschitz，则复合映射 $\Q\xrightarrow{f} \Q \to \RC$ 也 Lipschitz。

\begin{lem}\label{RC-extend-Q-Lipschitz}
  设 $f : \Q \to \RC$ 是 Lipschitz，常数为 $L : \Qp$。
  则存在 Lipschitz 映射 $\bar{f} : \RC \to \RC$（常数同样为 $L$），并且对所有 $q:\Q$ 有 $\bar{f}(\rcrat(q)) \jdeq f(q)$。
\end{lem}

\begin{proof}
  % Uniqueness follows directly from \cref{RC-continuous-eq}.
  我们用 $(\RC,\closesym)$-递归定义 $\bar{f}$，其陪域取 $A\defeq \RC$。
  定义关系 $\mathord{\bsim}: \RC \to \RC \to \Qp \to \prop$ 为
  \begin{align*}
    (u \bsim_\epsilon v) &\defeq (u \close{L\epsilon} v).
  \end{align*}
  对 $q : \Q$，定义
  %
  \begin{equation*}
    \bar{f}(\rcrat(q)) \defeq \rcrat(f(q)).
  \end{equation*}
  %
  对 Cauchy 近似 $x : \Qp \to \RC$，定义
  %
  \begin{equation*}
    \bar{f}(\rclim(x)) \defeq \rclim(\lamu{\epsilon : \Qp} \bar{f}(x_{\epsilon/L})).
  \end{equation*}
  %
  为使该定义成立，须验证 $y \defeq \lamu{\epsilon : \Qp} \bar{f}(x_{\epsilon/L})$ 是 Cauchy 近似。
  但该步归纳假设给出：对任意 $\delta,\epsilon:\Qp$，有 $\bar{f}(x_\delta) \bsim_{\delta+\epsilon} \bar{f}(x_\epsilon)$，即 $\bar{f}(x_\delta) \close{L\delta+L\epsilon} \bar{f}(x_\epsilon)$。
  因而有
  \[y_\delta \jdeq f(x_{\delta/L}) \close{\delta + \epsilon} f(x_{\epsilon/L})   \jdeq y_\epsilon. \]

  证明分离性时，只需注意：$\fall{\epsilon:\Qp} a\bsim_\epsilon b$ 意味着 $\fall{\epsilon:\Qp} a\close{L\epsilon} b$，这蕴含 $\fall{\epsilon:\Qp}a\close\epsilon b$，从而 $a=b$。

  为完成 $(\RC,\closesym)$-递归，还需验证 $\bsim$ 的四个条件。
  这基本等同于在 $\closesym$ 的四个构造子情形下证明 $\bar f$ 都是 Lipschitz。
  \begin{enumerate}
  \item 当 $u$ 为 $\rcrat(q)$、$v$ 为 $\rcrat(r)$，且 $-\epsilon < |q-r| <\epsilon$ 时，$f$ 的 Lipschitz 性给出 $f(q) \close{L\epsilon} f(r)$，故按定义有 $\bar{f}(\rcrat(q)) \bsim_\epsilon \bar{f}(\rcrat(r))$。
  \item 当 $u$ 为 $\rclim(x)$、$v$ 为 $\rcrat(q)$ 且 $x_\eta \close{\epsilon - \eta} \rcrat(q)$ 时，
      归纳假设为 $\bar{f}(x_\eta) \close{L \epsilon - L \eta} \rcrat(f(q))$，据此由 $\closesym$ 的第三个构造子可得
      \narrowequation{\bar{f}(\rclim(x)) \close{L \epsilon} \bar{f}(\rcrat(q))}
  \item 对称情形（$u$ 为有理点、$v$ 为极限点）本质相同。
  \item 当 $u$ 为 $\rclim(x)$、$v$ 为 $\rclim(y)$，且存在 $\delta, \eta : \Qp$ 使 $x_\delta \close{\epsilon - \delta - \eta} y_\eta$ 时，
      归纳假设为 $\bar{f}(x_\delta) \close{L \epsilon - L \delta - L \eta} \bar{f}(y_\eta)$，由 $\closesym$ 的第四个构造子得 $\bar{f}(\rclim(x)) \close{L
        \epsilon} \bar{f}(\rclim(y))$。
  \end{enumerate}
  这就完成了 $(\RC,\closesym)$-递归，也即 $\bar f$ 的构造。
  所需等式 $\bar f(\rcrat(q))\jdeq f(q)$ 正是 $(\RC,\closesym)$-递归的第一条计算规则，而附加结论~\eqref{eq:RC-sim-recursion-extra} 恰好说明 $\bar f$ 以 $L$ 为 Lipschitz 常数。
\end{proof}

到这里，在没有更好刻画 $\closesym$ 之前，我们能推进的基本都做完了。
在 $\closesym$ 的构造子里，我们规定了两种不同形态的 Cauchy 实数在何种条件下应当 $\epsilon$-接近。
但我们如何知道，在最终得到的归纳-归纳类型族中，这些就是该事实的\emph{全部}见证？
我们已经看到，归纳类型族（如恒等类型，见 \cref{sec:identity-systems}）和高阶归纳类型常会“长出比放进去更多的内容”，因此这并非杞人忧天。

为了更精确刻画 $\closesym$，我们将\emph{递归地}定义 $\RC$ 上的一族关系 $\approx_\epsilon$，使其在构造子上可计算，并证明这族关系与 $\close\epsilon$ 等价。

\begin{thm}\label{defn:RC-approx}
  存在一族纯关系 $\mathord\approx:\RC\to\RC\to\Qp\to\prop$，满足
  \begin{align}
    (\rcrat(q) \approx_\epsilon \rcrat(r))  &\defeq
    (-\epsilon < q - r < \epsilon)\label{eq:RCappx1}\\
    (\rcrat(q) \approx_\epsilon \rclim(y)) &\defeq
    \exis{\delta : \Qp} \rcrat(q) \approx_{\epsilon - \delta} y_\delta\label{eq:RCappx2}\\
    (\rclim(x) \approx_\epsilon \rcrat(r)) &\defeq
    \exis{\delta : \Qp} x_\delta \approx_{\epsilon - \delta} \rcrat(r)\label{eq:RCappx3}\\
    (\rclim(x) \approx_\epsilon \rclim(y)) &\defeq
    \exis{\delta, \eta : \Qp} x_\delta \approx_{\epsilon - \delta - \eta} y_\eta.\label{eq:RCappx4}
  \end{align}
  此外，还有
  \begin{gather}
    (u \approx_\epsilon v) \Leftrightarrow \exis{\theta:\Qp} (u \approx_{\epsilon-\theta} v) \label{RC-sim-rounded}\\
    (u \approx_\epsilon v) \to (v\close\delta w) \to (u\approx_{\epsilon+\delta} w)\label{eq:RC-sim-rtri}\\
    (u \close\epsilon v) \to (v\approx_\delta w) \to (u\approx_{\epsilon+\delta} w)\label{eq:RC-sim-ltri}.
  \end{gather}
\end{thm}

附加条件~\eqref{RC-sim-rounded}--\eqref{eq:RC-sim-ltri} 实际上是让归纳定义顺利进行所必需的。
条件~\eqref{RC-sim-rounded} 称为\define{圆整（rounded）}。
\indexsee{relation!rounded}{rounded relation}%
\indexdef{rounded!relation}%
把它从右往左读得到 $\approx$ 的\define{单调性}，
\index{monotonicity}%
\index{relation!monotonic}%
%
\begin{equation*}
  (\delta < \epsilon) \land (u \approx_\delta v) \Rightarrow (u \approx_\epsilon v)
\end{equation*}
%
而从左往右读得到 $\approx$ 的\define{开性}，
\index{open!relation}%
\index{relation!open}%
%
\begin{equation*}
  (u \approx_\epsilon v) \Rightarrow \exis{\delta : \Qp} (\delta < \epsilon) \land (u \approx_\delta v).
\end{equation*}
%
条件~\eqref{eq:RC-sim-rtri} 与~\eqref{eq:RC-sim-ltri} 是三角不等式的两种形式，表示 $\approx$ 在左右两侧都是 $\closesym$ 的一个“模”。

\begin{proof}
  我们将用双重 $(\RC,\closesym)$-递归来定义 $\mathord\approx:\RC\to\RC\to\Qp\to\prop$。
  首先应用一次 $(\RC,\closesym)$-递归，陪域取为 $\RC\to\Qp\to\prop$ 的子集，其中元素是圆整并满足某一侧三角不等式的谓词族。
  把这些谓词看作二元关系的一半，我们记作 $(u,\epsilon) \mapsto (\hapx_\epsilon u)$，符号 $\hapname$ 指整条关系。
  于是可把 $A$ 精确写为
  \begin{multline*}
    A \defeq\; \Bigg\{ \hapname :\RC\to\Qp\to\prop \;\bigg|\; \\
    \Big(\fall{u:\RC}{\epsilon:\Qp}
    \big((\hapx_\epsilon u) \Leftrightarrow \exis{\theta:\Qp} (\hapx_{\epsilon-\theta} u)\big)\Big)  \\
    \land \Big(\fall{u,v:\RC}{\eta,\epsilon:\Qp} (u\close\epsilon v) \to\\
    \big((\hapx_\eta u) \to (\hapx_{\eta+\epsilon} v) \big) \land \big((\hapx_\eta v) \to (\hapx_{\eta+\epsilon} u) \big)\Big)\Bigg\}
  \end{multline*}
  与子集类型的惯例一样，我们对 $A$ 的一个元素与其第一分量 $\hapname$ 使用同一记号。
  作为 $(\RC,\closesym)$-递归所需的关系族，我们取下式，它将保证另一侧三角不等式：
  \begin{narrowmultline*}
    (\hapname \bsim_\epsilon \hapbname ) \defeq \narrowbreak
    \fall{u:\RC}{\eta:\Qp} ((\hapx_\eta u) \to (\hapxb_{\epsilon+\eta} u))
    \land \narrowbreak
    ((\hapxb_\eta u) \to (\hapx_{\epsilon+\eta} u)).
  \end{narrowmultline*}
  可见这些关系是分离的。
  假设
  \narrowequation{\fall{\epsilon:\Qp} (\hapname \bsim_\epsilon \hapbname),}
  要证 $\hapname = \hapbname$，只需证对所有 $u:\RC$ 有 $(\hapx_\epsilon u) \Leftrightarrow (\hapxb_\epsilon u)$。
  但由圆整性，$\hapx_\epsilon u$ 蕴含对某个 $\theta$ 有 $\hapx_{\epsilon-\theta} u$，再结合 $\hapname \bsim_\epsilon \hapbname$ 得 $\hapxb_\epsilon u$；反向同理。

  现在递归原理要求的前两项数据是：
  \begin{itemize}
  \item 对任意 $q:\Q$，给出 $A$ 的一个元素，记为 $(\rcrat(q)\approx_{(\blank)} \blank)$。
  \item 对任意 Cauchy 近似 $x$，若已定义函数 $\Qp \to A$，记作 $\epsilon \mapsto (x_\epsilon \approx_{(\blank)} \blank)$，并满足
    % \[ \fall{u,v:\RC}{\delta,\epsilon,\eta:\Qp} (x_\delta \approx_\eta u) \to (u\close{\delta+\epsilon} v) \to (x_\epsilon \approx_{\eta+\delta+\epsilon} v) \]
    \begin{equation}
      \fall{u:\RC}{\delta,\epsilon,\eta:\Qp} (x_\delta \approx_\eta u) \to (x_\epsilon \approx_{\eta+\delta+\epsilon} u),\label{eq:appxrec2}
    \end{equation}
    则需给出 $A$ 的一个元素，记作 $(\rclim(x)\approx_{(\blank)} \blank)$。
  \end{itemize}
  在这两种情形中，我们都通过嵌套一次 $(\RC,\closesym)$-递归来给出定义，陪域取 $\Qp\to\prop$ 的子集，其中元素是纯命题的圆整族。
  把这些命题看作二元关系的“零半边”，记作 $\epsilon \mapsto (\tap{\epsilon})$，符号 $\tapname$ 表示整族。
  于是内层递归的陪域可精确写为
  \begin{narrowmultline*}
    C \defeq
    \bigg\{ \tapname :\Qp\to\prop \;\;\Big|\;\; \narrowbreak
    \fall{\epsilon:\Qp} \Big((\tap\epsilon) \Leftrightarrow \exis{\theta:\Qp} (\tap{\epsilon-\theta})\Big)\bigg\}
  \end{narrowmultline*}
  所需关系族取为三角不等式剩余部分：
  \begin{narrowmultline*}
    (\tapname \bbsim_\epsilon \tapbname) \defeq
    \fall{\eta:\Qp} ((\tap\eta) \to (\tapb{\epsilon+\eta})) \land
    \narrowbreak
    ((\tapb\eta) \to (\tap{\epsilon+\eta})).
  \end{narrowmultline*}
  这些关系与 $\bsim$ 一样可由同样论证得到分离性，依赖于 $C$ 中元素都圆整。

  注意，若这样的内层递归成功，它会产生谓词族 $\hapname : \RC\to\Qp\to \prop$，其圆整性来自它在 $\Qp\to\prop$ 中的像落在 $C$，并满足
  \[ \fall{u,v:\RC}{\epsilon:\Qp} (u\close\epsilon v) \to \big((\hapx_{(\blank)} u) \bbsim_\epsilon (\hapx_{(\blank)} u)\big). \]
  展开 $\bbsim$ 的定义，恰好得到 $\hapname$ 属于 $A$ 所需的第三个条件；这正是我们需要的。

  到这里，我们即可给出定义~\eqref{eq:RCappx1}--\eqref{eq:RCappx4}，它们作为两个内层递归各自的前两条分支（对应有理点和极限点）。
  在每种情形下都必须验证该关系是圆整的，从而属于 $C$。
  在有理-有理情形~\eqref{eq:RCappx1} 这很明显；其余情形由归纳假设给出。
  （在~\eqref{eq:RCappx2} 中，相关归纳假设是 $(\rcrat(q) \approx_{(\blank)} y_\delta) : C$；在~\eqref{eq:RCappx3} 与~\eqref{eq:RCappx4} 中，相关归纳假设是 $(x_\delta \approx_{(\blank)} \blank) : A$。）

  子递归剩余的数据是证明 \eqref{eq:RCappx1}--\eqref{eq:RCappx4} 相对于 $\closesym$ 构造子满足右侧三角不等式。
  共有八种情形——每个子递归四种——对应于~\eqref{eq:RC-sim-rtri} 中 $u,v,w$ 取有理点或极限点的八种组合。
  先看 $u$ 为 $\rcrat(q)$ 的情形。
  \begin{enumerate}
  \item 假设 $\rcrat(q)\approx_\phi \rcrat(r)$ 且 $-\epsilon<|r-s|<\epsilon$，需证 $\rcrat(q)\approx_{\phi+\epsilon} \rcrat(s)$。
    但按 $\approx$ 的定义，这直接化为有理数的三角不等式。
  \item 假设 $\phi,\epsilon,\delta:\Qp$ 满足 $\rcrat(q)\approx_\phi \rcrat(r)$ 且 $\rcrat(r) \close{\epsilon-\delta} y_\delta$，并归纳地有
    \begin{equation}
      \fall{\psi:\Qp}(\rcrat(q) \approx_{\psi} \rcrat(r)) \to (\rcrat(q) \approx_{\psi+\epsilon-\delta} y_\delta).\label{eq:RCappx-rtri-rrl1}
    \end{equation}
    还假设 $\psi,\delta\mapsto (\rcrat(q) \approx_{\psi} y_\delta)$ 关于 $\bbsim$ 构成 Cauchy 近似，即
    \begin{equation}
      \fall{\psi,\xi,\zeta:\Qp} (\rcrat(q) \approx_{\psi} y_\xi) \to (\rcrat(q) \approx_{\psi+\xi+\zeta} y_\zeta),\label{eq:RCappx-rtri-rrl2}
    \end{equation}
    尽管本情形并不需要此假设。
    实际上，取 $\psi\defeq \phi$ 代入 \eqref{eq:RCappx-rtri-rrl1} 立得 $\rcrat(q) \approx_{\phi+\epsilon-\delta} y_\delta$，故按 $\approx$ 定义有 $\rcrat(q) \approx_{\phi+\epsilon} \rclim(y)$。
  \item 假设 $\phi,\epsilon,\delta:\Qp$ 满足 $\rcrat(q)\approx_\phi \rclim(y)$ 与 $y_\delta \close{\epsilon-\delta} \rcrat(r)$，并归纳地有
    \begin{gather}
      \fall{\psi:\Qp}(\rcrat(q) \approx_{\psi} y_\delta) \to (\rcrat(q) \approx_{\psi+\epsilon-\delta} \rcrat(r)).\label{eq:RCappx-rtri-rlr1}\\
      \fall{\psi,\xi,\zeta:\Qp} (\rcrat(q) \approx_{\psi} y_\xi) \to (\rcrat(q) \approx_{\psi+\xi+\zeta} y_\zeta).\label{eq:RCappx-rtri-rlr2}
    \end{gather}
    由定义，$\rcrat(q)\approx_\phi \rclim(y)$ 意味着存在 $\theta:\Qp$ 使 $\rcrat(q) \approx_{\phi-\theta} y_\theta$。
    因此由假设~\eqref{eq:RCappx-rtri-rlr2} 可得 $\rcrat(q) \approx_{\phi+\delta} y_\delta$，再由~\eqref{eq:RCappx-rtri-rlr1} 推出 $\rcrat(q) \approx_{\phi+\epsilon} \rcrat(r)$，如愿。
  \item 假设 $\phi,\epsilon,\delta,\eta:\Qp$ 满足 $\rcrat(q)\approx_\phi \rclim(y)$ 与 $y_\delta \close{\epsilon-\delta-\eta} z_\eta$，并归纳地有
    \begin{gather}
      \fall{\psi:\Qp}(\rcrat(q) \approx_{\psi} y_\delta) \to (\rcrat(q) \approx_{\psi+\epsilon-\delta-\eta} z_\eta), \label{eq:RCappx-rtri-rll1}\\
      \fall{\psi,\xi,\zeta:\Qp} (\rcrat(q) \approx_{\psi} y_\xi) \to (\rcrat(q) \approx_{\psi+\xi+\zeta} y_\zeta), \label{eq:RCappx-rtri-rll2}\\
      \fall{\psi,\xi,\zeta:\Qp} (\rcrat(q) \approx_{\psi} z_\xi) \to (\rcrat(q) \approx_{\psi+\xi+\zeta} z_\zeta). \label{eq:RCappx-rtri-rll3}
    \end{gather}
    同样，$\rcrat(q)\approx_\phi \rclim(y)$ 意味着存在 $\xi:\Qp$ 使 $\rcrat(q) \approx_{\phi-\xi} y_\xi$；由~\eqref{eq:RCappx-rtri-rll2} 得 $\rcrat(q) \approx_{\phi+\delta} y_\delta$，再由~\eqref{eq:RCappx-rtri-rll1} 得 $\rcrat(q) \approx_{\phi+\epsilon-\eta} z_\eta$。
    由 $\approx$ 的定义即推出 $\rcrat(q) \approx_{\phi+\epsilon} \rclim(z)$，正是所求。
  \end{enumerate}
  下面转到 $u$ 为 $\rclim(x)$（$x$ 为 Cauchy 近似）的情形。
  此时定义 $(\rclim(x) \approx_{(\blank)} {\blank}) : A$ 的外层归纳假设是已给出 ${(x_\delta \approx_{(\blank)} {\blank})}: A$，因此除圆整性外，还满足右侧三角不等式。
  \begin{enumerate}\setcounter{enumi}{4}
  \item 假设 $\rclim(x)\approx_\phi \rcrat(r)$ 且 $-\epsilon<|r-s|<\epsilon$，需证 $\rclim(x)\approx_{\phi+\epsilon} \rcrat(s)$。
    由 $\approx$ 定义，前者意即 $x_\delta \approx_{\phi-\delta} \rcrat(r)$，据上面的三角不等式得 $x_\delta \approx_{\epsilon+\phi-\delta} \rcrat(s)$，故 $\rclim(x)\approx_{\phi+\epsilon} \rcrat(s)$。
  \item 假设 $\phi,\epsilon,\delta:\Qp$ 满足 $\rclim(x)\approx_\phi \rcrat(r)$ 与 $\rcrat(r) \close{\epsilon-\delta} y_\delta$，另有两个未用到的归纳假设。
    %
    由定义可得某个 $\eta:\Qp$ 使 $x_\eta \approx_{\phi-\eta} \rcrat(r)$，于是归纳三角不等式给出 $x_\eta \approx_{\phi+\epsilon-\eta-\delta} y_\delta$。
    由 $\approx$ 的定义立即得到 $\rclim(x) \approx_{\phi+\epsilon} \rclim(y)$。
  \item 假设 $\phi,\epsilon,\delta:\Qp$ 满足 $\rclim(x)\approx_\phi \rclim(y)$ 与 $y_\delta \close{\epsilon-\delta} \rcrat(r)$，另有两个未用到的归纳假设。
    由定义可得 $\xi,\theta:\Qp$ 使 $x_\xi \approx_{\phi-\xi-\theta} y_\theta$。
    由于 $y$ 是 Cauchy 近似，有 $y_\theta \close{\theta+\delta} y_\delta$，故由归纳三角不等式得 $x_\xi \approx_{\phi+\delta-\xi} y_\delta$，再得 $x_\xi \close{\phi+\epsilon-\xi} \rcrat(r)$。
    由 $\approx$ 的定义得 $\rclim(x) \approx_{\phi+\epsilon}\rcrat(r)$，如愿。
  \item 最后，假设 $\phi,\epsilon,\delta,\eta:\Qp$ 满足 $\rclim(x)\approx_\phi \rclim(y)$ 与 $y_\delta \close{\epsilon-\delta-\eta} z_\eta$。
    如前可得 $\xi,\theta:\Qp$ 使 $x_\xi \approx_{\phi-\xi-\theta} y_\theta$，再应用两次三角不等式即可同样完成证明。
  \end{enumerate}

  由此完成两个内层递归，也就完成了关系族 $(\rcrat(q)\approx_{(\blank)}\blank)$ 与 $(\rclim(x)\approx_{(\blank)}\blank)$ 的定义。
  由于它们都属于 $A$，故均圆整，并且相对于 $\closesym$ 满足右侧三角不等式。
% , and satisfy~\eqref{eq:appxrec2}.
  剩下要验证的是与 $\bsim$ 有关的条件，也就是这些关系在\emph{左侧}相对于 $\closesym$ 构造子满足三角不等式。
  四种情形对应~\eqref{eq:RC-sim-ltri} 中 $u,v$ 取有理点或极限点的四种选择；又因为它们都是纯命题，可再对 $w$ 用 $\RC$-归纳，假设 $w$ 也为有理点或极限点。
  这又产生八种子情形，其证明与刚才本质相同；我们就不再重复折磨读者了。
\end{proof}

现在可证明：

\begin{thm}\label{thm:RC-sim-characterization}
  对任意 $u,v:\RC$ 与 $\epsilon:\Qp$，有 $(u\close\epsilon v) = (u\approx_\epsilon v)$。
\end{thm}
\begin{proof}
  两边都是纯命题，因此只需双向蕴含。
  对左推右，应用 $\closesym$-归纳到 $C(u,v,\epsilon)\defeq (u\approx_\epsilon v)$。
  因而只需考察 $\closesym$ 的四个构造子。
  每一情形中，$u,v$ 都专门化为有理点或极限点，于是 $\approx$ 的定义可化简，且归纳假设总能应用。

  对右推左，用 $\RC$-归纳假设 $u,v$ 为有理点或极限点，以便 $\approx$ 计算化简。
  此时 $\approx$ 的定义与归纳假设恰好给出相应 $\closesym$ 构造子所需的数据。
\end{proof}

\index{encode-decode method}%
稍微牵强地说，也可把 $\approx$ 看作 $\closesym$ 的一束“编码”纤维，上述证明的两个方向分别是 \encode 与 \decode。
由 $\approx$ 的定义以及 \cref{thm:RC-sim-characterization}，得到如下等价：
\begin{align*}
  (\rcrat(q) \close\epsilon \rcrat(r))  &=
  (-\epsilon < q - r < \epsilon)\\
  (\rcrat(q) \close\epsilon \rclim(y)) &=
  \exis{\delta : \Qp} \rcrat(q) \close{\epsilon - \delta} y_\delta\\
  (\rclim(x) \close\epsilon \rcrat(r)) &=
  \exis{\delta : \Qp} x_\delta \close{\epsilon - \delta} \rcrat(r)\\
  (\rclim(x) \close\epsilon \rclim(y)) &=
  \exis{\delta, \eta : \Qp} x_\delta \close{\epsilon - \delta - \eta} y_\eta.
\end{align*}
我们的证明还给出如下附加信息。

\begin{cor}
  \index{triangle!inequality for R@inequality for $\RC$}%
  \indexsee{inequality!triangle}{triangle inequality}%
  $\closesym$ 是圆整\index{rounded!relation}的，并满足三角不等式：
    \begin{gather}
      \eqvspaced{
        (u \close\epsilon v)
      }{
        \exis{\theta : \Qp} u \close{\epsilon - \theta} v
      }\\
      (u\close\epsilon v) \to (v\close\delta w) \to (u\close{\epsilon+\delta} w). \label{item:RC-sim-triangle}
    \end{gather}
\end{cor}
% \begin{proof}
%   The construction of $\approx$ showed simultaneously that it is rounded, and satisfies ``triangle inequalities'' such as
%   \[ (u\approx_\epsilon v) \to (v\close\delta w) \to (u\approx_{\epsilon+\delta} w). \]
%   Thus, both properties follow from \cref{thm:RC-sim-characterization}.
% \end{proof}

有了三角不等式，我们可以证明：Cauchy 近似的“极限”确实按极限应有的方式工作。

\begin{lem}\label{thm:RC-sim-lim}
  对任意 $u:\RC$、Cauchy 近似 $y$ 与 $\epsilon,\delta:\Qp$，若 $u\close\epsilon y_\delta$，则 $u\close{\epsilon+\delta} \rclim(y)$。
\end{lem}
\begin{proof}
  对 $u$ 用 $\RC$-归纳。
  若 $u$ 为 $\rcrat(q)$，这正是 $\closesym$ 的第二个构造子。
  现设 $u$ 为 $\rclim(x)$，并归纳假设每个 $x_\eta$ 都满足：对任意 $y,\epsilon,\delta$，若 $x_\eta\close\epsilon y_\delta$ 则 $x_\eta \close{\epsilon+\delta} \rclim(y)$。
  特别取 $y\defeq x$、$\delta\defeq\eta$，可得对任意 $\eta,\theta:\Qp$，有 $x_\eta \close{\eta+\theta} \rclim(x)$。

  现在任取 $y,\epsilon,\delta$ 并假设 $\rclim(x) \close\epsilon y_\delta$。
  由圆整性，存在 $\theta$ 使 $\rclim(x) \close{\epsilon-\theta} y_\delta$。
  再由上面的观察，对任意 $\eta$ 有 $x_\eta \close{\eta+\theta/2} \rclim(x)$，故由三角不等式得 $x_\eta \close{\epsilon+\eta-\theta/2} y_\delta$。
  因而由 $\closesym$ 的第四个构造子得到 $\rclim(x) \close{\epsilon+2\eta+\delta-\theta/2} \rclim(y)$。
  于是取 $\eta \defeq \theta/4$ 即可得到结论。
\end{proof}

\begin{lem}\label{thm:RC-sim-lim-term}
  对任意 Cauchy 近似 $y$ 与任意 $\delta,\eta:\Qp$，有 $y_\delta \close{\delta+\eta} \rclim(y)$。
\end{lem}
\begin{proof}
  在前一引理中取 $u\defeq y_\delta$、$\epsilon\defeq \eta$。
\end{proof}

\begin{rmk}
  我们也许会期待有 $y_\delta \close{\delta} \rclim(y)$，但在实例中这并不成立。
  例如取 $x$ 满足 $x_\epsilon \defeq \epsilon$。
  它的极限显然是 $0$，但我们并没有 $|\epsilon - 0 |<\epsilon$，只能得到 $\le$。
\end{rmk}

作为应用，\cref{thm:RC-sim-lim-term} 使我们能够说明：\cref{RC-extend-Q-Lipschitz} 中 Lipschitz 函数的延拓是唯一的。

\begin{lem}\label{RC-continuous-eq}
  \index{function!continuous}%
  设 $f,g:\RC\to\RC$ 连续，即
  \[ \fall{u:\RC}{\epsilon:\Qp}\exis{\delta:\Qp}\fall{v:\RC} (u\close\delta v) \to (f(u) \close\epsilon f(v)) \]
  $g$ 同理。
  若对所有 $q:\Q$ 都有 $f(\rcrat(q))=g(\rcrat(q))$，则 $f=g$。
\end{lem}
\begin{proof}
  我们对所有 $u$ 用 $\RC$-归纳证明 $f(u)=g(u)$。
  有理情形正是前提。
  因此设对所有 $\delta$ 有 $f(x_\delta)=g(x_\delta)$。
  我们将证明对所有 $\epsilon$，有 $f(\rclim(x))\close\epsilon g(\rclim(x))$，从而可用 $\RC$ 的路径构造子。

  由 $f,g$ 连续，存在 $\theta,\eta$ 使对所有 $v$，有
  \begin{align*}
    (\rclim(x)\close\theta v) &\to (f(\rclim(x)) \close{\epsilon/2} f(v))\\
    (\rclim(x)\close\eta v) &\to (g(\rclim(x)) \close{\epsilon/2} g(v)).
  \end{align*}
  取 $\delta < \min(\theta,\eta)$，由 \cref{thm:RC-sim-lim-term} 得到同时有 $\rclim(x)\close\theta y_\delta$ 与 $\rclim(x)\close\eta y_\delta$。
  因而
  \[ f(\rclim(x)) \close{\epsilon/2} f(y_\delta) = g(y_\delta) \close{\epsilon/2} g(\rclim(x))\]
  再由三角不等式得 $f(\rclim(x))\close\epsilon g(\rclim(x))$。
\end{proof}

\subsection{Cauchy 实数的代数结构}
\label{sec:algebr-struct-cauchy}

我们先定义加法结构 $(\RC, 0, {+}, {-})$。显然，加法单位元
$0$ 就是 $\rcrat(0)$，而加法逆元 ${-} : \RC \to \RC$ 由
加法逆元 ${-} : \Q \to \Q$ 的延拓得到，使用 \cref{RC-extend-Q-Lipschitz}
且 Lipschitz 常数为~$1$。加法本身则需要多做一些工作。

\begin{lem} \label{RC-binary-nonexpanding-extension}
  设 $f : \Q \times \Q \to \Q$ 满足对所有 $q, r, s : \Q$，
  %
  \begin{equation*}
    |f(q, s) - f(r, s)| \leq |q - r|
    \qquad\text{and}\qquad
    |f(q, r) - f(q, s)| \leq |r - s|.
  \end{equation*}
  %
  则存在函数 $\bar{f} : \RC \times \RC \to \RC$，满足
  对所有 $q, r : \Q$ 有 $\bar{f}(\rcrat(q), \rcrat(r)) = f(q,r)$。并且
  对所有 $u, v, w : \RC$ 与 $q : \Qp$，有
  %
  \begin{equation*}
    u \close\epsilon v \Rightarrow \bar{f}(u,w) \close\epsilon \bar{f}(v,w)
    \quad\text{and}\quad
    v \close\epsilon w \Rightarrow \bar{f}(u,v) \close\epsilon \bar{f}(u,w).
  \end{equation*}
\end{lem}

\begin{proof}
  我们用 $(\RC, {\closesym})$-递归来构造 $\bar{f}$ 的柯里化形式，视为映射
  $\RC \to A$，其中 $A$ 是取值于实数的非扩张\index{function!non-expanding}\index{non-expanding function}函数空间：
  %
  \begin{equation*}
    A \defeq
    \setof{ h : \RC \to \RC |
      \fall{\epsilon : \Qp} \fall{u, v : \RC}
      u \close\epsilon v \Rightarrow h(u) \close\epsilon h(v)
    }.
  \end{equation*}
  %
  我们还需要在 $A$ 上给出合适的 $\bsim_\epsilon$，定义为
  %
  \begin{equation*}
    (h \bsim_\epsilon k) \defeq \fall{u : \RC} h(u) \close\epsilon k(u).
  \end{equation*}
  %
  显然，若 $\fall{\epsilon : \Qp} h \bsim_\epsilon k$，则对所有 $u :
  \RC$ 有 $h(u) = k(u)$，因此 $\bsim$ 是分离的。

  基础情形中，对 $q : \Q$，定义 $\bar{f}(\rcrat(q)) : A$ 为
  Lipschitz 映射 $\lam{r} f(q,r)$ 从 $\Q \to \Q$ 到 $\RC \to \RC$ 的延拓，
  即 \cref{RC-extend-Q-Lipschitz} 在 Lipschitz 常数~$1$ 下构造的映射。接着，对
  Cauchy 近似 $x$，定义 $\bar{f}(\rclim(x)) : \RC \to \RC$ 为
  %
  \begin{equation*}
    \bar{f}(\rclim(x))(v) \defeq \rclim (\lam{\epsilon} \bar{f}(x_\epsilon)(v)).
  \end{equation*}
  %
  为使定义合法，$\lam{\epsilon} \bar{f}(x_\epsilon)(v)$ 应是
  Cauchy 近似。故取任意 $\delta, \epsilon : \Q$，由假设
  $\bar{f}(x_\delta) \bsim_{\delta + \epsilon} \bar{f}(x_\epsilon)$，从而
  $\bar{f}(x_\delta)(v) \close{\delta + \epsilon} \bar{f}(x_\epsilon)(v)$。此外，
  $\bar{f}(\rclim(x))$ 之所以非扩张，是因为按归纳假设 $\bar{f}(x_\epsilon)$ 非扩张。
  确实，若 $u \close\epsilon v$，则对所有 $\epsilon : \Q$ 有
  %
  \begin{equation*}
    \bar{f}(x_{\epsilon/3})(u) \close{\epsilon/3} \bar{f}(x_{\epsilon/3})(v),
  \end{equation*}
  %
  因而由 $\closesym$ 的第四个构造子得 $\bar{f}(\rclim(x))(u) \close\epsilon \bar{f}(\rclim(x))(v)$。

  还需再检查四个条件，我们只示范其中一个。设
  $\epsilon : \Qp$，并且对某个 $\delta : \Qp$ 有 $\rcrat(q) \close{\epsilon - \delta}
  y_\delta$ 与 $\bar{f}(\rcrat(q)) \bsim_{\epsilon - \delta} \bar{f}(y_\delta)$。要证
  $\bar{f}(\rcrat(q)) \bsim_\epsilon \bar{f}(\rclim(y))$，取任意 $v : \RC$ 并观察
  %
  \begin{equation*}
    \bar{f}(\rcrat(q))(v) \close{\epsilon - \delta} \bar{f}(y_\delta)(v).
  \end{equation*}
  %
  因此，由 $\closesym$ 的第二个构造子得
  \narrowequation{\bar{f}(\rcrat(q))(v) \close\epsilon \bar{f}(\rclim(y))(v)}
  如所需。
\end{proof}

我们可将 \cref{RC-binary-nonexpanding-extension} 应用于任一在每个变量上分别非扩张的二元有理函数。加法就是这样的函数，因此得到 ${+} : \RC \times \RC \to \RC$。
\indexdef{addition!of Cauchy reals}%
此外，只要要求延拓在每个变量上都非扩张，该延拓就是唯一的；并且与一元情形一样，有理数上的恒等式可延拓为实数上的恒等式。由于非扩张映射的复合仍非扩张，我们可推出加法满足通常性质，例如交换律与结合律。
\index{associativity!of addition!of Cauchy reals}%
因此，$(\RC, 0, {+}, {-})$ 是交换
群。

我们还可将 \cref{RC-binary-nonexpanding-extension} 应用于函数 $\min : \Q \times
\Q \to \Q$ 与 $\max : \Q \times \Q \to \Q$，从而使 $\RC$ 成为格。偏
序 $\leq$ 在 $\RC$ 上由 $\max$ 定义为
%
\symlabel{leq-RC}
\index{order!non-strict}%
\index{non-strict order}%
\begin{equation*}
  (u \leq v) \defeq (\max(u, v) = v).
\end{equation*}
%
关系 $\leq$ 是偏序，因为在 \Q 上它就是偏序，而偏序公理可用 $\min$ 与 $\max$ 的等式表达，因此可传递到
$\RC$。

\index{absolute value}%
另一个可用同样方法延拓到 $\RC$ 的函数是绝对值 $|{\blank}|$。
同样，它具有预期性质，因为这些性质可从 \Q 传递到 $\RC$。

\symlabel{lt-RC}
由 $\leq$ 可得严格序 $<$：
\index{strict!order}%
\index{order!strict}%
%
\begin{equation*}
  (u < v) \defeq \exis{q, r : \Q} (u \leq \rcrat(q)) \land (q < r) \land (\rcrat(r) \leq v).
\end{equation*}
%
也就是说，当且仅当仅存在一对有理数 $q < r$ 使 $x \leq
\rcrat(q)$ 且 $\rcrat(r) \leq v$ 时，$u < v$ 成立。不难检查，$<$ 是反自反且传递的，并满足有序域所期望的其他性质。
阿基米德原理可直接由~$<$ 的定义推出。

\index{ordered field!archimedean}%
\begin{thm}[$\RC$ 的阿基米德原理] \label{RC-archimedean}
  %
  对每个 $u, v : \RC$，若 $u < v$，则仅存在 $q : \Q$ 使 $u < q < v$。
\end{thm}

\begin{proof}
  由 $u < v$，仅可得 $r, s : \Q$ 使 $u \leq r < s \leq v$，取 $q
  \defeq (r + s) / 2$ 即可。
\end{proof}

现在我们在 $\RC$ 上已有足够结构，可以用标准概念表达 $u \close\epsilon v$。

\begin{lem}\label{thm:RC-le-grow}
  若 $q:\Q$ 与 $u:\RC$ 满足 $u\le \rcrat(q)$，则对任意 $v:\RC$ 与 $\epsilon:\Qp$，若 $u\close\epsilon v$，则 $v\le \rcrat(q+\epsilon)$。
\end{lem}
\begin{proof}
  注意函数 $\max(\rcrat(q),\blank):\RC\to\RC$ 的 Lipschitz 常数为 $1$。
  先考虑 $u=\rcrat(r)$ 为有理点的情形。
  这里对 $v$ 做归纳。
  若 $v$ 是有理点，结论显然。
  若 $v$ 是 $\rclim(y)$，归纳假设为：对任意 $\epsilon,\delta$，若 $\rcrat(r)\close\epsilon y_\delta$，则 $y_\delta \le \rcrat(q+\epsilon)$，即 $\max(\rcrat(q+\epsilon),y_\delta)=\rcrat(q+\epsilon)$。

  现在给定 $\epsilon$ 并假设 $\rcrat(r)\close\epsilon \rclim(y)$，则存在 $\theta$ 使 $\rcrat(r)\close{\epsilon-\theta} \rclim(y)$，从而在 $\delta<\theta$ 时有 $\rcrat(r)\close\epsilon y_\delta$。
  因此归纳假设给出对这类 $\delta$ 有 $\max(\rcrat(q+\epsilon),y_\delta)=\rcrat(q+\epsilon)$。
  但按定义，
  \[\max(\rcrat(q+\epsilon),\rclim(y)) \jdeq \rclim(\lam{\delta} \max(\rcrat(q+\epsilon),y_\delta)).\]
  由于一个最终常值的 Cauchy 近似的极限就是该常值，我们有
  \[\max(\rcrat(q+\epsilon),\rclim(y)) = \rcrat(q+\epsilon),\] 故 $\rclim(y)\le \rcrat(q+\epsilon)$。

  再考虑一般 $u:\RC$。
  因为 $u\le \rcrat(q)$ 意味着 $\max(\rcrat(q),u)=\rcrat(q)$，由假设 $u\close\epsilon v$ 与 $\max(\rcrat(q),-)$ 的 Lipschitz 性可得 $\max(\rcrat(q),v) \close\epsilon \rcrat(q)$。
  因此又由 $\rcrat(q)\le \rcrat(q)$，第一种情形推出 $\max(\rcrat(q),v) \le \rcrat(q+\epsilon)$，再由 $\le$ 的传递性得 $v\le \rcrat(q+\epsilon)$。
\end{proof}

\begin{lem}\label{thm:RC-lt-open}
  设 $q:\Q$ 与 $u:\RC$ 满足 $u<\rcrat(q)$。则：
  \begin{enumerate}
  \item 对任意 $v:\RC$ 与 $\epsilon:\Qp$，若 $u\close\epsilon v$，则 $v< \rcrat(q+\epsilon)$。\label{item:RCltopen1}
  \item 存在 $\epsilon:\Qp$，使对任意 $v:\RC$，若 $u\close\epsilon v$ 则有 $v<\rcrat(q)$。\label{item:RCltopen2}
  \end{enumerate}
\end{lem}
\begin{proof}
  按定义，$u<\rcrat(q)$ 意味着存在 $r:\Q$ 使 $r<q$ 且 $u\le \rcrat(r)$。
  由 \cref{thm:RC-le-grow}，对任意 $\epsilon$，若 $u\close\epsilon v$ 则 $v\le \rcrat(r+\epsilon)$。
  结论~\ref{item:RCltopen1} 立即成立，因为 $r+\epsilon<q+\epsilon$；而对~\ref{item:RCltopen2}，取任意满足 $\epsilon <q-r$ 的 $\epsilon$ 即可。
\end{proof}

现在我们可以说明，辅助关系 $\closesym$ 正是我们所设想的关系。

\begin{thm} \label{RC-sim-eqv-le}
  \index{distance}%
  $\eqv{(u \close\epsilon v)}{(|u - v| < \rcrat(\epsilon))}$
  对所有 $u, v : \RC$ 与 $\epsilon : \Qp$ 都成立。
\end{thm}
\begin{proof}
  减法与绝对值的 Lipschitz 性蕴含：若 $u\close\epsilon v$，则 $|u-v| \close\epsilon |u-u| = 0$。
  因此，对左推右方向，只需证明若 $u\close\epsilon 0$，则 $|u|<\rcrat(\epsilon)$。
  我们对 $u$ 做 $\RC$-归纳。

  若 $u$ 是有理点，结论立刻成立，因为绝对值与次序在 $\Q$ 上延拓自标准定义。
  若 $u$ 是 $\rclim(x)$，由圆整性可得 $\theta:\Qp$ 使 $\rclim(x)\close{\epsilon-\theta} 0$。
  因此由三角不等式有 $x_{\theta/3} \close{\epsilon-2\theta/3} 0$，归纳假设给出 $|x_{\theta/3}|<\rcrat(\epsilon-2\theta/3)$。
  但 $x_{\theta/3} \close{2\theta/3} \rclim(x)$，故由 Lipschitz 性得 $|x_{\theta/3}| \close{2\theta/3} |\rclim(x)|$，于是 \cref{thm:RC-lt-open}\ref{item:RCltopen1} 推出 $|\rclim(x)|<\rcrat(\epsilon)$。

  对另一方向，我们对 $u,v$ 做 $\RC$-归纳。
  若二者皆为有理点，这就是 $\closesym$ 的第一个构造子。

  若 $u$ 是 $\rcrat(q)$ 且 $v$ 是 $\rclim(y)$，归纳假设是：对任意 $\epsilon,\delta$，若 $|\rcrat(q)-y_\delta|<\rcrat(\epsilon)$ 则 $\rcrat(q) \close{\epsilon} y_\delta$。
  固定一个 $\epsilon$，并假设 $|\rcrat(q) - \rclim(y)|<\rcrat(\epsilon)$。
  由于 \Q 在 $\RC$ 中序稠密，存在 $\theta<\epsilon$ 使 $|\rcrat(q) - \rclim(y)|<\rcrat(\theta)$。
  现在任取 $\delta,\eta$，有 $\rclim(y)\close{2\delta} y_\delta$，故由 Lipschitz 性得
  \[ |\rcrat(q) - \rclim(y)| \close{\delta+\eta} |\rcrat(q) - y_\delta|. \]
  因此由 \cref{thm:RC-lt-open}\ref{item:RCltopen1} 得 $|\rcrat(q) - y_\delta| < \rcrat(\theta+2\delta)$。
  再由归纳假设得 $\rcrat(q) \close{\theta+2\delta} y_\delta$，从而由三角不等式得 $\rcrat(q)\close{\theta+4\delta} \rclim(y)$。
  因而只需取 $\delta \defeq (\epsilon-\theta)/4$。

  剩余两种情形完全类似。
\end{proof}

\indexdef{multiplication!of Cauchy reals}%
接下来我们希望赋予 $\RC$ 乘法结构。对每个 $q : \Q$，映射
$r \mapsto q \cdot r$ 的 Lipschitz 常数是\footnote{我们把 Lipschitz
  常数定义为\emph{正}有理数。} $|q| + 1$，因此可延拓为实数上的按 $q$ 乘法。于是 $\RC$ 成为 \Q 上的向量空间\index{vector!space}。
一般地，实数乘法可定义为
%
\begin{equation}
  u \cdot v \defeq
  {\textstyle \frac{1}{2}} \cdot ((u + v)^2 - u^2 - v^2),\label{mult-from-square}
\end{equation}
%
因此只需把平方\index{squaring function} $u \mapsto u^2$ 定义为映射 $\RC \to \RC$。平方映射并非
Lipschitz，但在每个有界域上都是 Lipschitz，这使我们可以把它拼接起来。定义开区间与闭区间
%
\indexdef{interval!open and closed}%
\indexdef{open!interval}%
\indexdef{closed!interval}%
\begin{equation*}
  [u,v] \defeq \setof{ x : \RC | u \leq x \leq v }
  \qquad\text{and}\qquad
  (u,v) \defeq \setof{ x : \RC | u < x < v }.
\end{equation*}
%
严格说来，$[u,v]$ 或 $(u,v)$ 的元素是一个 Cauchy 实数加上一份证明；但后者处于纯命题中，因此没有额外信息。
因此，与子集类型的常见做法一样，只要 $x:\RC$ 且满足 $u\leq x \leq v$，我们通常直接写 $x:[u,v]$，其余同理。

\begin{thm} \label{RC-squaring}
  %
  存在唯一函数 ${(\blank)}^2 : \RC \to \RC$，延拓了有理数平方 $q \mapsto
  q^2$，并满足
  %
  \begin{equation*}
    \fall{n : \N}
    \fall{u, v : [-n, n]}
    |u^2 - v^2| \leq 2 \cdot n \cdot |u - v|.
  \end{equation*}
\end{thm}

\begin{proof}
  先注意：对每个 $u : \RC$，仅存在 $n : \N$ 使 $-n
  \leq u \leq n$，见 \cref{ex:traditional-archimedean}，所以映射
  %
  \begin{equation*}
    e : \Parens{\sm{n : \N} [-n, n]} \to \RC
    \qquad\text{defined by}\qquad
    e(n, x) \defeq x
  \end{equation*}
  %
  是满射。接着，对每个 $n : \N$，平方映射
  %
  \begin{equation*}
    s_n : \setof{ q : \Q | -n \leq q \leq n } \to \Q
    \qquad\text{defined by}\qquad
    s_n(q) \defeq q^2
  \end{equation*}
  %
  以 $2 n$ 为 Lipschitz 常数，因此可由 \cref{RC-extend-Q-Lipschitz}
  延拓为映射 $\bar{s}_n : [-n, n] \to \RC$，Lipschitz 常数同为 $2 n$，细节见
  \cref{RC-Lipschitz-on-interval}。这些映射 $\bar{s}_n$ 彼此相容：若某些 $m, n : \N$ 满足
  $m < n$，则 $s_n$ 在 $[-m, m]$ 上的限制必与 $s_m$
  一致，因为二者都 Lipschitz，因此按
  \cref{RC-continuous-eq} 的意义连续。故由 \cref{lem:images_are_coequalizers}，映射
  %
  \begin{equation*}
    \Parens{\sm{n : \N} [-n, n]} \to \RC,
    \qquad\text{given by}\qquad
    (n, x) \mapsto s_n(x)
  \end{equation*}
  %
  唯一地经由 $\RC$ 因子化，从而得到所需函数。
\end{proof}

到这里我们已经有了实数的环结构与阿基米德次序。要把 $\RC$ 建立为阿基米德有序域，还需要逆元。

\begin{thm}
  \index{apartness}%
  一个 Cauchy 实数可逆，当且仅当它与零分离。
\end{thm}

\begin{proof}
  先设 $u : \RC$ 有逆元 $v : \RC$。由阿基米德原理，存在 $q :
  \Q$ 使 $|v| < q$。于是 $1 = |u v| < |u| \cdot v < |u| \cdot q$，从而 $|u| >
  1/q$，即 $u \apart 0$。

  反向地，我们构造逆映射
  %
  \begin{equation*}
    ({\blank})^{-1} : \setof{ u : \RC | u \apart 0 } \to \RC
  \end{equation*}
  %
  方法与 \cref{RC-squaring} 中平方构造类似，也是通过拼接函数。这里只概述主要步骤。对每个 $q : \Q$，定义
  %
  \begin{equation*}
    [q, \infty) \defeq \setof{u : \RC | q \leq u}
    \qquad\text{and}\qquad
    (-\infty, q] \defeq \setof{u : \RC | u \leq -q}.
  \end{equation*}
  %
  当 $q$ 遍历 $\Qp$ 时，类型 $(-\infty, q]$ 与 $[q, \infty)$ 共同覆盖
  $\setof{u : \RC | u \apart 0}$。在每个 $[q, \infty)$ 与 $(-\infty, q]$ 上，
  逆函数都可由 \cref{RC-extend-Q-Lipschitz}
  以 Lipschitz 常数 $1/q^2$ 得到。最后，\cref{lem:images_are_coequalizers}
  保证该逆函数唯一地经由 $\setof{ u : \RC | u
    \apart 0 }$ 因子化。
\end{proof}

我们用一个定理总结 $\RC$ 的代数结构。

\begin{thm} \label{RC-archimedean-ordered-field}
  Cauchy 实数构成阿基米德有序域。
\end{thm}

\subsection{Cauchy 实数是 Cauchy 完备的}
\label{sec:cauchy-reals-cauchy-complete}

我们是通过在 Cauchy 近似的极限下封闭 \Q 来构造 $\RC$ 的，因此
$\RC$ 理应是 Cauchy 完备的。借助 \cref{RC-sim-eqv-le}，在构造
$\RC$ 时定义的 Cauchy 近似 $x : \Qp \to \RC$，与 \cref{defn:cauchy-approximation}
（适配到 $\RC$）意义下的 Cauchy 近似并无区别。

因此，给定 Cauchy 近似 $x : \Qp \to \RC$，自然会期望
$\rclim(x)$ 就是它的极限，而极限概念按
\cref{defn:cauchy-approximation} 定义。由 \cref{RC-sim-eqv-le} 与
\cref{thm:RC-sim-lim-term}，确实如此。我们已经证明：

\begin{thm}
  $\RC$ 中每个 Cauchy 近似都有极限。
\end{thm}

每个 Cauchy 近似都有极限的阿基米德有序域称为
\define{Cauchy 完备}。
\indexdef{Cauchy!completeness}%
\indexdef{complete!ordered field, Cauchy}%
\index{ordered field}%
Cauchy 实数是这样的域中最小的一个。

\begin{thm} \label{RC-initial-Cauchy-complete}
  Cauchy 实数可嵌入到每一个 Cauchy 完备的阿基米德有序域中。
\end{thm}

\begin{proof}
  \index{limit!of a Cauchy approximation}%
  设 $F$ 是一个 Cauchy 完备的阿基米德有序域。由于极限唯一，
  存在算子 $\lim$，把 $F$ 中的 Cauchy 近似送到其极限。我们用 $(\RC, {\closesym})$-递归定义嵌入 $e : \RC \to F$：
  %
  \begin{equation*}
    e(\rcrat(q)) \defeq q
    \qquad\text{and}\qquad
    e(\rclim(x)) \defeq \lim (e \circ x).
  \end{equation*}
  %
  在 $F$ 上可取合适的 $\bsim$ 为
  %
  \begin{equation*}
    (a \bsim_\epsilon b) \defeq |a - b| < \epsilon.
  \end{equation*}
  %
  由于 $F$ 是阿基米德的，该关系是分离的。其余
  $(\RC, {\closesym})$-递归所需分支都容易验证。还需检查 $e$
  是保持有理数不动的有序域嵌入。
\end{proof}

\index{real numbers!Cauchy|)}%
\section{Cauchy 实数与 Dedekind 实数的比较}
\label{sec:comp-cauchy-dedek}

\index{real numbers!Dedekind|(}%
\index{real numbers!Cauchy|(}%
\index{depression|(}

我们也来谈谈 Cauchy 实数与 Dedekind 实数之间的关系。由
\cref{RC-archimedean-ordered-field}，$\RC$ 是一个阿基米德有序域。它对 $\Omega$ 也是可容许的\index{ordered field!admissible}，这一点很容易检验。
（若 $\Omega$ 是初始
$\sigma$-frame
\index{initial!sigma-frame@$\sigma$-frame}%
\index{sigma-frame@$\sigma$-frame!initial}%
则需一个简单归纳；在其他情形下则是直接的。）
因此，由 \cref{RD-final-field} 存在一个固定有理数的有序域嵌入
%
\begin{equation*}
  \RC \to \RD
\end{equation*}
%
（也可由 \cref{RC-initial-Cauchy-complete,RD-cauchy-complete} 得到这一点。）
一般而言，在不加额外假设时，我们并不期望 $\RC$ 与 $\RD$ 重合。

\begin{lem} \label{lem:untruncated-linearity-reals-coincide}
  %
  若对每个 $x : \RD$ 仅仅存在
  %
  \begin{equation}
    \label{eq:untruncated-linearity}
    c : \prd{q, r : \Q} (q < r) \to (q < x) + (x < r)
  \end{equation}
  %
  则 Cauchy 实数与 Dedekind 实数重合。
\end{lem}

\begin{proof}
  注意~\eqref{eq:untruncated-linearity} 中的类型是~\eqref{eq:RD-linear-order} 的一个未截断变体，后者断言 $<$ 是弱线序。
  我们已知 $\RC$ 可嵌入到 $\RD$，因此只需证明每个 Dedekind 实数仅仅是某个有理数 Cauchy 序列\index{Cauchy!sequence}的极限。

  取任意 $x : \RD$。由假设仅仅存在引理陈述中的 $c$，并且由切割的可居留性\index{cut!Dedekind}仅仅存在 $a, b : \Q$ 使得 $a < x < b$。
  我们递归构造序列\index{sequence}
  $f : \N \to \setof{ \pairr{q, r} \in \Q \times \Q | q < r }$：
  %
  \begin{enumerate}
  \item 令 $f(0) \defeq \pairr{a, b}$。
  \item 若 $f(n)$ 已定义为满足 $q_n < r_n$ 的 $\pairr{q_n, r_n}$，定义
    $s \defeq (2 q_n + r_n)/3$ 与 $t \defeq (q_n + 2 r_n)/3$。
    然后 $c(s,t)$ 在 $s < x$ 与 $x < t$ 之间作出判定。
    若判定为 $s < x$，则令 $f(n+1) \defeq \pairr{s, r_n}$；否则令 $f(n+1) \defeq \pairr{q_n, t}$。
  \end{enumerate}
  %
  记序列 $f$ 的第 $n$ 项为 $\pairr{q_n, r_n}$。容易看出对所有 $n : \N$，都有
  $q_n < x < r_n$ 且 $|q_n - r_n| \leq (2/3)^n \cdot |q_0 - r_0|$。
  因此，$q_0, q_1, \ldots$ 与 $r_0, r_1, \ldots$ 都是收敛到 Dedekind 切割 $x$ 的 Cauchy 序列。
  这表明：对每个 $x : \RD$，仅仅存在一个收敛到 $x$ 的 Cauchy 序列。
\end{proof}

该引理推出：可数选择或排中律任一成立，都足以保证 $\RC$ 与 $\RD$ 重合。

\begin{cor} \label{when-reals-coincide}
  \index{axiom!of choice!countable}%
  \index{excluded middle}%
  若排中律或可数选择成立，则 $\RC$ 与 $\RD$ 等价。
\end{cor}

\begin{proof}
  若排中律成立，则可证明 $(x < y) \to (x < z) + (z < y)$：要么 $x < z$，要么 $\lnot (x < z)$。
  前者已得结论；后者由 $z \leq x < y$ 得 $z < y$。
  因此可得到~\eqref{eq:untruncated-linearity}，从而可应用 \cref{lem:untruncated-linearity-reals-coincide}。

  现在设可数选择成立。集合
  $S = \setof{ \pairr{q, r} \in \Q \times \Q | q <
    r }$
  与 $\N$ 等价，因此可把可数选择应用于“$x$ 是可定位的”这一命题：
  %
  \begin{equation*}
    \fall{\pairr{q, r} : S} (q < x) \lor (x < r).
  \end{equation*}
  %
  注意 $(q < x) \lor (x < r)$ 可表述为存在语句
  $\exis{b :
    \bool} (b = \bfalse \to q < x) \land (b = \btrue \to x < r)$。
  于是选择函数（柯里化形式）恰好就是~\eqref{eq:untruncated-linearity}，因此同样可应用
  \cref{lem:untruncated-linearity-reals-coincide}。
\end{proof}

\index{real numbers!Dedekind|)}%
\index{real numbers!Cauchy|)}%
\index{real numbers!agree}%

\index{depression|)}
\section{区间的紧致性}
\label{sec:compactness-interval}

\index{mathematics!classical|(}%
\index{mathematics!constructive|(}%

我们已经指出，我们对实数的构造与经典逻辑完全相容。
因此，只要假设排中律~\eqref{eq:lem} 与选择公理~\eqref{eq:ac}，就可以发展经典分析，\index{classical!analysis}\index{analysis!classical}其本质上就是把任一本标准分析教材照搬过来。

\index{analysis!constructive}%
\index{constructive!analysis}%
尽管如此，任何关心计算的人（例如数值分析师）都应当有兴趣在具有计算意义的环境中发展分析学。
分析学在构造性语境中确实可行，这一点由~\cite{Bishop1967} 证明。
为了展示经典分析与构造性分析的差异与相似，我们只简要讨论一个主题：闭区间
$[0,1]$ 的紧致性及其周边的几个定理。

紧致性也不例外，仍体现了构造性数学中的常见现象：在经典上等价的概念会发生分叉。
最常用的三种紧致性概念是：
%
\indexdef{compactness}%
\begin{enumerate}
\item \define{度量紧致：}“Cauchy 完备且全有界”，
  \indexdef{metrically compact}%
  \indexdef{compactness!metric}%
\item \define{Bolzano--Weierstra\ss{} 紧致：}“每个序列都有收敛子序列”，
  \index{compactness!Bolzano--Weierstrass@Bolzano--Weierstra\ss{}}%
  \indexsee{Bolzano--Weierstrass@Bolzano--Weierstra\ss{}}{compactness}%
  \index{sequence}%
\item \define{Heine--Borel 紧致：}“每个开覆盖都有有限子覆盖”。
  \index{compactness!Heine--Borel}%
  \indexsee{Heine--Borel}{compactness}%
\end{enumerate}
%
在经典数学中它们彼此等价。
我们来看看它们在同伦类型论中的表现。我们可以使用 Dedekind 实数或 Cauchy 实数，因此统一把实数记为~$\R$。
先回顾几个基本定义。

\indexsee{space!metric}{metric space}
\index{metric space|(}%

\begin{defn} \label{defn:metric-space}
  \define{度量空间}
  \indexdef{metric space}%
  $(M, d)$ 是一个集合 $M$，配备映射 $d : M \times M \to \R$，
  并且对所有 $x, y, z : M$ 满足
  %
  \begin{align*}
    d(x,y) &\geq 0, &
    d(x,y) &= d(y,x), \\
    d(x,y) &= 0 \Leftrightarrow x = y, &
    d(x,z) &\leq d(x,y) + d(y,z).
  \end{align*}
  %
\end{defn}

\begin{defn} \label{defn:complete-metric-space}
  $M$ 中的一个\define{Cauchy 近似}
  \index{Cauchy!approximation}%
  是满足下式的序列 $x : \Qp \to M$：
  %
  \begin{equation*}
    \fall{\delta, \epsilon} d(x_\delta, x_\epsilon) < \delta + \epsilon.
  \end{equation*}
  %
  \index{limit!of a Cauchy approximation}%
  Cauchy 近似 $x : \Qp \to M$ 的\define{极限}是点 $\ell : M$，满足
  %
  \begin{equation*}
    \fall{\epsilon, \theta : \Qp} d(x_\epsilon, \ell) < \epsilon + \theta.
  \end{equation*}
  %
  \indexdef{metric space!complete}%
  \indexdef{complete!metric space}%
  \define{完备度量空间}指每个 Cauchy 近似都存在极限的度量空间。
\end{defn}

\begin{defn} \label{defn:total-bounded-metric-space}
  对正有理数 $\epsilon$，度量空间 $(M,
  d)$ 中的一个\define{$\epsilon$-网}
  \indexdef{epsilon-net@$\epsilon$-net}%
  是下述类型的一个元素：
  %
  \begin{equation*}
    \sm{n : \N}{x_1, \ldots, x_n : M}
    \fall{y : M} \exis{k \leq n} d(x_k, y) < \epsilon.
  \end{equation*}
  %
  换言之，它是一个有限点列 $x_1, \ldots, x_n$，使得 $M$ 中每个点仅仅地位于某个~$x_k$ 的 $\epsilon$ 邻域内。

  度量空间 $(M, d)$ 称为\define{全有界}
  \indexdef{totally bounded metric space}%
  \indexdef{metric space!totally bounded}%
  ，若它对任意精度都存在 $\epsilon$-网：
  %
  \begin{equation*}
    \prd{\epsilon : \Qp}
    \sm{n : \N}{x_1, \ldots, x_n : M}
    \fall{y : M} \exis{k \leq n} d(x_k, y) < \epsilon.
  \end{equation*}
\end{defn}

\begin{rmk}
  在全有界性的定义里，我们用了不严格的记号 $\sm{n : \N}{x_1, \ldots, x_n : M}$。形式上应写为 $\sm{x : \lst{M}}$，
  其中 $\lst{M}$ 是 \cref{sec:bool-nat} 中的有限列表归纳类型\index{type!of lists}。
  但那会使后续表述稍显繁琐。
\end{rmk}

注意，在全有界的定义里我们要求的是 $\epsilon$-网的纯存在，而不是仅仅存在。
如此便得到一个函数，它给每个 $\epsilon : \Qp$ 指派一个具体的 $\epsilon$-网。
这个函数可称为“全有界模”。
一般而言，把经典度量概念迁移到同伦类型论时，应谨慎使用命题截断，通常是为了避免要求从 $\R$ 到 $\Q$ 或 $\N$ 的非常值映射。
例如，下面是一致连续的“正确”定义。

\begin{defn} \label{defn:uniformly-continuous}
  度量空间上的映射 $f : M \to \R$ 称为\define{一致连续}
  \indexdef{function!uniformly continuous}%
  \indexdef{uniformly continuous function}%
  ，若
  %
  \begin{equation*}
    \prd{\epsilon : \Qp}
    \sm{\delta : \Qp}
    \fall{x, y : M}
    d(x,y) < \delta \Rightarrow |f(x) - f(y)| < \epsilon.
  \end{equation*}
  %
  特别地，一致连续映射带有一致连续模\indexdef{modulus!of uniform continuity}，即给每个 $\epsilon$ 指派一个相应 $\delta$ 的函数。
\end{defn}

下面说明 $[0,1]$ 在第一种意义下是紧的。

\begin{thm} \label{analysis-interval-ctb}
  \index{compactness!metric}%
  \index{interval!open and closed}%
  闭区间 $[0,1]$ 是完备且全有界的。
\end{thm}

\begin{proof}
  给定 $\epsilon : \Qp$，存在 $k : \N$ 使得 $2/k < \epsilon$，于是可取 $\epsilon$-网 $x_i = i/k$（$i = 0, \ldots, k$）。
  这是一个 $\epsilon$-网，因为对每个 $y : [0,1]$，仅仅存在 $i$ 使得 $0 \leq i \leq k$ 且 $(i -
  1)/k < y < (i+1)/k$，从而 $|y - x_i| < 2/k < \epsilon$。

  对于 $[0,1]$ 的完备性，取一个 Cauchy 近似 $x : \Qp \to
  [0,1]$，令 $\ell$ 为它在 $\R$ 中的极限。
  由于 $\max$ 与 $\min$ 都是 Lipschitz 映射，故由
  $r(x) \defeq \max(0, \min(1, x))$ 定义的回缩 $r : \R \to [0,1]$ 与 Cauchy 近似的极限可交换，因此
  %
  \begin{equation*}
    r(\ell) =
    r (\lim x) =
    \lim (r \circ x) =
    \lim x =
    \ell,
  \end{equation*}
  %
  即 $0 \leq \ell \leq 1$，如所需。
\end{proof}

由此我们在同伦类型论中至少得到一种良好的紧致性概念。
但它局限于度量空间，因为全有界是度量概念。
稍后我们会讨论另外两种概念；不过先证明：在全有界空间上，一致连续映射有\define{上确界}，
\indexsee{least upper bound}{supremum}%
即一个上界，且小于等于所有其他上界。

\begin{thm} \label{ctb-uniformly-continuous-sup}
  %
  \indexdef{supremum!of uniformly continuous function}%
  在全有界度量空间 $(M, d)$ 上，一致连续映射 $f : M \to \R$
  存在上确界 $m : \R$。并且对每个 $\epsilon : \Qp$，存在 $u : M$ 使得
  $|m - f(u)| < \epsilon$。
\end{thm}

\begin{proof}
  令 $h : \Qp \to \Qp$ 为 $f$ 的一致连续模。
  下面定义一个近似 $x : \Qp \to \R$：对任意 $\epsilon : \Q$，由 $M$ 的全有界性可得一个 $h(\epsilon)$-网 $y_0, \ldots, y_n$。
  定义
  %
  \begin{equation*}
    x_\epsilon \defeq \max (f(y_0), \ldots, f(y_n)).
  \end{equation*}
  %
  我们断言 $x$ 是 Cauchy 近似。取任意 $\epsilon, \eta : \Q$，设
  %
  \begin{equation*}
    x_\epsilon \jdeq \max (f(y_0), \ldots, f(y_n))
    \quad\text{and}\quad
    x_\eta \jdeq \max (f(z_0), \ldots, f(z_m))
  \end{equation*}
  %
  其中 $y_0, \ldots, y_n$ 是某个 $h(\epsilon)$-网，$z_0, \ldots, z_m$ 是某个 $h(\eta)$-网。
  每个 $z_i$ 仅仅与某个 $y_j$ 在 $h(\epsilon)$ 意义下足够接近，因此 $|f(z_i) - f(y_j)| <
  \epsilon$，于是可推出
  %
  \begin{equation*}
    f(z_i) < \epsilon + f(y_j) \leq \epsilon + x_\epsilon,
  \end{equation*}
  %
  从而 $x_\eta < \epsilon + x_\epsilon$。对称地可得 $x_\epsilon < \eta +
  x_\eta$，因此 $|x_\eta - x_\epsilon| < \eta + \epsilon$。

  断言 $m \defeq \lim x$ 是 $f$ 的上确界。要证对所有 $x : M$ 有 $f(x) \leq m$，只需证
  $\lnot (m < f(x))$。
  反设 $m <
  f(x)$。则存在 $\epsilon : \Qp$ 使得 $m + \epsilon < f(x)$。
  而在定义 $x_\epsilon$ 时参与的某个 $y_i$ 仅仅满足 $|f(x) - f(y_i)| <
  \epsilon$，于是 $m < f(x) - \epsilon < f(y_i) \leq m$，矛盾。

  最后证明 $m$ 满足定理第二部分；这将自动给出其为最小上界。
  给定任意 $\epsilon : \Qp$，一方面有
  $|m - f(x_{\epsilon/2})| < 3 \epsilon/4$，另一方面在定义 $x_{\epsilon/2}$ 时参与的某个 $y_i$ 仅仅满足
  $|f(x_{\epsilon/2}) - f(y_i)| <
  \epsilon/4$，因此取 $u \defeq y_i$，由三角不等式得 $|m - f(u)| < \epsilon$。
\end{proof}

现在，若在 \cref{ctb-uniformly-continuous-sup} 中还知道 $M$ 完备，我们就希望把假设从一致连续弱化为连续，并把结论加强为上确界在某点处取得。
这些改进的通常证明依赖于以下事实（在完备且全有界空间中）：
%
\begin{enumerate}
\item 连续蕴含一致连续；
\item 每个序列都有收敛子序列。
\end{enumerate}
%
第一条可由 Heine--Borel 紧致性直接得到，第二条正是 Bolzano--Weierstra\ss{} 紧致性。
\index{compactness!Bolzano--Weierstrass@Bolzano--Weierstra\ss{}}%
不幸的是，这两者都有些问题。先说明 Bolzano--Weierstra\ss{} 紧致性蕴含排中律的一个实例，即\define{有限全知原理}：
\indexsee{axiom!limited principle of omniscience}{limited principle of omniscience}%
\indexdef{limited principle of omniscience}%
对每个 $\alpha : \N \to \bool$，
%
\begin{equation} \label{eq:lpo}
  \Parens{\sm{n : \N} \alpha(n) = \btrue} +
  \Parens{\prd{n : \N} \alpha(n) = \bfalse}.
\end{equation}
%
从计算角度看，我们并不期望该原理成立，因为它要求我们判定一个函数是否在无穷多个取值上都为~$\bfalse$。

\begin{thm} \label{analysis-bw-lpo}
  %
  $[0,1]$ 的 Bolzano--Weierstra\ss{} 紧致性蕴含有限全知原理。
  \index{compactness!Bolzano--Weierstrass@Bolzano--Weierstra\ss{}}%
\end{thm}

\begin{proof}
  给定任意 $\alpha : \N \to \bool$，定义序列\index{sequence} $x : \N \to [0,1]$：
  %
  \begin{equation*}
    x_n \defeq
    \begin{cases}
      0 & \text{if $\alpha(k) = \bfalse$ for all $k < n$,}\\
      1 & \text{if $\alpha(k) = \btrue$ for some $k < n$}.
    \end{cases}
  \end{equation*}
  %
  若 Bolzano--Weierstra\ss{} 性质成立，则存在严格递增的 $f : \N \to
  \N$，使得 $x \circ f$ 是 Cauchy 序列\index{Cauchy!sequence}。
  当 $n :
  \N$ 足够大时，第 $n$ 项 $x_{f(n)}$ 与其极限相距小于 $1/6$。
  此时要么 $x_{f(n)} < 2/3$，要么 $x_{f(n)} > 1/3$。
  若 $x_{f(n)} < 2/3$，则~$x_n$ 收敛到 $0$，从而 $\prd{n : \N}
  \alpha(n) = \bfalse$。
  若 $x_{f(n)} > 1/3$，则 $x_{f(n)} = 1$，因此 $\sm{n : \N}
  \alpha(n) = \btrue$。
\end{proof}

对于 Bolzano--Weierstra\ss{} 紧致性，也许我们不会太惋惜它的缺失；但 Heine--Borel 紧致性似乎更难舍弃，经典数学与 Brouwer 直觉主义都接受了它，这点可作证明。
因为我们不想过深进入一般拓扑，这里只讨论基本开集。
在 $\R$ 的情形，基本开集就是有理端点开区间。
这类区间按某类型~$I$ 索引的一族，将是一个映射
%
\begin{equation*}
  \mathcal{F} : I \to \setof{(q, r) : \Q \times \Q | q < r},
\end{equation*}
%
其含义是：满足 $q < r$ 的有理数对 $(q, r)$ 对应类型 $\setof{ x : \R | q < x < r}$。
稍微更方便的是也允许退化区间，因此我们把\define{基本区间族}
\indexdef{family!of basic intervals}%
\indexdef{interval!family of basic}%
定义为映射
%
\begin{equation*}
  \mathcal{F} : I \to \Q \times \Q.
\end{equation*}
%
更精确地说，一个族是依值对 $(I, \mathcal{F})$，而不只是
$\mathcal{F}$。
\define{有限基本区间族}是指索引类型为 $\setof{ m :
  \N | m < n}$（某个 $n : \N$）的区间族；我们通常把它记作有限列表 $[(q_0, r_0), \ldots,
(q_{n-1}, r_{n-1})]$。
最后，$(I, \mathcal{F})$ 的一个\define{有限子族}\indexdef{subfamily, finite, of intervals}由一个索引列表 $[i_1, \ldots, i_n]$ 给出，相应确定有限族
$[\mathcal{F}(i_1), \ldots, \mathcal{F}(i_n)]$。

只要我们注意区分有理数对 $(q, r)$ 与相应区间 $\setof{ x : \R | q < x < r}$，就可以安全地对两者都使用记号 $(q, r)$。
区间的交\indexdef{intersection!of intervals}与包含\indexdef{inclusion!of intervals}\indexdef{containment!of intervals}可由端点表达：
%
\symlabel{interval-intersection}
\symlabel{interval-subset}
\begin{align*}
  (q, r) \cap (s, t) &\ \defeq\  (\max(q, s), \min(r, t)),\\
  (q, r) \subseteq (s, t) &\ \defeq\ (q < r \Rightarrow s \leq q < r \leq t).
\end{align*}
%
我们称 $\intfam{i}{I}{(q_i, r_i)}$\define{（逐点地）覆盖 $[a,b]$}
\indexdef{interval!pointwise cover}%
\indexdef{cover!pointwise}%
\indexdef{pointwise!cover}%
，当且仅当
%
\begin{equation} \label{eq:cover-pointwise-truncated}
  \fall{x : [a,b]} \exis{i : I} q_i < x < r_i.
\end{equation}
%
\define{$[0,1]$ 的 Heine--Borel 紧致性}
\indexdef{compactness!Heine--Borel}%
断言：$[0,1]$ 的每个覆盖族仅仅有一个仍覆盖 $[0,1]$ 的有限子族。

\index{depression}
\begin{thm} \label{classical-Heine-Borel}
  \index{excluded middle}%
  若排中律成立，则 $[0,1]$ 是 Heine--Borel 紧的。
\end{thm}

\begin{proof}
  为导出矛盾，设区间族 $\intfam{i}{I}{(a_i,
    b_i)}$ 覆盖 $[0,1]$，但没有任何有限子族覆盖它。
  我们构造闭区间序列 $[q_n, r_n]$，使其嵌套、长度趋于~$0$，并且每个都不能被 $\intfam{i}{I}{(a_i, b_i)}$ 的有限子族覆盖。

  令 $[q_0, r_0] \defeq [0,1]$。已构造 $[q_n, r_n]$ 后，设
  $s \defeq (2 q_n + r_n)/3$，$t \defeq (q_n + 2 r_n)/3$。
  $[q_n, t]$ 与 $[s, r_n]$ 都被 $\intfam{i}{I}{(a_i, b_i)}$ 覆盖，但它们不可能同时有有限子覆盖，否则 $[q_n, r_n]$ 也会有。
  由排中律，$[q_n, t]$ 要么有有限子覆盖，要么没有。
  若有，则令 $[q_{n+1}, r_{n+1}] \defeq [s, r_n]$；否则令 $[q_{n+1},
  r_{n+1}] \defeq [q_n, t]$。

  序列 $q_0, q_1, \ldots$ 与 $r_0, r_1, \ldots$ 都是 Cauchy 的，并收敛到某个点 $x : [0,1]$，且 $x$ 位于每个 $[q_n, r_n]$ 中。
  仅仅存在 $i : I$ 使得 $a_i < x < b_i$。
  由于区间 $[q_n, r_n]$ 的长度趋于零，存在 $n : \N$ 使得 $a_i < q_n \leq x
  \leq r_n < b_i$。
  这意味着 $[q_n, r_n]$ 被单个区间 $(a_i,
  b_i)$ 覆盖；但同时它又没有有限子覆盖，矛盾。
\end{proof}

若没有排中律，或不加入一点 Brouwer 式直觉主义，我们似乎就卡住了。
不过，$[0,1]$ 的 Heine--Borel 紧致性\emph{确实可以}在构造性语境中恢复，而且仍与经典数学相容！
为此我们需要重新审视“覆盖”的概念。
\eqref{eq:cover-pointwise-truncated} 的问题在于：截断存在允许一种杂乱无章的覆盖方式，从计算上看，我们几乎不可能仅仅提取出一个有限子覆盖。
去掉截断可得
%
\begin{equation} \label{eq:cover-pointwise}
  \prd{x : [0,1]} \sm{i : I} q_i < x < r_i,
\end{equation}
%
这也许有帮助，但它对覆盖的要求又过强。
按这个定义，我们甚至无法证明 $(0,3)$ 与 $(2,5)$ 覆盖 $[1,4]$，因为那等价于构造一个非常值映射 $[1,4] \to \bool$，见
\cref{ex:reals-non-constant-into-Z}。
这里可以借鉴“无点拓扑”\index{pointfree topology}%
\index{topology!pointfree}%
（即 locale 理论）\index{locale}%
的思想：覆盖概念应当用开集来表达，而不诉诸点。
这种对空间的“整体式”观点将允许我们分析覆盖概念，并恢复 Heine--Borel 紧致性。
locale 理论使用幂集，\index{power set}%
若假设命题缩放即可获得；\index{propositional!resizing}%
但我们也可以借用 locale 理论的直谓对应物中的思想，\index{mathematics!predicative}%
即“形式拓扑”。\index{formal!topology}%

\index{acceptance|(}

设我们有一个族 $\pairr{I, \mathcal{F}}$ 与一个区间 $(a, b)$。
如何在不引用点的前提下表达“$(a,b)$ 被该族覆盖”？
一种方式是：若 $(a, b)$ 等于某个 $\mathcal{F}(i)$，则它被该族覆盖。
另一种方式是：若 $(a,b)$ 被另一个族 $(J, \mathcal{G})$ 覆盖，且每个 $\mathcal{G}(j)$ 又被 $\pairr{I, \mathcal{F}}$ 覆盖，则 $(a,b)$ 被
$\pairr{I, \mathcal{F}}$ 覆盖。
注意我们是在列出可用来\emph{推导}“$\pairr{I, \mathcal{F}}$ 覆盖 $(a,b)$”的\emph{规则}。
我们应找到足够好的规则，并把它们变成归纳定义。

\begin{defn} \label{defn:inductive-cover}
  %
  \define{归纳覆盖关系 $\cover$}
  \indexdef{inductive!cover}%
  \indexdef{cover!inductive}%
  是一个纯关系
  %
  \begin{equation*}
    {\cover} : (\Q \times \Q) \to \Parens{\sm{I : \type} (I \to \Q \times \Q)} \to \prop
  \end{equation*}
  %
  它由以下规则归纳定义，其中 $q, r, s, t$ 为有理数，
  $\pairr{I, \mathcal{F}}$ 与 $\pairr{J, \mathcal{G}}$ 是基本区间族：
  %
  \begin{enumerate}

  \item \emph{自反性：}
    \index{reflexivity!of inductive cover}%
    对任意 $i : I$，有 $\mathcal{F}(i) \cover \pairr{I, \mathcal{F}}$，

  \item \emph{传递性：}
    \index{transitivity!of inductive cover}%
    若 $(q, r) \cover \pairr{J, \mathcal{G}}$ 且 $\fall{j : J} \mathcal{G}(j) \cover \pairr{I,\mathcal{F}}$
    则 $(q, r) \cover \pairr{I, \mathcal{F}}$，

  \item \emph{单调性：}
    \index{monotonicity!of inductive cover}%
    若 $(q, r) \subseteq (s, t)$ 且 $(s,t) \cover \pairr{I, \mathcal{F}}$，则 $(q, r) \cover
    \pairr{I, \mathcal{F}}$，

  \item \emph{局部化：}
    \index{localization of inductive cover}%
    若 $(q, r) \cover (I, \mathcal{F})$，则 $(q, r) \cap (s, t) \cover
    \intfam{i}{I}{(\mathcal{F}(i) \cap (s, t))}$。

  \item \label{defn:inductive-cover-interval-1}
    若 $q < s < t < r$，则 $(q, r) \cover [(q, t), (s, r)]$，

  \item \label{defn:inductive-cover-interval-2}
    $(q, r) \cover \intfam{u}{\setof{ (s,t) : \Q \times \Q | q < s < t < r}}{u}$。
  \end{enumerate}
\end{defn}

这个定义应被读作一个高阶归纳类型：上列规则是点构造子，而该类型是 $(-1)$-截断的。
前四条是一般性质，直观上较清楚。
后两条是实直线特有的：一条说两个相交区间可以覆盖一个区间；另一条说一个区间可以由其内部覆盖。
顺便一提，若 $r
\leq q$，则由最后一条可知 $(q, r)$ 被空族覆盖。

归纳覆盖满足 Heine--Borel 性质，其证明需要一个引理。

\begin{lem} \label{reals-formal-topology-locally-compact}
  若 $q < s < t < r$ 且 $(q, r) \cover \pairr{I, \mathcal{F}}$，则仅仅存在 $\pairr{I, \mathcal{F}}$ 的一个有限子族，归纳地覆盖 $(s, t)$。
\end{lem}

\begin{proof}
  对 $(q, r) \cover \pairr{I, \mathcal{F}}$ 做归纳证明。共有六种情形：
  %
  \begin{enumerate}

  \item 自反性：若 $(q, r) = \mathcal{F}(i)$，则由单调性可知 $(s, t)$ 被有限子族 $[\mathcal{F}(i)]$ 覆盖。

  \item 传递性：
    设 $(q, r) \cover \pairr{J, \mathcal{G}}$ 且 $\fall{j : J} \mathcal{G}(j) \cover
    \pairr{I, \mathcal{F}}$。由归纳假设，仅仅存在
    $[\mathcal{G}(j_1), \ldots, \mathcal{G}(j_n)]$ 覆盖 $(s, t)$。
    再由归纳假设，每个 $\mathcal{G}(j_k)$ 都被 $\pairr{I, \mathcal{F}}$ 的某个有限子族覆盖；把这些收集起来，即得覆盖 $(s, t)$ 的有限子族。

  \item 单调性：
    若 $(q, r) \subseteq (u, v)$ 且 $(u, v) \cover \pairr{I, \mathcal{F}}$，则可把归纳假设应用于 $(u, v) \cover \pairr{I, \mathcal{F}}$，因为 $u <
    s < t < v$。

  \item 局部化：
    设 $(q', r') \cover \pairr{I, \mathcal{F}}$ 且 $(q, r) = (q', r') \cap (a, b)$。
    因为 $q' < s < t < r'$，由归纳假设有 $(s, t)$ 的有限子覆盖
    $[\mathcal{F}(i_1), \ldots, \mathcal{F}(i_n)]$。
    又有 $a < s <
    t < b$，故 $(s, t) = (s, t) \cap (a, b)$ 被
    $[\mathcal{F}(i_1) \cap (a,b), \ldots, \mathcal{F}(i_n) \cap (a,b)]$ 覆盖，
    这是 $\intfam{i}{I}{(\mathcal{F}(i) \cap (a, b))}$ 的一个有限子族。

  \item 若对某些 $q < u < v < r$ 有 $(q, r) \cover [(q, v), (u, r)]$，则由单调性
    $(s, t) \cover [(q, v), (u, r)]$。

  \item 最后，$(s, t) \cover \intfam{z}{\setof{ (u,v):\Q \times \Q | q < u < v < r}}{z}$ 由自反性直接得到。\qedhere
  \end{enumerate}
\end{proof}

称 $\pairr{I, \mathcal{F}}$\define{归纳覆盖
  $[a, b]$}，若仅仅存在 $\epsilon : \Qp$ 使得 $(a - \epsilon, b +
\epsilon) \cover \pairr{I, \mathcal{F}}$。

\begin{cor} \label{interval-Heine-Borel}
  \index{compactness!Heine-Borel}%
  \index{interval!open and closed}%
  闭区间对归纳覆盖满足 Heine--Borel 紧致性。
\end{cor}

\begin{proof}
  设 $[a, b]$ 被 $\pairr{I, \mathcal{F}}$ 归纳覆盖，则仅仅存在
  $\epsilon : \Qp$ 使得 $(a - \epsilon, b + \epsilon) \cover \pairr{I, \mathcal{F}}$。
  由 \cref{reals-formal-topology-locally-compact}，$(a - \epsilon/2, b + \epsilon/2)$ 有一个有限子覆盖，故它也是 $[a, b]$ 的有限子覆盖。
\end{proof}

形式拓扑\index{topology!formal}的经验表明，归纳覆盖的这些规则足以支撑点自由拓扑的构造性发展。
但我们也可以给出自己的证据，说明它是一个合理概念。

\begin{thm} \label{inductive-cover-classical}
  \mbox{}
  %
  \begin{enumerate}
  \item 归纳覆盖也是逐点覆盖。
  \item 在排中律下，逐点覆盖也是归纳覆盖。
  \end{enumerate}
\end{thm}

\begin{proof}
  \mbox{}
  %
  \begin{enumerate}

  \item
    取一个基本区间族 $\pairr{I, \mathcal{F}}$，其中记 $(q_i,
    r_i) \defeq \mathcal{F}(i)$；再取一个被 $\pairr{I,
      \mathcal{F}}$ 归纳覆盖的区间 $(a,b)$，以及满足 $a < x < b$ 的点 $x$。
    %
    对 $(a,b) \cover \pairr{I, \mathcal{F}}$ 归纳证明：仅仅存在 $i : I$ 使得 $q_i < x < r_i$。
    多数情形都很直接，这里只展示两种。
    若 $(a,b) \cover \pairr{I, \mathcal{F}}$ 由自反性得到，则仅仅存在某个 $i : I$ 使得 $(a,b) = (q_i, r_i)$，故 $q_i < x < r_i$。
    若 $(a,b)
    \cover \pairr{I, \mathcal{F}}$ 由经 $\intfam{j}{J}{(s_j, t_j)}$ 的传递性得到，则由归纳假设仅仅存在 $j : J$ 使得 $s_j < x < t_j$；而又因 $(s_j, t_j) \cover \pairr{I, \mathcal{F}}$，再由归纳假设仅仅存在
    $i : I$ 使得 $q_i < x < r_i$。
    其余情形同样“精彩”。

  \item 设 $\intfam{i}{I}{(q_i, r_i)}$ 逐点覆盖 $(a, b)$。由
    \cref{defn:inductive-cover} 的 \cref{defn:inductive-cover-interval-2}
    可知，只需证明：当 $a < c < d < b$ 时，$\intfam{i}{I}{(q_i, r_i)}$ 归纳覆盖 $(c, d)$。
    于是取这样的 $c,d$。
    由 \cref{classical-Heine-Borel}，存在有限子族 $[i_1, \ldots, i_n]$ 已经逐点覆盖 $[c,
    d]$，从而也覆盖 $(c,d)$。
    令 $\epsilon : \Qp$ 为
    \index{Lebesgue number}
    $(q_{i_1}, r_{i_1}), \ldots, (q_{i_n}, r_{i_n})$ 的一个 Lebesgue 数（见
    \cref{ex:finite-cover-lebesgue-number}）。
    存在正整数 $k : \N$ 使得 $2 (d - c)/k
    < \min(1, \epsilon)$。对 $0 \leq i \leq k$，令
    %
    \begin{equation*}
      c_k \defeq ((k - i) c + i d) / k.
    \end{equation*}
    %
    区间 $(c_0, c_2)$、$(c_1, c_3)$、\dots、$(c_{k-2}, c_k)$ 可通过反复使用传递性与
    \cref{defn:inductive-cover} 中的~\cref{defn:inductive-cover-interval-1}
    归纳覆盖 $(c,d)$。
    由于这些区间宽度都小于 $\epsilon$，每个都包含在某个 $(q_i, r_i)$ 中；再由传递性和单调性可得
    $\intfam{i}{I}{(q_i, r_i)}$ 归纳覆盖 $(c, d)$。\qedhere
  \end{enumerate}
\end{proof}

上一条定理的要点是：就经典数学而言，逐点覆盖与归纳覆盖没有差别。
特别地，由于在同伦类型论中假设排中律是一致的，我们无法给出“归纳覆盖但非逐点覆盖”的例子。
换句话说，二者差异不在于覆盖了什么，而在于“覆盖成立”的\emph{证明}。

若继续展开，这里又可以再写一本书；但我们就此打住，并希望上述内容已足以支持“分析学可以在同伦类型论中发展”这一主张。
感兴趣的读者可参见 \cref{ex:mean-value-theorem}，其中给出了介值定理的构造性版本。

\index{acceptance|)}

\index{mathematics!classical|)}%
\index{mathematics!constructive|)}%
\section{超现实数}
\label{sec:surreals}

\index{surreal numbers|(}%

本节我们讨论另一个高阶归纳-归纳类型的例子，它把前文许多线索汇聚在一起：Conway 的\emph{超现实数}域 \NO~\cite{conway:onag}。
超现实数是（Dedekind）实数（\cref{sec:dedekind-reals}）与序数（\cref{sec:ordinals}）的自然共同推广。
Conway 在带排中律与选择公理的经典\index{mathematics!classical}数学中，把一个超现实数定义为一对超现实数的\emph{集合}，记作 $\surr L R$，并要求 $L$ 的每个元素都严格小于 $R$ 的每个元素。
这显然像一个归纳定义，但把它当作归纳定义有三个问题。

首先，该定义需要超现实数之间的（严格）不等关系，因此该关系必须与超现实数类型 \NO 同时定义。
（Conway 通过先定义\emph{博弈}\index{game!Conway}来回避这个问题；博弈类似超现实数，但去掉了 $L$ 与 $R$ 的相容条件。）
与 Cauchy 实数的关系 $\closesym$ 一样，这种同时定义\emph{先验地}既可能是归纳-归纳，也可能是归纳-递归。
出于与 $\closesym$ 相同的理由，我们选择归纳-归纳方式。

此外，我们会分别（并互归纳地）定义超现实数的严格不等 $<$ 与非严格不等 $\le$。
Conway 把 $<$ 用 $\le$ 来定义，这在经典上合理，但在构造性语境中不合理。
\index{mathematics!constructive}%
并且，若用负定义来定义 $<$，它将不能作为高阶归纳类型构造子的假设（见 \cref{sec:strictly-positive}）。

其次，Conway 说 $\surr L R$ 中的 $L,R$ 应是“超现实数的集合”，但若把它朴素理解为谓词 $\NO\to\prop$，它并非正性的，因此不能作为归纳构造子的输入。
然而，这样的理解本来也不是 Conway 原意的好类型论翻译，因为在集合论里超现实数形成一个真类，而 $L,R$ 是真正的（小）集合，而非 \NO 的任意子类。
在类型论中，这意味着 \NO 将相对于某个宇宙 \UU 定义，但它自身属于下一层更高宇宙 $\UU'$，类似序数与基数的集合 \ord、\card，集合论累积层级 $V$，或在没有命题缩放时的 Dedekind 实数。
\index{propositional!resizing}%
于是我们会要求超现实数“集合” $L,R$ 是 \UU-小的，因此很自然地把它们表示为由某个 \UU-小类型索引的超现实数\emph{族}。
（这与我们在 \cref{sec:cumulative-hierarchy} 对累积层级所做的处理完全一致。）
也就是说，超现实数构造子的类型将是
\[ \prd{\LL,\RR:\UU} (\LL\to\NO) \to (\RR\to \NO) \to (\text{some condition}) \to \NO \]
这确实是严格正性的。\index{strict!positivity}

最后，在给出 \NO 及其序关系的互定义后，Conway 宣称：若 $x\le y$ 且 $y\le x$，则两个超现实数 $x,y$ \emph{相等}。
这自然可读作：把“预超现实数”集合按某个等价关系取商。
%(In set-theoretic foundations, one has to us an additional trick to deal with large equivalence classes.)
然而，在没有选择公理时，这种商与通常 Cauchy 实数构造中的商面临同样问题：我们不再能保证一对\emph{超现实数族}总能给出新的超现实数 $\surr L R$，因为不一定能把 $L,R$“提升”为预超现实数族。
当然，我们可以像处理 Cauchy 实数时一样，用\emph{高阶}归纳-归纳定义来解决这个问题。

\begin{defn}\label{defn:surreals}
  类型 \NO（\define{超现实数}）
  \indexdef{surreal numbers}%
  \indexsee{number!surreal}{surreal numbers}%
  及关系 $\mathord<:\NO\to\NO\to\type$ 与 $\mathord\le:\NO\to\NO\to\type$，按如下方式高阶归纳-归纳定义。
  类型 \NO 具有下列构造子。
  \begin{itemize}
  \item 对任意 $\LL,\RR:\UU$ 与函数 $\LL\to \NO$、$\RR\to \NO$，其值分别记为 $x^L$ 与 $x^R$（其中 $L:\LL$、$R:\RR$），若 $\fall{L:\LL}{R:\RR} x^L<x^R$，则有一个超现实数 $x$。
  \item 对任意 $x,y:\NO$，若 $x\le y$ 且 $y\le x$，则有 $\noeq(x,y):x=y$。
  \end{itemize}
  我们把第一构造子的输入称为一个\define{切割}。
  \indexdef{cut!of surreal numbers}%
  若 $x$ 由某个切割构造，则记号 $x^L$ 默认假设 $L:\LL$，同理 $x^R$ 默认假设 $R:\RR$。
  这样通常可避免显式命名索引类型 $\LL,\RR$，在同时讨论多个切割时尤为方便。
  依 Conway 的术语，我们称 $x^L$ 为 $x$ 的\emph{左选项}\indexdef{option of a surreal number}，$x^R$ 为\emph{右选项}。

  路径构造子意味着不同切割可以定义同一个超现实数。
  因而，除非还知道某个超现实数 $x$ 由特定切割给出，否则谈论其左/右选项是没有意义的。
  所以下文我们会说“给定一个定义超现实数 $x$ 的切割”，而不是仅说“给定一个超现实数 $x$”。

  关系 $\le$ 有以下构造子。
  \index{non-strict order}%
  \index{order!non-strict}%
  \begin{itemize}
  \item 给定定义两个超现实数 $x,y$ 的切割，若对所有 $L$ 有 $x^L<y$，且对所有 $R$ 有 $x<y^R$，则 $x\le y$。
  \item 命题截断：
    对任意 $x,y:\NO$，若 $p,q:x\le y$，则 $p=q$。
  \end{itemize}
  而关系 $<$ 有以下构造子。
  \index{strict!order}%
  \index{order!strict}%
  \begin{itemize}
    % Don't technically need x in the first one and y in the second one to be defined by cuts?
  \item 给定定义两个超现实数 $x,y$ 的切割，若存在某个 $L$ 使得 $x\le y^L$，则 $x<y$。
  \item 给定定义两个超现实数 $x,y$ 的切割，若存在某个 $R$ 使得 $x^R\le y$，则 $x<y$。
  \item 命题截断：对任意 $x,y:\NO$，若 $p,q:x<y$，则 $p=q$。
  \end{itemize}
\end{defn}

\noindent
把这与 Conway 的定义对照：
\begin{itemize}\footnotesize
\item[-] 若 $L,R$ 是任意两个数集，且 $L$ 中没有元素 $\ge$ $R$ 中任何元素，则存在一个数 $\surr L R$。
  所有数都以这种方式构造。
\item[-] $x\ge y$ 当且仅当（不存在 $x^R\le y$ 且 $x\le$ 不存在 $y^L$）。
\item[-] $x=y$ 当且仅当（$x \ge y$ 且 $y\ge x$）。
\item[-] $x>y$ 当且仅当（$x\ge y$ 且 $y\not\ge x$）。
\end{itemize}
若对象全是[超现实]数而非“博弈”\index{game!Conway}，则在 $x>y$ 的定义中包含 $x\ge y$ 是多余的。
因此 Conway 的 $<$ 只是其 $\ge$ 的否定，于是他对 $\surr L R$ 成为超现实数的条件与我们的条件相同。
把 Conway 的 $\le$ 取否并消去双重否定，就得到我们的 $<$ 定义；随后可把他的 $\le$ 改写为不含否定、仅用 $<$ 的形式。

我们可以立即在 \NO 中填入许多超现实数。
依 Conway 的记号，写作
\symlabel{surreal-cut}
\[\surr{x,y,z,\dots}{u,v,w,\dots}\]
表示由某个切割定义的超现实数，其中 $\LL\to\NO$ 与 $\RR\to\NO$ 分别由 $x,y,z,\dots$ 和 $u,v,w,\dots$ 描述。
当然，若 $\LL$ 或 $\RR$ 为 $\emptyt$，就把对应部分留空。
这与子集标准记号 $\setof{x:A | P(x)}$ 有不幸冲突，但本节不会使用后者。
\begin{itemize}
\item 我们递归定义 $\iota_{\nat}:\nat\to\NO$：
  \begin{align*}
    \iota_{\nat}(0) &\defeq \surr{}{},\\
    \iota_\nat(\suc(n)) &\defeq \surr{\iota_\nat(n)}{}.
  \end{align*}
  即 $\iota_\nat(0)$ 由切割 $\emptyt\to\NO$ 与 $\emptyt\to\NO$ 定义。
  同理，$\iota_\nat(\suc(n))$ 由 $\unit\to\NO$（选出 $\iota_\nat(n)$）与 $\emptyt\to\NO$ 定义。
\item 同样地，我们用符号分情形递归原理（\cref{thm:sign-induction}）定义 $\iota_{\Z}:\Z\to\NO$：
  \begin{align*}
    \iota_{\Z}(0) &\defeq \surr{}{},\\
    \iota_\Z(n+1) &\defeq \surr{\iota_\Z(n)}{} & &\text{$n\ge 0$,}\\
    \iota_\Z(n-1) &\defeq \surr{}{\iota_\Z(n)} & &\text{$n\le 0$.}
  \end{align*}
\item \define{二进有理数}
  \indexdef{rational numbers!dyadic}%
  \indexsee{dyadic rational}{rational numbers, dyadic}%
  指一对 $(a,n)$，其中 $a:\Z$、$n:\nat$，并要求当 $n>0$ 时 $a$ 为奇数。
  我们写作 $a/2^n$，并将其认同为对应有理数。
  若 $\Q_D$ 表示二进有理数集合，则按 $n$ 归纳定义 $\iota_{\Q_D}:\Q_D\to\NO$：
  \begin{align*}
    \iota_{\Q_D}(a/2^0) &\defeq \iota_\Z(a),\\
    \iota_{\Q_D}(a/2^n) &\defeq \surr{\iota_{\Q_D}(a/2^n - 1/2^n)}{\iota_{\Q_D}(a/2^n + 1/2^n)},
    \quad \text{for $n>0$.}
  \end{align*}
  这里用到：若 $n>0$ 且 $a$ 为奇数，则 $a/2^n \pm 1/2^n$ 仍是二进有理数，且分母比 $a/2^n$ 的更小。
\item 我们定义 $\iota_{\RD}:\RD\to\NO$，\label{reals-into-surreals}其中 $\RD$ 是（任一版本的）\cref{sec:dedekind-reals} 中的 Dedekind 实数：
  \begin{align*}
    \iota_{\RD}(x) &\defeq
    \surr{q\in\Q_D \text{ such that } q<x}{q\in\Q_D \text{ such that } x<q}.
  \end{align*}
  与前面几种情形不同，把二进有理数视作 Dedekind 实数时，此定义是否扩张了 $\iota_{\Q_D}$ 并不显然。
  这将由简单性定理（\cref{thm:NO-simplicity}）推出。
\item 回忆 \cref{sec:ordinals} 的\emph{序数}\index{ordinal}类型 \ord，它在关系 $<$ 下是良序的，其中 $A<B$ 表示对某个 $b:B$ 有 $A = \ordsl B b$。
  我们沿 \ord 的良基递归（\cref{thm:wfrec}）定义 $\iota_{\ord}:\ord\to\NO$\label{ord-into-surreals}：
  \begin{equation*}
    \iota_{\ord}(A) \defeq
    \surr{\iota_\ord(\ordsl A a) \text{ for all } a:A}{}.
  \end{equation*}
  由简单性定理还可推出：$\iota_\ord$ 在有限序数上的限制与 $\iota_\nat$ 一致。
  （但提醒读者：不同于上述例子，除非限制到更小的一类序数，否则 $\iota_\ord$ 在构造性上并非单射；见 \cref{ex:ord-into-surreals,ex:hiit-plump}。）
\item Conway 中另外几个有趣例子：
  \begin{align*}
    \omega &\defeq \surr{0,1,2,3,\dots}{} \qquad\text{(also an ordinal)}\\
    -\omega &\defeq \surr{}{\dots,-3,-2,-1,0}\\
    1/\omega &\defeq \textstyle\surr{0}{1,\frac12,\frac14,\frac18,\dots}\\
    \omega-1 &\defeq \surr{0,1,2,3,\dots}{\omega}\\
    \omega/2 &\defeq \surr{0,1,2,3,\dots}{\dots,\omega-2,\omega-1,\omega}.
  \end{align*}
\end{itemize}

在识别由不同切割给出的超现实数时，下述简单观察很有用。

\begin{thm}[Conway 的简单性定理]\label{thm:NO-simplicity}
  \index{simplicity theorem}%
  \index{theorem!Conway's simplicity}%
  设 $x,z$ 是由切割定义的超现实数，且满足：
  \begin{itemize}
  \item 对所有 $L,R$，有 $x^L < z < x^R$。
  \item 对 $z$ 的每个左选项 $z^L$，存在 $x$ 的某个左选项 $x^{L'}$ 使得 $z^L\le x^{L'}$。
  \item 对 $z$ 的每个右选项 $z^R$，存在 $x$ 的某个右选项 $x^{R'}$ 使得 $x^{R'}\le z^R$。
  \end{itemize}
  则 $x=z$。
\end{thm}
\begin{proof}
  应用 \NO 的路径构造子，只需证明 $x\le z$ 与 $z\le x$。
  对前者，需要对所有 $L$ 证明 $x^L<z$（这是已知的），并对所有 $R$ 证明 $x<z^R$。
  但按假设，对任意 $z^R$ 存在 $x^{R'}$ 使得 $x^{R'}\le z^R$，从而 $x<z^R$。
  故 $x\le z$；$z\le x$ 的证明对称。
\end{proof}

\index{induction principle!for surreal numbers}
然而，要进一步讨论超现实数，我们需要它们的归纳原理。
$(\NO,\le,<)$ 的互归纳原理适用于三族类型：
\begin{align*}
  A &: \NO\to\type\\
  B &: \prd{x,y:\NO}{a:A(x)}{b:A(y)} (x\le y) \to \type\\
  C &: \prd{x,y:\NO}{a:A(x)}{b:A(y)} (x<y) \to \type.
\end{align*}
与 Cauchy 实数的归纳原理类似，把 $B,C$ 看作类型 $A(x)$ 与 $A(y)$ 之间的关系族会更有帮助。
\symlabel{NO-recursion}
因此，把 $B(x,y,a,b,\xi)$ 记为 $(x,a) \ble^\xi (y,b)$，把 $C(x,y,a,b,\xi)$ 记为 $(x,a) \blt^\xi (y,b)$。
同样地，我们通常省略 $\xi$，因为它只居留于一个纯命题而不含实质信息；$x,y$ 也常可省略，只写 $a\ble b$ 或 $a\blt b$。
用这些记号，归纳原理的假设如下。
\begin{itemize}
\item 对任意定义超现实数 $x$ 的切割，若给定
  \begin{enumerate}
  \item 每个 $L$ 的元素 $a^L:A(x^L)$，以及
  \item 每个 $R$ 的元素 $a^R:A(x^R)$，并且
  \item 对所有 $L,R$ 有 $(x^L,a^L) \blt (x^R,a^R)$，
  \end{enumerate}
  则有指定元素 $f_a:A(x)$。
  我们把这类数据称为定义 $x$ 的切割上的\define{依值切割}
  \indexdef{cut!of surreal numbers!dependent}%
  \indexdef{dependent!cut}%
  。
\item 对任意 $x,y:\NO$ 及 $a:A(x), b:A(y)$，若 $x\le y$ 且 $y\le x$，并且 $(x,a) \ble (y,b)$
  与 $(y,b) \ble (x,a)$ 都成立，
  则有 $\dpath{A}{\noeq}{a}{b}$。
\item 给定定义超现实数 $x,y$ 的切割，以及 $x$ 上的依值切割 $a$ 与 $y$ 上的依值切割 $b$，若对所有 $L$ 有 $x^L<y$ 且 $(x^L,a^L)\blt (y,f_b)$，
  且对所有 $R$ 有 $x<y^R$ 且 $(x,f_a) \blt (y^R,b^R)$，
  则 $(x,f_a) \ble (y,f_b)$。
\item $\ble$ 取值于纯命题。
\item 给定定义超现实数 $x,y$ 的切割、依值切割 $a,b$，以及某个 $L_0$ 使得 $x\le y^{L_0}$ 且 $(x,f_a) \ble (y^{L_0},b^{L_0})$，
  则有 $(x,f_a) \blt (y,f_b)$。
\item 给定定义超现实数 $x,y$ 的切割、依值切割 $a,b$，以及某个 ${R_0}$ 使得 $x^{R_0}\le y$ 且 $(x^{R_0},a^{R_0}),\ble (y,f_b)$，
  则有 $(x,f_a) \blt (y,f_b)$。
\item $\blt$ 取值于纯命题。
\end{itemize}
在这些假设下可推出函数 $f:\prd{x:\NO} A(x)$，并满足
\begin{align}
  f(x) &\;\jdeq\; f_{f[x]} \label{eq:noind1}\\
  (x\le y) &\;\Rightarrow\; (x,f(x)) \ble (y,f(y)) \notag\\
  (x< y) &\;\Rightarrow\; (x,f(x)) \blt (y,f(y)). \notag
\end{align}
在点构造子的计算规则~\eqref{eq:noind1} 中，$x$ 是由某个切割定义的超现实数，$f[x]$ 表示把 $f$ 作用到 $x$ 所定义的依值切割（并利用 $f$ 把 $<$ 送到 $\blt$ 这一事实）。
照例，我们一般使用模式匹配记号：在切割 $\surr{x^L}{x^R}$ 上定义 $f$ 时，可使用符号 $f(x^L),f(x^R)$，并假设它们形成依值切割。

与 Cauchy 实数一样，通过把 $A,\ble,\blt$ 的部分成分平凡化可得到若干特例。
取 $\ble,\blt$ 都恒为 \unit，即得 \define{\NO-归纳}；为简洁起见，我们只陈述纯命题版本：
\begin{itemize}
\item 给定 $P:\NO\to\prop$，若每当 $x$ 由一个切割定义且对所有 $L,R$ 都有 $P(x^L),P(x^R)$ 时可推出 $P(x)$，则对所有 $x:\NO$ 都有 $P(x)$。
\end{itemize}
这可与 Conway 的一段话对照：
\begin{quote}\footnotesize
  一般地，当我们希望对所有数 $x$ 建立命题 $P(x)$ 时，我们将通过归纳来证明：由所有命题 $P(x^L)$ 与 $P(x^R)$ 的成立推出 $P(x)$。
  我们认为“所有数都以这种方式构造”这句话足以说明这种做法的正当性。
\end{quote}
借助 \NO-归纳，我们可证明

\begin{thm}[Conway 的定理 0]\label{thm:NO-refl-opt}\
  \index{theorem!Conway's 0}%
  \begin{enumerate}
  \item 对任意 $x:\NO$，有 $x\le x$。\label{item:NO-le-refl}
  \item 对任意由切割定义的 $x:\NO$，对所有 $L,R$ 有 $x^L <x$ 且 $x<x^R$。\label{item:NO-lt-opt}
  \end{enumerate}
\end{thm}
\begin{proof}
  先注意：若 $x\le x$，则每当 $x$ 作为某个切割 $y$ 的左选项出现时，由 $<$ 的第一构造子可得 $x<y$；同理，若 $x$ 作为右选项出现，则由第二构造子得 $y<x$。
  特别地，~\ref{item:NO-le-refl}$\Rightarrow$\ref{item:NO-lt-opt}。

  下面对 $x$ 用 \NO-归纳证明~\ref{item:NO-le-refl}。
  设 $x$ 由一个切割定义，且对所有 $L,R$ 都有 $x^L\le x^L$ 与 $x^R \le x^R$。
  由上面的观察，这些假设推出对所有 $L,R$ 都有 $x^L<x$ 与 $x<x^R$，进而由 $\le$ 的构造子得到 $x\le x$。
\end{proof}

\begin{cor}\label{thm:NO-set}
  \NO 是一个 0-类型。
%  (As with $V$, it might be confusing to say that it is a ``set''.)
\end{cor}
\begin{proof}
  纯关系 $R(x,y)\defeq (x\le y) \land (y\le x)$ 由 \NO 的路径构造子蕴含恒等，并由 \cref{thm:NO-refl-opt}\ref{item:NO-le-refl} 包含对角线。
  因此可应用 \cref{thm:h-set-refrel-in-paths-sets}。
\end{proof}

相比之下，Conway 的 Theorem 1（$\le$ 的传递性）在我们的定义下要更难建立；见 \cref{thm:NO-unstrict-transitive}。

% Of course, we also have:

% \begin{lem}
%   Every surreal number is merely defined by a cut.
% \end{lem}
% \begin{proof}
%   Obvious by $\NO$-induction.
% \end{proof}

我们还需要联合递归原理，\define{$(\NO,\le,<)$-recursion}。
下面用如下方式陈述更方便。
设 $A$ 是一个类型，配备关系 $\mathord\ble:A\to A\to\prop$ 与 $\mathord\blt:A\to A\to\prop$。
则可按以下步骤定义 $f:\NO\to A$。
\begin{enumerate}
\item 对任意由切割定义的 $x$，在已定义 $f(x^L)$ 与 $f(x^R)$ 且对所有 $L,R$ 满足 $f(x^L)\blt f(x^R)$ 的前提下，我们必须定义 $f(x)$。（这称为递归的\emph{主子句}。）\label{item:NO-rec-primary}
\item 证明 $\ble$ 是\emph{反对称的}\index{relation!antisymmetric}：若 $a\ble b$ 且 $b\ble a$，则 $a=b$。
\item 对由切割定义的 $x,y$，若对所有 $L$ 有 $x^L<y$、对所有 $R$ 有 $x<y^R$，并归纳地假设对所有 $L$ 有 $f(x^L)\blt f(y)$、对所有 $R$ 有 $f(x)\blt f(y^R)$，且对所有 $L,R$ 也有 $f(x^L)\blt f(x^R)$ 与 $f(y^L)\blt f(y^R)$，则必须证明 $f(x)\ble f(y)$。
\item 对由切割定义的 $x,y$ 与某个 $L_0$，若 $x\le y^{L_0}$，并归纳地假设 $f(x)\ble f(y^{L_0})$，且对所有 $L,R$ 有 $f(x^L)\blt f(x^R)$ 与 $f(y^L)\blt f(y^R)$，则必须证明 $f(x)\blt f(y)$。
\item 对由切割定义的 $x,y$ 与某个 $R_0$，若 $x^{R_0}\le y$，并归纳地假设 $f(x^{R_0})\ble f(y)$，且对所有 $L,R$ 有 $f(x^L)\blt f(x^R)$ 与 $f(y^L)\blt f(y^R)$，则必须证明 $f(x)\blt f(y)$。\label{item:NO-rec-last}
\end{enumerate}
最后三条也可更简洁地说成：必须证明按第一条定义的 $f$ 把 $\le$ 映到 $\ble$，把 $<$ 映到 $\blt$。
我们将这一性质称为\emph{$f$ 保持不等关系}。
此外，在证明 $f$ 保持不等关系时，可假设当前的 $\le$ 或 $<$ 来自其某个构造子，并可使用归纳假设：$f$ 保持该构造子输入中出现的所有不等关系。

若我们完成了上面的~\ref{item:NO-rec-primary}--\ref{item:NO-rec-last}，就得到 $f:\NO\to A$：它在切割上的计算如~\ref{item:NO-rec-primary} 所指定，且保持所有不等关系：
%
\begin{narrowmultline*}
  \fall{x,y:\NO}\Big((x\le y) \to (f(x)\ble f(y))\Big) \land
  \narrowbreak
  \Big((x< y) \to (f(x)\blt f(y))\Big).
\end{narrowmultline*}
%
和 Cauchy 实数的 $(\RC,\closesym)$-递归一样，这个递归原理对定义 \NO 上函数至关重要，因为我们无法先在“预超现实数”上定义函数，再事后证明它尊重相等关系。

\begin{eg}
  下面定义\emph{取负}函数 $\NO\to\NO$。
  我们使用联合递归原理，取 $A\defeq\NO$，并令 $(x\ble y)\defeq (y\le x)$、$(x\blt y)\defeq (y< x)$。
  显然此时 $\ble$ 是反对称的。

  对定义的主子句，设 $x$ 由某个切割定义，且 $-x^L$、$-x^R$ 已定义并满足对所有 $L,R$ 有 $-x^L \blt -x^R$。
  按定义这意味着对所有 $L,R$ 有 $-x^R< -x^L$，所以可由切割 $\surr{-x^R}{-x^L}$ 定义 $-x$。
  这一记号遵循 Conway：其左选项由索引 $x$ 右选项的类型 $\RR$ 给出，右选项由索引 $x$ 左选项的类型 $\LL$ 给出；相应族 $\RR\to\NO$ 与 $\LL\to\NO$ 由 $x$ 的对应族与取负复合得到。

  现在需验证 $f$ 保持不等关系。
  \begin{itemize}
  \item 对 $x\le y$，可设对所有 $L$ 有 $x^L<y$、对所有 $R$ 有 $x < y^R$，需证 $-y\le -x$。
    但归纳地可设 $-y <-x^L$ 且 $-y^R<-x$，由 $-y$、$-x$ 的定义及 $\le$ 的构造子即得结论。
  \item 对 $x<y$，第一种情形来自某个 $x\le y^{L_0}$，可归纳假设 $-y^{L_0} \le -x$，于是由 $<$ 的构造子得 $-y<-x$。
  \item 同理，若 $x<y$ 来自 $x^{R_0}\le y$，则归纳假设是 $-y \le -x^R$，同样推出 $-y<-x$。
  \end{itemize}
\end{eg}

不过若要做更多事情，我们需要像在 \cref{thm:RC-sim-characterization} 对 Cauchy 实数那样，更显式地刻画关系 $\le$ 与 $<$。
同样地，我们还必须同步证明这些关系的若干关键性质，归纳才能推进。

\begin{thm}\label{defn:No-codes}
  存在关系 $\mathord\preceq:\NO\to\NO\to\prop$ 与 $\mathord\prec:\NO\to\NO\to\prop$，使得当 $x,y$ 是由切割定义的超现实数时，
  \begin{align*}
    (x\preceq y) &\defeq
    \big(\fall{L} x^L\prec y\big) \land \big(\fall{R} x\prec y^R\big)\\
    (x\prec y) &\defeq
    \big(\exis{L} x\preceq y^L\big) \lor \big(\exis{R} x^R \preceq y\big).
  \end{align*}
  另外有
  \begin{equation}\label{eq:NO-codes-unstrict}
    (x\prec y) \to (x\preceq y)
  \end{equation}
  以及所有合理的传递性质，使得 $\prec$ 与 $\preceq$ 构成关于 $\le$ 与 $<$ 的一个“二模”\index{bimodule}：
  \begin{equation}\label{eq:NO-codes-transitivity}
    \begin{array}{c@{\hspace{1cm}}c}
      (x \le y) \to (y\preceq z) \to (x\preceq z) &
      (x \preceq y) \to (y\le z) \to (x\preceq z) \\
      (x \le y) \to (y\prec z) \to (x\prec z) &
      (x \preceq y) \to (y< z) \to (x\prec z) \\
      (x < y) \to (y\preceq z) \to (x\prec z) &
      (x \prec y) \to (y\le z) \to (x\prec z).
  \end{array}
  \end{equation}
\end{thm}

\begin{proof}
  我们对 $x,y$ 做双重 $(\NO,\le,<)$-归纳来定义 $\preceq$ 与 $\prec$。
  第一层归纳是简单递归，其陪域是 $(\NO\to\prop)\times (\NO\to\prop)$ 的一个子集 $A$：其中元素是一对谓词，一个蕴含另一个，并满足“右传递性”，即~\eqref{eq:NO-codes-unstrict} 以及~\eqref{eq:NO-codes-transitivity} 右列（把 $(x\preceq \blank)$ 与 $(x\prec \blank)$ 替换为给定的两个谓词）。
  与 \cref{defn:RC-approx} 的证明一样，我们把这些谓词看作二元关系的一半，记为 $y\mapsto (\hle y)$ 与 $y\mapsto (\hlt y)$，并用 $\hlname$ 表示这对关系。
  % The precise definition of $A$ is
  % \begin{align*}
  %   A\defeq \bigg\{ \hlname : (\NO\to\prop)\times (\NO\to\prop) \;\bigg|\;\\
  %   \begin{split}
  %     \fall{y,z:\NO}
  %     &\Big( (\hle y) \to (y\le z) \to (\hle z) \Big)\\
  %     \land\; &\Big( (\hle y) \to (y< z) \to (\hlt z) \Big)\\
  %     \land\; &\Big( (\hlt y) \to (y\le z) \to (\hlt z) \Big)\\
  %     \land\; &\Big( (\hlt y) \to (y< z) \to (\hlt z) \Big) \bigg\}
  %   \end{split}
  % \end{align*}
  我们给 $A$ 配上如下两个关系：
  \begin{align*}
    (\hlname \ble \hlbname) &\defeq
    \fall{y:\NO} \Big( (\hleb y) \to (\hle y) \Big) \land
    \Big( (\hltb y) \to (\hlt y) \Big),\\
    (\hlname \blt \hlbname) &\defeq
    \fall{y:\NO} \Big( (\hleb y) \to (\hlt y) \Big).
    %\land \Big( (\hltb y) \to (\hlt y) \Big)
  \end{align*}
  注意 $\ble$ 是反对称的：若 $\hlname \ble \hlbname$ 且 $\hlbname \ble \hlname$，则对所有 $y$ 有 $(\hleb y) \Leftrightarrow (\hle y)$ 与 $(\hltb y) \Leftrightarrow (\hlt y)$，故由纯命题上的泛等与函数外延性得 $\hlname=\hlbname$。
  另外，说函数 $\NO\to A$ 保持不等关系，恰好等价于说把它看作 \NO 上一对二元关系时满足“左传递性”（即~\eqref{eq:NO-codes-transitivity} 左列）。

  现在看外层递归的主子句：设给定由某切割定义的 $x$，并且对所有 $L,R$ 已有关系 $(x^L \prec \blank)$、$(x^R \prec \blank)$、$(x^L \preceq \blank)$、$(x^R \preceq \blank)$，其中严格者蕴含非严格者，满足右传递性，并且
  \begin{equation}\label{eq:NO-prec-outer-IH}
    \fall{L,R}{y:\NO}\Big( (x^R\preceq y) \to (x^L \prec y) \Big).
    % \land\Big( (x^R \prec y) \to (x^L \prec y) \Big)
  \end{equation}
  我们现在要对所有 $y$ 定义 $(x\prec y)$ 与 $(x\preceq y)$。
  与 \cref{defn:RC-approx} 不同，这里不用内层递归，而用内层归纳，以便可归纳地使用关于不等式 $x^L<x$ 与 $x<x^R$ 的左传递性。
  定义 $A':\NO\to\type$，其中 $A'(y)$ 取为 $\prop\times\prop$ 的子集 $A'$，由两个纯命题组成，分别记作 $\tle y$ 与 $\tlt y$（记对应元素为 $\tlname:A'(y)$），并满足
  \begin{gather}
    (\tlt y) \to (\tle y)\\
    \fall{L} (\tle y)\to (x^L\prec y) \label{eq:NO-prec-IHL}\\
    \fall{R} (x^R \preceq y) \to (\tlt y) \label{eq:NO-prec-IHR}.
  \end{gather}
  类似于 $\ble$ 与 $\blt$ 的记号，对 $\tlname:A'(y)$ 与 $\tlbname:A'(z)$ 定义关系
  \begin{align*}
    (\tlname \bble \tlbname) &\defeq
    \Big((\tle y) \to (\tleb z)\Big) \land \Big((\tlt y) \to (\tltb z)\Big)\\
    (\tlname \bblt \tlbname) &\defeq
    \Big((\tle y) \to (\tltb z)\Big). % \land \Big(\tlt \to \tltb\Big).
  \end{align*}
  % (These are the type families $B$ and $C$ in the general induction principle.)
  同样地，$\bble$ 在适当意义下显然反对称。
  此外，一个保持不等关系的函数 $\prd{y:\NO} A'(y)$，恰好是一对满足“严格蕴含非严格”、右传递性，以及相对于不等式 $x^L<x$ 与 $x<x^R$ 的左传递性的谓词。
  因而，这个内层归纳将给出完成外层递归主子句所需的一切。

  对内层归纳的主子句，再设给定由切割定义的 $y$，以及对所有 $L,R$ 的性质 $(x\prec y^L)$、$(x\prec y^R)$、$(x\preceq y^L)$、$(x\preceq y^R)$，其中严格者蕴含非严格者，且在 $x^L<x$ 与 $x<x^R$ 意义下左传递，并在 $y^L<y^R$ 意义下右传递。
  % \begin{equation}
  %   \fall{L,R}\Big((x \preceq y^L) \to (x \prec y^R)\Big) % \land \Big((x \prec y^L) \to (x\prec y^R)\Big).
  %   \label{eq:NO-prec-inner-IH}
  % \end{equation}
  现在可给出定理中指定的定义：
  \begin{align}
    (x\preceq y) &\defeq
    (\fall{L} x^L\prec y) \land (\fall{R} x\prec y^R), \label{eq:NO-preceq-def}\\
    (x\prec y) &\defeq
    (\exis{L} x\preceq y^L) \lor (\exis{R} x^R \preceq y).\label{eq:NO-prec-def}
  \end{align}
  要使这给出 $A'(y)$ 的一个元素，首先要证明 $(x\prec y) \to (x\preceq y)$。
  假设 $x\prec y$，有两种情形。
  一方面，若存在 $L_0$ 使 $x\preceq y^{L_0}$，则由关于 $y^{L_0}<y^R$ 的右传递性，对所有 $R$ 有 $x\prec y^R$。
  又由关于 $x^L<x$ 的左传递性，对任意 $L$ 有 $x^L \prec y^{L_0}$，再由右传递性得 $x^L\prec y$。
  因而 $x\preceq y$。

  另一方面，若存在 $R_0$ 使 $x^{R_0}\preceq y$，则由~\eqref{eq:NO-prec-outer-IH} 可得对所有 $L$ 有 $x^L \prec y$。
  同时由关于 $x<x^{R_0}$ 与 $y<y^R$ 的左右传递性，对任意 $R$ 有 $x\prec y^R$。
  因而 $x\preceq y$。

  还要证明这些定义在 $x^L<x$ 与 $x<x^R$ 意义下左传递。
  但若 $x\preceq y$，按定义对所有 $L$ 立即有 $x^L\prec y$；而若 $x^R\preceq y$，按定义也有 $x\prec y$。

  因此，~\eqref{eq:NO-preceq-def} 与~\eqref{eq:NO-prec-def} 的确定义了 $A'(y)$ 的一个元素。
  接着要验证该定义作为到 $A'$ 的依值函数保持不等关系，即这些关系在右侧是传递的。
  记住：每个情形中都可归纳地假设它们对于不等式构造子输入中出现的所有不等关系都具右传递性。
  \begin{itemize}
  \item 设 $x\preceq y$ 且 $y\le z$，后者来自对所有 $L,R$ 的 $y^L<z$ 与 $y<z^R$。
    则把（内层递归的）归纳假设应用于 $y<z^R$，得对任意 $R$ 有 $x\prec z^R$。
    又由定义，$x\preceq y$ 蕴含对任意 $L$ 有 $x^L \prec y$，故由外层递归的归纳假设得 $x^L \prec z$。
    因而 $x\preceq z$。
  \item 设 $x\preceq y$ 且 $y<z$。
    第一种情形：$y<z$ 来自 $y\le z^{L_0}$。
    则把内层归纳假设应用于 $y\le z^{L_0}$，得 $x \preceq z^{L_0}$，从而 $x\prec z$。

    第二种情形：$y<z$ 来自 $y^{R_0}\le z$。
    由定义，$x\preceq y$ 蕴含 $x\prec y^{R_0}$，再把内层归纳假设用于 $y^{R_0}\le z$，得 $x\prec z$。
  \item 设 $x\prec y$ 且 $y\le z$，后者来自对所有 $L,R$ 的 $y^L<z$ 与 $y<z^R$。
    按定义，$x\prec y$ 蕴含仅仅存在某个 $R_0$ 使 $x^{R_0}\preceq y$，或某个 $L_0$ 使 $x\preceq y^{L_0}$。
    若 $x^{R_0}\preceq y$，则由外层归纳假设得 $x^{R_0}\preceq z$，故 $x\prec z$。
    若 $x\preceq y^{L_0}$，则对 $y^{L_0}<z$（它由 $y\le z$ 的构造子给出）应用内层归纳假设，得 $x\prec z$。
  % \item Suppose $x\prec y$ and $y<z$.
  %   First, suppose $y<z$ arises from $y\le z^{L_0}$.
  %   Then the inner inductive hypothesis for $y\le z^{L_0}$ yields $x\prec z^{L_0}$, hence $x\preceq z^{L_0}$; thus $x\prec z$.

  %   Second, suppose $y<z$ arises from $y^{R_0}\le z$.
  %   Then by definition, $x\prec y$ implies there merely exists $R_1$ with $x^{R_1}\preceq y$ or $L_1$ with $x\preceq y^{L_1}$.
  %   If $x^{R_1}\preceq y$, then the outer inductive hypothesis implies $x^{R_1}\prec z$, hence $x^{R_1}\preceq z$, and thus $x\prec z$.
  %   And if $x\preceq y^{L_1}$, then the inner inductive hypothesis applied to $y^{L_1}<y^{R_0}$ (which comes from $y$ being defined as a cut) and $y^{R_0}\le z$ yields $x\prec z$.
  \end{itemize}
  这完成了内层归纳。
  因而，对任意由切割定义的 $x$，我们已由~\eqref{eq:NO-preceq-def} 与~\eqref{eq:NO-prec-def} 定义了 $(x\prec \blank)$ 与 $(x\preceq \blank)$，并且它们在右侧是传递的。

  为完成外层递归，还需验证这些定义在左侧也传递。
  对 $z$ 再做一次 \NO-归纳后，会得到三个与上面“右传递”几乎完全相同的情形。
  因而略去。
\end{proof}

\begin{thm}\label{thm:NO-encode-decode}
  对任意 $x,y:\NO$，有 $(x<y)=(x\prec y)$ 且 $(x\le y)=(x\preceq y)$。
\end{thm}
\begin{proof}
  从左到右，用 $(\NO,\le,<)$-归纳，其中 $A(x)\defeq\unit$，并由 $\preceq,\prec$ 提供关系 $\ble,\blt$。
  在所有构造子情形里，$x,y$ 都由切割定义，所以 $\preceq,\prec$ 的定义可化简，归纳假设可直接应用。

  从右到左，用 \NO-归纳假设 $x,y$ 由切割定义。
  此时 $\preceq,\prec$ 的定义及归纳假设恰好提供了 $\le$ 与 $<$ 相应构造子所需的数据。
  % From right to left, we first prove by $\NO$-induction on $x$ that for any $y:\NO$ we have $(x\prec y) \to (x<y)$ and $(x\preceq y) \to (x\le y)$.
  % Thus, we assume this to be true for all $x^L$ and $x^R$ in a cut, and show it for the resulting $x:\NO$.
  % Next, we prove by $\NO$-induction on $y$ that $(x\prec y) \to (x<y)$ and $(x\preceq y) \to (x\le y)$, hence we assume it to be true for all $y^L$ and $y^R$ in a cut, and show it for the resulting $y:\NO$.
  % Now since $x$ and $y$ are both defined by cuts, $x\preceq y$ means that $x^L\prec y$ and $x\prec y^R$ for all $L$ and $R$.
  % By the inductive hypotheses, this gives $x^L<y$ and $x<y^R$, hence $x\le y$ by the constructor of $\le$.
  % Similarly, $x\prec y$ yields merely an $R_0$ with $x^{R_0}\preceq y$ or an $L_0$ with $x\preceq y^{L_0}$.
  % Hence merely $x^{R_0}\le y$ or $x\le y^{L_0}$ by the inductive hypothesis, so $x<y$ by a constructor.
\end{proof}

\begin{cor}\label{thm:NO-unstrict-transitive}
  \NO 上的关系 $\le$ 与 $<$ 满足
  \[ \fall{x,y:\NO} (x<y) \to (x\le y) \]
  且具有传递性：
  \index{transitivity!of . for surreals@of $<$ for surreals}
  \index{transitivity!of . for surreals@of $\leq$ for surreals}
  \begin{gather*}
    (x\le y) \to (y\le z) \to (x\le z)\\
    (x\le y) \to (y< z) \to (x< z)\\
    (x< y) \to (y\le z) \to (x< z).
  \end{gather*}
\end{cor}

与 Cauchy 实数一样，在定义 \NO 上各种运算时，联合 $(\NO,\le,<)$-递归原理仍是必不可少的。

\begin{eg}\label{eg:surreal-addition}
\index{addition!of surreal numbers}%
我们把 $\mathord+:\NO\to\NO\to\NO$ 定义为：先对第一参数递归，再对第二参数归纳。
外层递归的陪域取为 $\NO\to\NO$ 的一个子集，包含那些函数 $g$，它们对所有 $x,y$ 满足 $(x<y) \to (g(x)<g(y))$ 与 $(x\le y) \to (g(x)\le g(y))$。
对这样的 $g,h$，定义 $(g\ble h)\defeq \fall{x:\NO} g(x)\le h(x)$ 与 $(g\blt h)\defeq \fall{x:\NO} g(x)< h(x)$。
显然 $\ble$ 是反对称的。

对递归主子句，设 $x$ 由某个切割定义，且函数 $(x^L+\blank)$ 与 $(x^R+\blank)$ 已定义、保持不等关系，并满足 $x^L+y<x^R+y$，我们需定义 $(x+\blank)$。
与 \cref{defn:No-codes} 一样，这里不用内层递归，而用内层归纳到族 $A:\NO\to\type$，其中 $A(y)$ 是满足“每个 $x^L + y < z$ 且每个 $x^R + y > z$”的那些 $z:\NO$ 构成的子集。
给 $A$ 赋予由 \NO 诱导来的关系 $\le$ 与 $<$，于是反对称性显然。
内层递归的主子句中，再设 $y$ 由某个切割定义，$x+y^L$ 与 $x+y^R$ 都已定义，并满足 $x^L+y^L < x+y^L$、$x^L+y^R < x+y^R$、$x+y^L < x^R + y^L$、$x+y^R < x^R+y^R$（这些来自对 $A(y)$ 元素施加的附加条件），另外还有 $x+y^L < x+y^R$（因为 $x+y^L$ 与 $x+y^R$ 在 $A(y)$ 中形成依值切割）。
现在给出 Conway 的定义：
\[ x+y \defeq \surr{x^L+y, x+y^L}{x^R+y,x+y^R}. \]
换言之，$x+y$ 的左选项是所有形如 $x^L+y$（某个左选项 $x^L$）或 $x+y^L$（某个左选项 $y^L$）的数。
我们必须证明这些左选项中的每一个都小于这些右选项中的每一个：
\begin{itemize}
\item 由外层归纳假设，$x^L+y < x^R+y$。
\item $x^L+y < x^L + y^R < x + y^R$，前者因 $(x^L+\blank)$ 保持不等关系，后者因 $x+y^R : A(y^R)$。
\item $x+y^L < x^R+ y^L < x^R + y$，前者因 $x+y^L : A(y^L)$，后者因 $(x^R+\blank)$ 保持不等关系。
\item 由内层归纳假设（具体即“形成依值切割”这一事实），$x+y^L < x+y^R$。
\end{itemize}
还要证明如此定义的 $x+y$ 属于 $A(y)$，即 $x^L + y < x+y$ 且 $x+y < x^R + y$；这由 \cref{thm:NO-refl-opt}\ref{item:NO-lt-opt} 直接给出。

接着要验证 $(x+\blank)$ 的定义保持不等关系：
\begin{itemize}
\item 若 $y\le z$ 来自“对所有 $L,R$ 有 $y^L<z$ 与 $y<z^R$”，则内层归纳假设给出 $x+y^L<x+z$ 与 $x+y < x+z^R$，外层归纳假设给出 $x^L+y \le x^L+z$ 与 $x^R+ y \le x^R+z$。
  另外，由定义 $x^R+y$ 是 $x+y$ 的右选项，所以 $x+y < x^R+y$。
  同理，$x^L+z$ 是 $x+z$ 的左选项，因此 $x^L+z < x+z$。
  于是用传递性可得 $x^L+y < x+z$ 且 $x+y < x^R+z$；故由 $\le$ 的构造子推出 $x+y \le x+z$。
\item 若 $y<z$ 来自某个 $L_0$ 满足 $y\le z^{L_0}$，则归纳地有 $x+y \le x+z^{L_0}$，而 $x+z^{L_0}$ 是 $x+z$ 的右选项，故 $x+y<x+z$。
\item 同理，若 $y<z$ 来自 $y^{R_0}\le z$，则由 $x+y^{R_0}\le x+z$ 可得 $x+y<x+z$。
\end{itemize}
这完成内层归纳。
外层递归还需验证 $+$ 在左侧也保持不等关系；再做一次 \NO-归纳即可，过程完全类似。
\end{eg}

\index{acceptance|(}%
\index{mathematics!formalized}%
在~\cite{conway:onag} Part Zero 的附录中，Conway 讨论了如何在 ZFC 集合论中形式化超现实数：要么沿序数迭代并在每个等价类中选取最低秩代表构成集合，要么用“符号展开”来表示数。
随后他说：
\begin{quote}\footnotesize
  这些构造令人惊讶的复杂性，与其说揭示了我们的数系本身，不如说更多揭示了 ZF 内部形式化的性质\dots
\end{quote}
并继续主张应建立一个关于“可允许构造方式”的一般理论，其中应包括
\begin{enumerate}\footnotesize
\item 对象可以由先前对象以任何合理构造性的方式生成。\label{item:conway1}
\item 新生成对象之间的相等性可以取任意希望的等价关系。\label{item:conway2}
\end{enumerate}
\noindent
条件~\ref{item:conway1} 很自然地可读作对\emph{归纳定义}一般原理的正当化，例如 \cref{sec:strictly-positive,sec:generalizations} 中给出的那些。
特别地，构造子严格正性的条件可看作“合理构造性”这一说法的形式化。
条件~\ref{item:conway2} 则提示我们应把这点扩展到各种\emph{高阶}归纳定义，在其中可通过路径构造子，以任意合理方式把对象判等。
例如 Conway 在下一段说：
\begin{quote}\footnotesize
  \dots 例如，我们也可以自由地创造一个新对象 $(x,y)$，并称之为 $x$ 与 $y$ 的有序对。
  我们还可以再创造一个有序对 $[x,y]$，它与 $(x,y)$ 不同，但与之并存\dots
  若我们反而希望把 $(x,y)$ 做成无序对，就可以把相等定义为如下等价关系：$(x,y)=(z,t)$ 当且仅当 $x=z,y=t$ \emph{或} $x=t,y=z$。
\end{quote}
引入由特定构造子生成、具有新名字的新对象的自由，正是我们在归纳定义理论中所拥有的。
正如 \cref{sec:appetizer-univalence} 中自然数 $\nat$ 与 $\nat'$ 那样，若我们把笛卡尔积类型 $A\times B$ 的定义再原样写一遍，就会得到一个不同的积类型 $A\times' B$，其典范元素可以自由写作 $[x,y]$。
而且我们可以通过加入合适路径构造子，把其中一个做成“无序对”类型。% (and perhaps 0-truncating).

当然，Conway 的本意并不是只批评 ZF，而是同时面向所有基础理论：
\begin{quote}\footnotesize
  \dots 这个提议并不是要拿某个特定理论来替代 ZF\dots{}
  真正的提议是：我们应赋予自己自由，去创建这类任意的数学理论，同时证明一个元定理，彻底保证任何这类理论都可用任一标准基础理论来形式化。
\end{quote}
可以回应说：实际上，单值基础并不是 Conway 所设想“标准基础理论”之一，而更像是可在其中表达“创建新理论能力”并证明 Conway 元定理的\emph{元理论}。
例如，超现实数正是 Conway 心目中的一种“数学理论”，而我们已看到它可在单值基础内部构造并获得正当化。
同样，Conway 更早还说过：
\begin{quote}\footnotesize
  \dots 集合论将是这样一种理论：集合由先前集合通过对应通常公理的过程构造出来，而相等关系就是“具有相同成员”。
\end{quote}
这种描述与 \cref{sec:cumulative-hierarchy} 中通过高阶归纳方式构造集合论累积层级高度吻合。
Conway 的元定理则对应于我们多次提到的事实：可在 ZFC 内部构造单值基础的一个模型（这超出本书范围）。

然而，单值基础自身已足够丰富而强大，把它仅仅降格为构造类集合理论的元理论是很不明智的。
我们已经看到，即便在集合（0-类型）层面，单值基础中的高阶归纳类型也给出了按泛性质直接构造对象的方法（\cref{sec:free-algebras}），例如 Cauchy 完备化的构造性理论（\cref{sec:cauchy-reals}）。
但更重要的是，能在基础系统内直接建模同伦论与范畴论（\cref{cha:homotopy,cha:category-theory}）这一潜力，使单值基础具有任何集合论基础都无法匹敌的优势。
\index{acceptance|)}%

\index{surreal numbers|)}%

\sectionNotes

在 Dedekind 实数上定义代数运算，尤其是乘法，既有些棘手也相当繁琐。
让算术运转起来有多种做法：各有优点，但似乎都需要一定技术工作。
例如 Richman~\cite{Richman:reals} 先在正切割上定义 Dedekind 实数的乘法，再以代数方式扩展到全部 Dedekind 切割；而 Conway~\cite{conway:onag} 指出，超现实数的乘法定义同样适用于 Dedekind 实数。

我们对 Dedekind 实数的处理借鉴了~\cite{BauerTaylor09} 的许多思想；该文在 Abstract Stone Duality 语境中构造了 Dedekind 实数。
\index{Abstract Stone Duality}%
这是一种（受限的）简单类型 $\lambda$-演算，带有一个区分对象 $\Sigma$ 用来分类开集，并由对偶也分类闭集。在~\cite{BauerTaylor09} 中还可以找到算术运算基本性质的详细证明。

\cref{RC-initial-Cauchy-complete} 证明了 $\RC$ 是最小的 Cauchy 完备阿基米德有序域，这表明我们的 Cauchy 实数很可能与 Escard{\'o}--Simpson 实数一致~\cite{EscardoSimpson:01}。
\index{real numbers!Escardo--Simpson@Escard\'o--Simpson}%
检验这点是否确实成立将会很有意思。Escard{\'o}--Simpson 实数（更准确地说是对应闭区间）的概念很有意义，因为它可在任何带有限积的范畴中陈述。

在加入“正则扩张公理”的构造性集合论中，也可以尝试通过超限迭代下对 Cauchy 序列极限封闭来定义 Cauchy 完备化。
检验该构造是否与我们的构造一致，同样很有意思。

在构造性数学中，Cauchy 实数与 Dedekind 实数重合需要依赖选择，这几乎是民间常识；而“可数选择已足够”则较少为人知。
回忆\define{依赖选择}
\indexdef{axiom!of choice!dependent}%
\index{axiom!of choice!countable}%
\index{total!relation}%
是指：对 $A$ 上一个总关系 $R$（即 $\fall{x : A} \exis{y : A} R(x,y)$），以及任意 $a : A$，仅仅存在 $f : \N \to A$ 满足 $f(0) = a$ 且对所有 $n : \N$ 有 $R(f(n), f(n+1))$。
我们的 \cref{when-reals-coincide} 使用了把依赖选择的一次应用转化为可数选择应用的典型技巧：先用一次可数选择预先做完所有可能出现的选择，再用所得选择函数避免依赖选择。

在构造性语境下，不同紧致性概念之间复杂关系见 \cite{bridges2002compactness}。
Palmgren~\cite{Palmgren:FT} 对逐点分析与无点拓扑做了很好的比较。

超现实数最初由~\cite{conway:onag} 定义，使用了某种归纳式定义，但并未以任何基础系统术语显式给出正当化。
因此，一些后续作者倾向于使用符号展开或其他更显式表示，以便更明显地编码进集合论。
在类型论中表示超现实数的想法最早由 Hancock 考虑；而
Setzer 与 Forsberg~\cite{forsbergfinite} 指出，超现实数及其不等关系 $<$ 与 $\le$ 自然形成一个归纳-归纳定义。
这里给出的\emph{高阶}归纳-归纳版本是新的：它把超现实数正确的相等概念直接内建进去。


\sectionExercises

\begin{ex}\label{ex:alt-dedekind-reals}
  给出 Dedekind 实数的另一种定义：先定义平方，再使用 \cref{mult-from-square}。
  检查所得结构是否为交换环。
\end{ex}

\begin{ex} \label{ex:RD-extended-reals}
  %
  设我们去掉 \cref{defn:dedekind-reals} 中
  \ref{defn:dedekind-reals-inhabited} 的有界性条件。
  则得到\define{扩展实数}
  \indexdef{real numbers!extended}%
  \indexdef{extended real numbers}%
  ，其中包含 $-\infty \defeq
  (\emptyt, \Q)$ 与 $\infty \defeq (\Q, \emptyt)$。
  对切割定义的哪些算术运算在扩展实数上仍有意义？可得到什么代数结构？
\end{ex}

\begin{ex} \label{ex:RD-lower-cuts}
  %
  通过考虑单侧切割可得到\define{下}与\define{上} Dedekind 实数，
  \indexdef{real numbers!Dedekind!upper}%
  \indexdef{real numbers!Dedekind!lower}%
  \indexdef{lower Dedekind reals}%
  \indexdef{upper Dedekind reals}%
  \index{cut!Dedekind}%
  分别如下。
  例如，一个下实数由谓词 $L : \Q \to \Omega$ 给出，它满足：
  %
  \begin{enumerate}
  \item \emph{可居留：} $\exis{q : \Q} L(q)$，以及
  \item \emph{圆整：} $L(q) = \exis{r : \Q} q < r \land L(r)$。
    \index{rounded!Dedekind cut}
  \end{enumerate}
  %
  （也可再要求 $\exis{r : \Q} \lnot L(r)$，以排除切割 $\infty \defeq
  \Q$。）
  你能在下实数上定义哪些算术运算？特别地，加法逆元会怎样？
\end{ex}

\begin{ex} \label{ex:RD-interval-arithmetic}
  %
  \index{interval!arithmetic}%
  设我们去掉 \cref{defn:dedekind-reals} 中的可定位性条件。
  则得到\define{区间域}
  \indexdef{interval!domain}%
  $\mathbb{I}$，因为此时切割允许出现“空隙”，即区间。
  在 $\mathbb{I}$ 上定义偏序 $\sqsubseteq$：
  %
  \begin{narrowmultline*}
    ((L, U) \sqsubseteq (L', U'))
    \defeq \narrowbreak
    (\fall{q : \Q} L(q) \Rightarrow L'(q)) \land
    (\fall{q : \Q} U(q) \Rightarrow U'(q)).
  \end{narrowmultline*}
  %
  相对于 $\mathbb{I}$，$\mathbb{I}$ 的极大元是什么？
  定义“端点”运算，把区间域元素映到其下端点与上端点。
  这些端点是实数、下实数还是上实数（见
  \cref{ex:RD-lower-cuts}）？
  对切割定义的哪些算术运算在区间域上仍有意义？
\end{ex}

\begin{ex} \label{ex:RD-lt-vs-le}
  证明对所有 $x, y : \RD$，
  %
  \begin{equation*}
    \lnot (x < y) \Rightarrow y \leq x
  \end{equation*}
  %
  以及
  %
  \begin{equation*}
    \eqv{(x \leq y)}{\Parens{\prd{\epsilon : \Qp} x < y + \epsilon}}.
  \end{equation*}
  %
  另外，$\lnot (x \leq y)$ 是否蕴含 $y < x$？
\end{ex}

\begin{ex} \label{ex:reals-non-constant-into-Z}
  \mbox{}
  %
  \begin{enumerate}
  \item
    在假设排中律下，构造一个非常值映射 $\RD \to \Z$。
  \item
    设 $f : \RD \to \Z$ 满足 $f(0) = 0$ 且对所有 $x >
    0$ 有 $f(x) \neq 0$。由此推出有限全知原理~\eqref{eq:lpo}。
\index{limited principle of omniscience}%
  \end{enumerate}
\end{ex}

\begin{ex} \label{ex:traditional-archimedean}
  \index{ordered field!archimedean}%
  证明在有序域 $F$ 中，$\Q$ 的稠密性与传统阿基米德公理等价：
  %
  \begin{equation*}
    (\fall{x, y : F} x < y \Rightarrow \exis{q : \Q} x < q < y)
    \Leftrightarrow
    (\fall{x : F} \exis{k : \Z} x < k).
  \end{equation*}
\end{ex}

\begin{ex} \label{RC-Lipschitz-on-interval} 设 $a, b : \Q$，且 $f : \setof{ q : \Q |
    a \leq q \leq b } \to \RC$ 是 Lipschitz 常数为~$L$ 的映射。证明存在唯一扩张
  $\bar{f} : [a,b] \to \RC$，它仍以~$L$ 为 Lipschitz 常数。
  提示：不必针对闭区间重做 \cref{RC-extend-Q-Lipschitz}，注意有回缩 $r : \RC \to [-n,n]$，并把
  \cref{RC-extend-Q-Lipschitz} 应用于 $f \circ r$。
\end{ex}

\begin{ex} \label{ex:metric-completion}
  \index{completion!of a metric space}%
  把 $\RC$ 的构造推广到任意度量空间的 Cauchy 完备化。先思考：度量空间的距离\index{distance}函数的陪域采用哪种实数概念最自然？这是否重要？然后完成以下两种构造细节：
  %
  \begin{enumerate}
  \item 仿照 Cauchy 实数构造，把度量空间完备化定义为一个在 Cauchy 序列极限下封闭的归纳-归纳类型。\index{Cauchy!sequence}
  \item 使用 Lawvere~\cite{lawvere:metric-spaces}\index{Lawvere} 与 Richman~\cite{Richman00thefundamental} 的如下构造：度量空间 $(M, d)$ 的完备化定义为\define{位置}的类型。
    \indexdef{location}%
    所谓位置是函数 $f : M \to \R$，满足
    %
    \begin{enumerate}
    \item 对所有 $x, y : M$，有 $f(x) \geq |f(y) - d(x,y)|$，以及
    \item $\inf_{x \in M} f(x) = 0$，其含义是 $\fall{\epsilon : \Qp} \exis{x : M} |f(x)| < \epsilon$ 且 $\fall{x : M} f(x) \geq 0$。
    \end{enumerate}
    %
    直观上，$f$ 看起来像是在测量到某个点的距离。
  \end{enumerate}
  %
  \index{universal!property!of metric completion}%
  最后证明度量完备化的如下泛性质：从某度量空间到某 Cauchy 完备度量空间的局部一致连续映射，可唯一扩张为其完备化上的局部一致连续映射。（映射称为\define{局部一致连续}
  \indexdef{function!locally uniformly continuous}%
  \indexdef{locally uniformly continuous map}%
  若它在开球上是一致连续的。）
\end{ex}

\index{metric space|)}%

\begin{ex} \label{ex:reals-apart-neq-MP}
  \define{Markov 原理}
  \indexdef{axiom!Markov's principle}%
  \indexdef{Markov's principle}%
  断言对所有 $f : \nat \to \bool$，
  %
  \begin{equation*}
    (\lnot \lnot \exis{n : \nat} f(n) = \btrue)
    \Rightarrow
    \exis{n : \nat} f(n) = \btrue.
  \end{equation*}
  %
  这是双重否定律~\eqref{eq:ldn} 的一个特例。
  证明
  $\fall{x, y: \RD} x \neq y \Rightarrow x \apart y$ 可推出 Markov 原理。
  反向是否也成立？
\end{ex}

\begin{ex} \label{ex:reals-apart-zero-divisors}
  \index{apartness}%
  验证实数满足如下“无零因子”性质：
  $x y \apart 0 \Leftrightarrow x \apart 0 \land y \apart 0$。
\end{ex}

\begin{ex} \label{ex:finite-cover-lebesgue-number}
  %
  设 $(q_1, r_1), \ldots, (q_n, r_n)$ 逐点覆盖 $(a, b)$。证明存在
  $\epsilon : \Qp$，使得每当 $a < x < y < b$ 且 $|x - y| < \epsilon$ 时，
  仅仅存在某个 $i$ 使得 $q_i < x < r_i$ 且 $q_i < y < r_i$。
  这样的 $\epsilon$ 称为该覆盖的\define{Lebesgue 数}
  \indexdef{Lebesgue number}%
  。
\end{ex}

\begin{ex} \label{ex:mean-value-theorem}
  %
  证明介值定理的如下近似版本：
  %
  \begin{quote}
    \emph{
      若 $f : [0,1] \to \R$ 一致连续，且 $f(0) < 0 < f(1)$，则
      对每个 $\epsilon : \Qp$，仅仅存在 $x : [0,1]$ 使得 $|f(x)| <
      \epsilon$。
    }
  \end{quote}
  %
  提示：不要尝试二分法，因为它会导向选择公理。
  取而代之，用分段线性映射逼近 $f$。应如何构造分段线性映射？
\end{ex}

\begin{ex}\label{ex:knuth-surreal-check}
  检查~\cite{knuth74:_surreal_number} 中的全部内容，是否都可用 \cref{sec:surreals} 的高阶归纳-归纳超现实数来完成。
\end{ex}

\begin{ex}\label{ex:reals-into-surreals}
  回忆第~\pageref{reals-into-surreals} 页定义的函数 $\iota_{\RD}:\RD\to\NO$。
  \begin{enumerate}
  \item 证明 $\iota_{\RD}$ 是单射。
  \item $\iota_{\RD}$ 有向扩展实数（\cref{ex:RD-extended-reals}）与区间域（\cref{ex:RD-interval-arithmetic}）的自然扩张。
    它们是否是单射？
  \end{enumerate}
\end{ex}

\begin{ex}\label{ex:ord-into-surreals}
  证明第~\pageref{ord-into-surreals} 页定义的函数 $\iota_{\ord}:\ord\to\NO$ 是单射，当且仅当 \LEM{} 成立。
\end{ex}

\begin{ex}\label{ex:hiit-plump}
  仿照 \NO 的定义，仅使用左选项，定义一个配备二元关系 $\le$ 与 $<$ 的类型 $\mathsf{POrd}$。
  \begin{enumerate}
  \item 构造映射 $j:\mathsf{POrd} \to \NO$，并证明它是嵌入。
  \item 证明在关系 $<$ 下，$\mathsf{POrd}$ 是一个序数（与 \ord 类似，位于下一层更高宇宙）。
  \item 假设命题缩放，证明 $\mathsf{POrd}$ 等价于 \cref{ex:plump-ordinals} 中 \ord 的子集
    \[\setof{A:\ord | \mathsf{isPlump}(A)}\]
    并推出：$\iota_{\ord}:\ord\to\NO$ 在 plump 序数上的限制是单射。
  \end{enumerate}
  在没有命题缩放时，我们仍可把 $\mathsf{POrd}$ 的元素（或其在 \NO 中的像）称为\define{plump 序数}。\index{ordinal!plump}\index{plump!ordinal}
\end{ex}

\begin{ex}\label{ex:pseudo-ordinals}
  定义：若某超现实数等于一个无右选项的切割 $\surr{x^L}{}$（但其左选项本身可有右选项），则称它为\define{伪序数}。\index{pseudo-ordinal}\index{ordinal!pseudo-}
  证明命题“每个伪序数都是 plump 序数”等价于 \LEM{}。
\end{ex}

\begin{ex}\label{ex:double-No-recursion}
  注意 \cref{defn:No-codes} 与 \cref{eg:surreal-addition} 都使用了类似模式来定义函数 $\NO \to \NO \to B$：
  先做一个外层 \NO-递归，其陪域是保序函数 $\NO\to B$ 的集合；再做一个内层 \NO-归纳到族 $A:\NO\to\type$，其中 $A(y)$ 是 $B$ 的一个子集，用以确保不等式 $x^L<x$ 与 $x<x^R$ 也被保持。
  试表述并证明一个一般性的“双重 \NO-递归”原理，以推广这些证明。
\end{ex}

\index{real numbers|)}%

%%% Local Variables:
%%% mode: latex
%%% TeX-master: "hott-online"
%%% End:
